\documentclass[11pt,letterpaper]{article}
\usepackage[T1]{fontenc}
\usepackage{lmodern}
\usepackage[margin=1in]{geometry}
\usepackage{amsmath,amssymb,amsthm,mathtools}
\usepackage{booktabs,array,longtable}
\usepackage{microtype}
\usepackage{enumitem}
\usepackage{xcolor}
\usepackage{hyperref}
\hypersetup{colorlinks=true,linkcolor=black,citecolor=black,urlcolor=blue!45!black,
 pdftitle={Corner polyhedra and Schnyder labelings: polynomial exponents and limiting amplitudes},
 pdfauthor={Research note and research addendum},pdfsubject={Two-sided bounds, non-D-finiteness and full asymptotic equivalents for OEIS A377922, A377920, A377921; Part III: a logarithmic cone theorem for walks with a finite internal state}}
\numberwithin{equation}{section}
\newtheorem{theorem}{Theorem}[section]
\newtheorem{proposition}[theorem]{Proposition}
\newtheorem{lemma}[theorem]{Lemma}
\newtheorem{corollary}[theorem]{Corollary}
\theoremstyle{remark}\newtheorem{remark}[theorem]{Remark}
\newcommand{\PP}{\mathbb P}
\newcommand{\EE}{\mathbb E}
\newcommand{\ZZ}{\mathbb Z}
\newcommand{\RR}{\mathbb R}
\newcommand{\Cov}{\operatorname{Cov}}
\newcommand{\Var}{\operatorname{Var}}
\newcommand{\diag}{\operatorname{diag}}
\newcommand{\SE}{\mathsf{SE}}
\newcommand{\e}{\mathrm e}
\newcommand{\oState}{\mathrm o}
\newcommand{\one}{\mathbf 1}
\newcommand{\norm}[1]{\lVert#1\rVert}
\setlist{itemsep=3pt,topsep=5pt}
\setlength{\emergencystretch}{2em}
\allowdisplaybreaks[1]
\raggedbottom
\newcommand{\repo}[1]{\nolinkurl{#1}}
\newenvironment{writenote}{\begin{quote}\small\sloppy\textbf{[write, 2026-10-02]}\ }{\end{quote}}
% Notes added when Part III (bundle Report 112, batch 103) was written
\newenvironment{writenotelog}{\begin{quote}\small\sloppy\textbf{[write, 2026-10-05]}\ }{\end{quote}}
% Part III (Report 112) macros
\newcommand{\gp}{\Gamma_P}\newcommand{\gs}{\Gamma_S}
\newcommand{\ap}{\alpha_P}\newcommand{\as}{\alpha_S}
\newcommand{\rad}{\operatorname{rad}}
\newcommand{\argu}{\operatorname{arg}}
\newenvironment{sourceabstract}[1]{\begin{quote}\small\textbf{Abstract of #1, as delivered.}\ }{\end{quote}}
\title{Corner polyhedra and Schnyder labelings:\\polynomial exponents and limiting amplitudes\\[4pt]
\large OEIS A377922, A377920 and A377921}
\author{Part I: research note; Part II: research addendum; Part III: research report 112\\[2pt]\small merged in ProveIt from two manuscripts (batch 77, sources 45 and 46)\\
\small and bundle Report 112 (batch 103, source 47; Part III)}
\date{2 October 2026; Part III added 5 October 2026}
\begin{document}
\maketitle
\begin{abstract}
Fusy, Narmanli and Schaeffer encode corner polyhedra (A377922) and 3-connected Schnyder
labelings (A377920) as quadrant walks whose steps depend on a two-state parity, and conjectured
equivalents $\kappa\mu^nn^{-\alpha}$ with irrational exponents. Part~I proves two-sided
$\Theta$ bounds with the conjectured exponents, non-D-finiteness of the generating functions of
A377922, A377920 and A377921, and first-threshold inverses with $O(1)$ error. Part~II proves the
full equivalents with finite positive amplitudes, including
$\widetilde s_n\sim(16/19)^3\kappa_S\mu_S^nn^{-\alpha_S}$ for A377921, and a Lambert-$W$ inverse
with a vanishing-width ceiling bracket. No digits of the amplitudes are certified. The two
manuscripts are printed in full; a guide fixes the notation shared between them.

[5 October 2026, batch 103.] Part~III adds bundle Report 112, which proves the exponents of
A377922, A377920 and A377921 again, at the logarithmic scale, as instances of a general cone
theorem for walks with a finite internal state and unbounded steps (Theorem~\ref{cps:log:thm:cone}).
New with it: that theorem, a six-face sleeve injection giving $s_{n-6}\le\widetilde s_n\le s_n$, a
non-D-finiteness criterion that needs only the logarithmic law, a fixed-endpoint sandwich for
A377922 and an explicit exponential-moment certificate.
\end{abstract}
\tableofcontents
\clearpage
\begingroup\sloppy
\section*{Guide to the merged report}\phantomsection\label{cps:sec:guide}
\addcontentsline{toc}{section}{Guide to the merged report}

\subsection*{What the report proves}
Let $p_n$ count corner polyhedra with $n+3$ flats (A377922), $s_n$ count 3-connected Schnyder
labelings with $n$ inner faces (A377920), and $\widetilde s_n$ count rigid orthogonal surfaces
(A377921), all indexed as in Fusy--Narmanli--Schaeffer \cite{FNS}. Put $\mu_P=9/2$,
$\mu_S=16/3$, $\alpha_P=1+\pi/\arccos(9/16)$ and $\alpha_S=1+\pi/\arccos(22/27)$.
\begin{itemize}
\item \textbf{Part~\ref{cps:part:exponents}} (source 45): $p_n=\Theta(\mu_P^nn^{-\alpha_P})$ and
 $s_n=\Theta(\mu_S^nn^{-\alpha_S})$ on the full integer sequence
 (Theorem~\ref{cps:thm:main}), which proves the exponents of \cite[Conjecture 25]{FNS}; the
 exponents are irrational (Lemma~\ref{cps:lem:irrational}); the generating functions of A377922,
 A377920 and A377921 are not D-finite (Corollary~\ref{cps:cor:nonD}); and the first threshold
 index satisfies $N(Y)=(\log Y+\alpha\log\log Y)/\log\mu+O(1)$ (Lemma~\ref{cps:lem:inverse}).
\item \textbf{Part~\ref{cps:part:amplitudes}} (source 46): finite positive $\kappa_P,\kappa_S$ with
 $p_n\sim\kappa_P\mu_P^nn^{-\alpha_P}$, $s_n\sim\kappa_S\mu_S^nn^{-\alpha_S}$ and
 $\widetilde s_n\sim(16/19)^3\kappa_S\mu_S^nn^{-\alpha_S}$
 (Theorem~\ref{cps:amp:thm:main}): the two positive-amplitude equivalents of
 \cite[Conjecture 25]{FNS} and a coefficientwise equivalent for A377921. On the way: a
 boundary-shift comparison (Lemma~\ref{cps:amp:lem:shift}), harmonic and endpoint limits for
 valid regeneration cycles (Theorem~\ref{cps:amp:thm:cyclelimits}), exact counted-time survival
 constants (Theorem~\ref{cps:amp:thm:exacttime}), a uniform killed local limit
 (Lemma~\ref{cps:amp:lem:killed}), fixed-endpoint bridge amplitudes
 (Theorem~\ref{cps:amp:thm:bridge}) with the universal Brownian factor evaluated
 (Section~\ref{cps:amp:sec:Brownian}), and a Lambert-$W_{-1}$ inverse with a vanishing-width
 two-ceiling bracket (Theorem~\ref{cps:amp:thm:inverse}).
\end{itemize}
Part~\ref{cps:part:amplitudes} supersedes none of Part~\ref{cps:part:exponents}: it uses Part I's
models, covariance calculations, cycle buffer, survival bound, free local theorem and uniform
endpoint estimate as inputs, and it does not reprove non-D-finiteness (which needs only the
$\Theta$ bounds and the irrationality of the exponents).

\begin{writenotelog}
\textbf{Part~\ref{cps:part:logcone}} (source 47, bundle Report 112, batch 103) proves
$p_n=\mu_P^nn^{-\alpha_P+o(1)}$, $s_n=\mu_S^nn^{-\alpha_S+o(1)}$ and the same for $\widetilde s_n$,
non-D-finiteness and inverses with error $o(\log\log Y)$ (Theorem~\ref{cps:log:thm:main}): a
second route to Part~\ref{cps:part:exponents}, weaker than Parts~I--II. New to the report: a
logarithmic cone theorem for Markov-additive walks with a finite internal state, unbounded steps
and a uniform exponential moment (Theorem~\ref{cps:log:thm:cone}); the sleeve injection
$s_{n-6}\le\widetilde s_n\le s_n$ (Proposition~\ref{cps:log:prop:sleeve}), which with
Theorem~\ref{cps:thm:main} gives $\widetilde s_n=\Theta(\mu_S^nn^{-\alpha_S})$
(Remark~\ref{cps:log:rem:theta}); a non-D-finiteness criterion that needs only the logarithmic
law (Remark~\ref{cps:log:rem:criterion}); the sandwich $e_{n-1}\le p_n\le e_{n+1}$; the moment
certificate $189/13$; the bounds $4<\alpha_P<5$, $6<\alpha_S<7$; and three references.
Section~\ref{cps:log:sec:guide} maps every result of Report 112 to Parts~I--II.
\end{writenotelog}

\subsection*{Sources and provenance}
Both manuscripts are dated 2 October 2026 and name no ProveIt commit; neither inspected the
repository beyond a GitHub code search (which then returned nothing for these sequences). They
arrived in batch 77 of \texttt{docs/incoming} (arrival commit \texttt{096ee7b87}) and were placed
in \texttt{aa7345800}. Author lines: ``Research note'' (source 45) and ``Research addendum''
(source 46); the PDF metadata of both say ``Research note''. Neither names a tool; the report is
unrefereed and not formalized.

\begin{center}\small
\begin{tabular}{@{}l>{\raggedright}p{0.4\textwidth}ll>{\raggedright\arraybackslash}p{0.2\textwidth}@{}}
\toprule
Source & Archive (wrapper directory, main file) & Pin & Placed & Printed as\\
\midrule
45 & \texttt{oeis-bimodal-cone-research.zip} (\texttt{oeis-bimodal-cone-report/}, \texttt{article.tex}, 15-page PDF) & none & \texttt{aa7345800} & Part I: Sections 1--10, Appendix A\\
46 & \texttt{oeis-bimodal-full-asymptotics.zip} (\texttt{oeis-bimodal-amplitude-report/}, \texttt{addendum.tex}, 14-page PDF) & none & \texttt{aa7345800} & Part II: Sections 11--21\\
\bottomrule
\end{tabular}
\end{center}

\begin{writenotelog}
Source 47 is bundle Report 112, \texttt{Geometric\_OEIS\_\allowbreak Sequences\_Cone\_\allowbreak Exponents\_and\_\allowbreak Inverses\_Source.zip}
(wrapper \texttt{report112/}, main file \texttt{report112.tex}, 19-page PDF), dated 2 October 2026,
author line ``Research report 112'', no pin; arrival commit \texttt{60f54ea06}, placed in
\texttt{9c995cefe} (batch 103), printed as Part~\ref{cps:part:logcone}, Sections 22--32. It was
written without knowledge of sources 45 and 46. Section~\ref{cps:log:sec:provenance} gives the
details.
\end{writenotelog}

\subsection*{How the two sources were merged}
\begin{itemize}
\item Source 46 is an addendum to source 45 and embeds an unchanged copy of all 30 files of
 source 45's package (verified byte-identical at placement). That copy is not shipped; source 45
 is printed once, as Part~\ref{cps:part:exponents}.
\item Part~\ref{cps:part:exponents} keeps source 45's numbering exactly (Sections 1--10,
 Appendix~\ref{cps:app:checks}), because this guide is unnumbered. So the numbers in source 46's
 ``Foundation Lemma 2.1, Lemma 4.1, Section 3, Section 4.1, Proposition 6.1, Proposition 7.2,
 equation (8.6)'' are still correct here; they are now hyperlinks. Source 46's Section $n$ is
 Section $n+10$ here, and its Theorem $n.m$ is Theorem $(n+10).m$.
\item Both manuscripts are printed in full, with their abstracts. Source 46's Section 2
 (Section~\ref{cps:amp:sec:inputs} here) restates Part I's inputs in its own notation; it is
 kept, because the rest of Part~\ref{cps:part:amplitudes} cites its displays. No result is proved
 in both Parts with the same strength: Part II's inverse theorem sharpens Part I's inverse lemma,
 and both are printed.
\item Source 45's appendix of exact checks is placed after Part~\ref{cps:part:amplitudes}, so that
 it stays Appendix~A.
\item Labels carry the prefix \texttt{cps:} in Part I and \texttt{cps:amp:} in Part II. Apart from
 labels, the only changes to the manuscripts' text are: source 46's citations of the
 Foundation, which became cross-references to Part I; its citation key \texttt{DW}, which became
 source 45's \texttt{DWrevisited} (the same paper); and the dated notes marked \textbf{[write]}.
 No statement, proof, symbol or number was changed.
\item One bibliography: source 45's, plus source 46's entry for Ba\~nuelos--Smits. Source 46's
 entries for Fusy--Narmanli--Schaeffer and Denisov--Wachtel are the same papers as source 45's.
 Its \texttt{Foundation} entry is replaced by Part I. Its uncited \texttt{OEIS} entry (A377922,
 A377920, A377921) duplicates source 45's two OEIS entries and is omitted. Source 45's uncited
 entry for Denisov--Wachtel 2015 is kept, as delivered.
\end{itemize}

\begin{writenotelog}
Part~\ref{cps:part:logcone} was added after Part~\ref{cps:part:amplitudes} and before
Appendix~\ref{cps:app:checks}, which keeps its letter. Its labels carry the prefix
\texttt{cps:log:}; its Section 22 is a guide written then, and Report 112's Section $k$ is
Section $k+22$. No earlier number or label changed. Section~\ref{cps:log:sec:printing} lists how
Report 112 is printed and how its bibliography was merged.
\end{writenotelog}

\subsection*{Notation shared and not shared}
The two manuscripts use the same normalizations: $\mu_P,\mu_S,\alpha_P,\alpha_S$, the lattice $L$,
the kernels $K_P,K_S$, the killed transition probabilities $q_n$ and $q_n^*$, the covariances
$\Sigma_P,\Sigma_S$ per counted step (Part I remarks that the matrices of \cite[Section 4.2]{FNS}
are $\mu\Sigma$), the angle $\theta$ and $p=\pi/\theta=\alpha-1$. \emph{No symbol was renamed and no
normalization was changed}; each Part keeps its own letters, and the letters below mean different
things in the two Parts, or in different places of one Part. In both Parts the roman
$\mathrm e$ is the even parity state, not Euler's number.

\begin{table}[ht]
\centering\small
\caption{Clashing letters. ``Reading to avoid'' is the tempting false identification.}\label{cps:tab:notation}
\begin{tabular}{@{}>{$}l<{$}>{\raggedright}p{0.3\textwidth}>{\raggedright}p{0.3\textwidth}>{\raggedright\arraybackslash}p{0.22\textwidth}@{}}
\toprule
\text{Letter} & Part I (source 45) & Part II (source 46) & Reading to avoid\\
\midrule
D & $D=(1-\frac1{3x^2})(1-\frac{y^2}3)$ in \eqref{cps:eq:PK}; $D_0$; operators $D_x=x\partial_x$, $D_y$ & $D$: valid complete-cycle kernel (Section~\ref{cps:amp:sec:kernels}), $D_i$ per parity; Part I's $D$ is $D_P$ & $D^m$ in Part II is not a power of Part I's denominator\\
H & $H=1-x^{-2}-y^2-2x^{-2}y^2$ in \eqref{cps:eq:SM} & $H_S$ = Part I's $H$; $\mathcal H=\lim D^mV$ harmonic function & $\mathcal H$ is not $H$\\
T & return times $T_k$; $T(z)=\sum\widetilde s_nz^n$ (proof of Corollary~\ref{cps:cor:nonD}) & whitening $T=(2\Sigma)^{-1/2}$; $T_k$; $T_Y=\log Y$; Part I's $T(z)$ is $\widetilde S(z)$ & $T$ in Part II is a matrix, not a series\\
K & tilted kernels $K_P,K_S$ & the same kernels; also the wedge $K=TC$, $K_b$, heat kernels $k_t^K$, the semigroup $K_t$ & \\
C & generic constants; conditional covariances $C_i$; a fixed ball $C$ (Section~\ref{cps:sec:gluing}) & open quadrant $C=(0,\infty)^2$; $C_B$; bridge amplitude $\mathcal C(a,b)$; shifts $c_P,c_S$ & \\
B & entry $B=x/y+x^{-1}y/H$; balls $B,B',B_n$ & continuity sets $B$; bridge factor $\mathcal B_\Sigma$; shifts $b_P,b_S$; coefficients $b_j$ & \\
L & lattice $L$; $L=\log Y$ in the proof of Lemma~\ref{cps:lem:inverse} & lattice $L$; $L_m=\sqrt m\log m$ & \\
A & generating function $A(z)$ (Lemma~\ref{cps:lem:nonD}) & survival amplitudes $A(a)$, $A^*(b)$ & \\
S & survival $S_n(a)$; $S(z)=\sum s_nz^n$ & $S(z)$ the same series; $S$ also names the model & \\
\beta & $\beta=k-\alpha$ (proof of Lemma~\ref{cps:lem:nonD}) & survival constants $\beta,\beta^*,\beta_B$ & \\
\lambda & principal eigenvalue $\lambda(t)$; exponential-moment parameter & $\lambda=\log\mu$ (Theorem~\ref{cps:amp:thm:inverse}) & \\
h & corrector $h_{\mathrm e},h_{\mathrm o}$ & corrector; also $h(x)=u(\Sigma^{-1/2}x)$ (Section~\ref{cps:amp:sec:Brownian}) & \\
a & stay probability $a=1/6$ or $1/8$ in \eqref{cps:eq:internal}; also points & points; sequences $a_n$ in Theorem~\ref{cps:amp:thm:inverse} & \\
r,\ g & coordinate index $r$; $r\in\{9/16,22/27\}$ & $r(Y)$ the Lambert centre; $g_t$ entrance density, $g(t)$ in a proof & \\
\gamma,\ \kappa & --- & $\kappa_P,\kappa_S$ amplitudes; $\gamma$ a generic amplitude ($=\kappa$) & \\
\bottomrule
\end{tabular}
\end{table}

\begin{writenotelog}
Part~\ref{cps:part:logcone} has its own table, Table~\ref{cps:log:tab:notation}. The most
dangerous clash: Report 112 works in half-scale coordinates ($x=2U+c$), so its $\Sigma_P$, $\Sigma_S$
and phase covariances are one quarter of the matrices with the same names in Parts~I--II
(Remark~\ref{cps:log:rem:scale}); the angles and exponents agree. Its $\Gamma_P,\Gamma_S$ are
$\mu_P,\mu_S$ and its $\nu$ is $p$.
\end{writenotelog}

\subsection*{What is not claimed}
Every limitation stated by either source is kept in its text. In summary:
\begin{itemize}
\item No numerical digits of $\kappa_P$ or $\kappa_S$, no rate of convergence and no correction
 term are certified (Part II, Sections 11 and 21). The amplitudes are characterized by positive
 killed-harmonic and Brownian heat-kernel expressions; only the universal Brownian factor is
 evaluated.
\item No unconditional rounding formula $N(Y)=\lceil r(Y)\rceil$ (Part II); Part I's inverse has
 $O(1)$ error and no next constant.
\item A377923, whose OEIS entry the delivery also exported, is a different sequence and is
 outside both theorems.
\item No general cone theorem for Markov-modulated walks (Part I, Section 1); the arguments use a
 fixed cycle buffer specific to these two models. No cone-killed local theorem for the two-state
 chain is assumed (Part II).
\item The counting bijections are Fusy--Narmanli--Schaeffer's, the iid cone input is
 Denisov--Wachtel's, and the Brownian cone kernel is Ba\~nuelos--Smits'. Neither source claims
 publication priority; both literature checks (\texttt{45-cone-literature-status.md},
 \texttt{46-amp-literature-status.md}) are bounded snapshots that do not establish novelty.
\item The finite computations check identities, normalizations and indexing (coefficients to
 $n=30$); they do not establish any limit. Neither source was externally peer reviewed or
 machine-formalized; the shipped mathematical reviews are the deliveries' own.
\end{itemize}

\begin{writenotelog}
The bullet ``No general cone theorem for Markov-modulated walks'' is true of
Parts~\ref{cps:part:exponents}--\ref{cps:part:amplitudes}, but no longer of the report:
Part~\ref{cps:part:logcone}'s Theorem~\ref{cps:log:thm:cone} is a general theorem for walks with a
finite internal state, at the logarithmic scale. It now holds only for amplitudes and $\Theta$
bounds: the report has no general amplitude or $\Theta$ theorem (Questions~\ref{cps:log:q:amplitude}
and~\ref{cps:log:q:theta}). Part~\ref{cps:part:logcone} keeps all of Report 112's non-claims: no
amplitude, no full equivalent, no rate, no claim at endpoints without seeds, no $O(1)$ inverse of
its own, and finite checks that prove no limit.
\end{writenotelog}

\subsection*{Inversion mechanics and formal status}
Part I's Lemma~\ref{cps:lem:inverse} and Part II's Theorem~\ref{cps:amp:thm:inverse} are instances of
three theorems of the repository's transseries volume
(\repo{Analysis/Transseries/docs/series-and-transseries/Transseries_And_Inversion/}):
\texttt{p0:thm:lambert-core}, \texttt{p0:thm:lambert-centered} and, for the integer bracket,
\texttt{p0:thm:staircase}. No novelty is claimed for the inversion mechanics. Nothing in this
report is formalized in Lean or Rocq, and its place in the collection confers no formal status.
The only related formal statement is the generic staircase lemma
\texttt{Fabius.staircase\_separation}
(\repo{Analysis/FabiusFunction/Lean/FabiusFunction/StaircaseInversion.lean}), which concerns an
arbitrary monotone interpolation, not these sequences.
\begin{writenotelog}
Part~\ref{cps:part:logcone}'s inverse \eqref{cps:log:eq:inverse-main} is the same instance of
\texttt{p0:thm:lambert-core} at logarithmic precision. Nothing in Part~\ref{cps:part:logcone} is
formalized either.
\end{writenotelog}
\endgroup
\clearpage
\part{Polynomial exponents (source 45)}\label{cps:part:exponents}
\begin{center}\large Polynomial exponents for corner polyhedra and Schnyder labelings\\[2pt]
\normalsize Research note, 2 October 2026\end{center}
\begin{sourceabstract}{source 45}
The quadrant-walk models of Fusy, Narmanli, and Schaeffer have a two-state parity dependence that prevents a direct application of an independent-increment cone theorem at the original time scale. We prove two-sided constant-factor estimates for the two enumerating sequences:
\[
 p_n=\Theta\!\left((9/2)^n n^{-1-\pi/\arccos(9/16)}\right),\qquad
 s_n=\Theta\!\left((16/3)^n n^{-1-\pi/\arccos(22/27)}\right).
\]
The proof regenerates at returns to a parity state. A deterministic one-unit bound on excursions within a regeneration cycle transfers independent-increment cone estimates to the original walk. Exponential concentration of the random clock yields lower bounds at every sufficiently large exact time; a Fourier local limit estimate and an interior gluing argument yield the extra factor $n^{-1}$. The estimates prove the conjectured polynomial exponents, but do not prove convergence to positive asymptotic amplitudes. They also imply non-D-finiteness of the two generating functions and, by a rational generating-function identity, of the related rigid-surface generating function. All symbolic and finite enumeration checks are supplied in a reproducibility package.
\end{sourceabstract}

\section{Results and scope}
Let $p_n$ be the number of polyhedral orientations with $n$ inner vertices, equivalently the number of combinatorial types of corner polyhedra with $n+3$ flats. Let $s_n$ count 3-connected Schnyder labelings with $n$ inner faces in the associated quadrangular dissection. These are the sequences \href{https://oeis.org/A377922}{A377922} and \href{https://oeis.org/A377920}{A377920}, with precisely the indexing of Fusy--Narmanli--Schaeffer \cite{FNS,OEISP,OEISS}. Put
\begin{equation}\label{cps:eq:parameters}
 \mu_P=\frac92,\quad \mu_S=\frac{16}3,\qquad
 \alpha_P=1+\frac{\pi}{\arccos(9/16)},\quad
 \alpha_S=1+\frac{\pi}{\arccos(22/27)}.
\end{equation}
Throughout, $a_n=\Theta(b_n)$ means that positive constants $c,C$ exist such that $cb_n\le a_n\le Cb_n$ for all sufficiently large integers $n$.

\begin{theorem}\label{cps:thm:main}
With the preceding conventions,
\begin{equation}\label{cps:eq:main}
 p_n=\Theta(\mu_P^n n^{-\alpha_P}),\qquad
 s_n=\Theta(\mu_S^n n^{-\alpha_S}).
\end{equation}
Both estimates hold on the full integer sequence, not merely on a congruence subsequence.
\end{theorem}

\begin{center}
\begin{tabular}{@{}llcc@{}}
\toprule
Sequence & Objects & Exponential rate & Polynomial exponent\\
\midrule
A377922 & Corner polyhedra & $9/2$ & $\alpha_P=4.227476082\ldots$\\
A377920 & Schnyder labelings & $16/3$ & $\alpha_S=6.080304846\ldots$\\
\bottomrule
\end{tabular}
\end{center}

The exponential rates were proved in \cite[Theorem 24]{FNS}. Their Conjecture 25 predicts asymptotic equivalents with positive amplitudes and the exponents in \eqref{cps:eq:parameters}. Theorem~\ref{cps:thm:main} establishes the exponents and stronger two-sided bounds, but leaves the amplitude limits unresolved. In particular,
\begin{equation}\label{cps:eq:loglimits}
 \lim_{n\to\infty}\frac{\log p_n-n\log\mu_P}{\log n}=-\alpha_P,
 \qquad
 \lim_{n\to\infty}\frac{\log s_n-n\log\mu_S}{\log n}=-\alpha_S.
\end{equation}
The manuscript relies on the published combinatorial bijections and on the independent-increment cone theorem stated in Section~\ref{cps:sec:iid}. It does not posit a new general cone theorem for arbitrary Markov-modulated walks. Its argument is specific to a geometric property of these two tandem-walk models.

\begin{writenotelog}
This disclaimer describes source 45 alone. Part~\ref{cps:part:logcone} (bundle Report 112) proves
a general cone theorem for walks with a finite internal state, Theorem~\ref{cps:log:thm:cone}:
under hypotheses (H1)--(H5), which allow unbounded steps with a uniform exponential moment and need
no cycle buffer, the killed fixed-endpoint kernel is $n^{-1-\pi/\theta+o(1)}$ at every integer
time. It is a logarithmic-scale statement; for amplitudes and for $\Theta$ bounds the disclaimer
still holds (Questions~\ref{cps:log:q:amplitude}--\ref{cps:log:q:theta}).
\end{writenotelog}

\begin{writenote}
The amplitude limits that this paragraph leaves unresolved are established in
Part~\ref{cps:part:amplitudes} (source 46): Theorem~\ref{cps:amp:thm:main} proves both
positive-amplitude equivalents of \cite[Conjecture 25]{FNS}, on the full integer sequence, and
an equivalent for A377921. Part~\ref{cps:part:amplitudes} characterizes $\kappa_P,\kappa_S$ by
positive harmonic and heat-kernel expressions and certifies no numerical digits of them.
\end{writenote}

\section{Exact walk models and normalization}\label{cps:sec:models}
A tandem step is either $\SE=(1,-1)$ or a \emph{face step} $(-i,j)$ with $i,j\ge0$. A walk is even when every increment has coordinates of the same parity. The common state space is
\begin{equation}\label{cps:eq:lattice}
 L=\{(a,b)\in\ZZ^2:a\equiv b\pmod2\}.
\end{equation}
The internal state $J$ is the parity of the second coordinate, denoted $\e$ or $\oState$. In a matrix of Laurent series, rows give the initial state, columns the final state, and the coefficient of $x^ay^b$ gives the weight of displacement $(a,b)$. Nonnegative coefficients are understood in the indicated convergent expansions.

\subsection{Corner polyhedra}
The $P$-admissible steps are
\begin{align}
 \mathcal S_P^{\e}
 &=\{(1,-1)\}\cup\{(-2i,2j):i\ge0,j\ge1\}
       \cup\{(-2i-1,2j+1):i,j\ge0\},\label{cps:eq:Psteps}\\
 \mathcal S_P^{\oState}
 &=\{(1,-1)\}\cup\{(-2i,2j):i\ge1,j\ge0\}
       \cup\{(-2i-1,2j+1):i,j\ge0\}.\notag
\end{align}
Every step has counting weight one. Proposition 13 of \cite{FNS} identifies $p_n$ with the number of $P$-admissible walks of length $n$ from $(0,0)$ to the horizontal axis, staying in the nonnegative quadrant. A nonempty such walk cannot return to $(0,0)$: its final step would have to be horizontal from a point on the horizontal axis, whereas no such face step is permitted at even height. Therefore the endpoints are $(2i,0)$, $i\ge1$.

Write $D_0=(1-x^{-2})(1-y^2)$. The unweighted matrix is
\begin{equation}
 M_P(x,y)=\begin{pmatrix}
 y^2/D_0&x/y+x^{-1}y/D_0\\
 x/y+x^{-1}y/D_0&x^{-2}/D_0
 \end{pmatrix}.
\end{equation}
Tilting and normalizing gives
\begin{equation}\label{cps:eq:PK}
 K_P(x,y)=\frac2{9}M_P(\sqrt3x,y/\sqrt3)
 =\begin{pmatrix}
 \dfrac{2y^2}{27D}&\dfrac{2x}{3y}+\dfrac{2y}{27xD}\\[5pt]
 \dfrac{2x}{3y}+\dfrac{2y}{27xD}&\dfrac{2}{27x^2D}
 \end{pmatrix},\quad
 D=\left(1-\frac1{3x^2}\right)\left(1-\frac{y^2}3\right).
\end{equation}
At $(1,1)$ the matrix is stochastic. Each admissible displacement $z$ has probability $\mu_P^{-1}3^{(z_1-z_2)/2}$. Let
\begin{equation}
 Q_P=L\cap\ZZ_{\ge0}^2,\qquad
 q_n^P(a,b)=\PP_a(X_n=b,\ X_t\in Q_P\text{ for }0\le t\le n).
\end{equation}
The tilt telescopes along a path, giving the exact identity
\begin{equation}\label{cps:eq:Pcount}
 p_n=\mu_P^n\sum_{i\ge1}3^{-i}q_n^P((0,0),(2i,0)).
\end{equation}

\subsection{Schnyder labelings}
An $S$-admissible face step is one of the following, with $\ell,r\ge0$ and weight $\binom{\ell+r}{r}$:
\begin{equation}\label{cps:eq:Ssteps}
 \begin{array}{c|c}
 \text{initial parity}&\text{face displacement}\\ \hline
 \text{either}&(-2\ell-2,\ 2r+2)\\
 \e&(-2\ell-1,\ 2r+3)\\
 \oState&(-2\ell-3,\ 2r+1).
 \end{array}
\end{equation}
Weights count distinct decorations. Write $s'_n$ for the subclass in which the two outer vertices labeled $G$ are non-isolated. Proposition 14 and Remark 9 of \cite{FNS} state that $s'_n$ counts weighted quadrant walks from $(0,2)$ to $(2,0)$ with $n+2$ southeast steps, and that
\begin{equation}\label{cps:eq:Sconversion}
 s_n=s'_n+2s'_{n-1}+s'_{n-2}\qquad(n\ge4).
\end{equation}
Each such walk begins and ends with a southeast step. Removing the last one leaves a walk from $(0,2)$ to $(1,1)$ with $n+1$ southeast steps.

An \emph{aggregate} consists of one southeast step followed by an arbitrary, possibly empty, string of face steps. The face-step matrix and southeast matrix are
\begin{equation}
 W(x,y)=\frac1{1-x^{-2}-y^2}
 \begin{pmatrix}x^{-2}y^2&x^{-1}y^3\\x^{-3}y&x^{-2}y^2\end{pmatrix},
 \qquad E(x,y)=\frac{x}{y}\begin{pmatrix}0&1\\1&0\end{pmatrix}.
\end{equation}
Thus $M_S=E(I-W)^{-1}$. With
\begin{equation}
 H=1-x^{-2}-y^2-2x^{-2}y^2,\qquad B=x/y+x^{-1}y/H,
\end{equation}
explicit inversion yields
\begin{equation}\label{cps:eq:SM}
 M_S(x,y)=\begin{pmatrix}x^{-2}/H&B\\B&y^2/H\end{pmatrix},\qquad
 K_S(x,y)=\frac3{16}M_S(2x,y/2).
\end{equation}
The expansion of $(I-W)^{-1}$ accounts for every decorated face string with its correct multiplicity. It converges at the tilting point $(2,1/2)$, and in a neighborhood thereof.

\begin{lemma}[Exact aggregate boundary]\label{cps:lem:Sboundary}
For an aggregate walk from $(0,2)$ to $(1,1)$, every underlying physical step stays in the nonnegative quadrant if and only if every aggregate endpoint stays in
\begin{equation}\label{cps:eq:SQ}
 Q_S=\{(a,b)\in L:a\ge0,\ b\ge1\}.
\end{equation}
\end{lemma}
\begin{proof}
If an aggregate begins at height at least one, its southeast step remains nonnegative. Subsequent face steps decrease the first coordinate and increase the second, so nonnegativity of the aggregate's final first coordinate controls the whole face string. Conversely, any aggregate except the last must end at height at least one to allow the following southeast step; the prescribed last endpoint has height one. The first coordinate is nonnegative at all endpoints whenever the physical path is valid.
\end{proof}
Let $q_N^S$ denote the $K_S$ chain killed outside $Q_S$, with $N$ counting aggregates. A decorated path of $N$ aggregates with displacement $(a,b)$ has probability $\mu_S^{-N}2^{a-b}$. The total displacement here is $(1,-1)$; hence
\begin{equation}\label{cps:eq:Scount}
 s'_n=\frac{\mu_S^{n+1}}4q_{n+1}^S((0,2),(1,1)).
\end{equation}
There is no physical-time offset in \eqref{cps:eq:Pcount}; the $n+1$ in \eqref{cps:eq:Scount} is essential.

\section{Correctors and regenerative covariance}\label{cps:sec:moments}
In this section $K$ denotes either $K_P$ or $K_S$. Time means a physical step for $P$ and an aggregate for $S$. In both cases the internal transition matrix is
\begin{equation}\label{cps:eq:internal}
 P=K(1,1)=\begin{pmatrix}a&1-a\\1-a&a\end{pmatrix},\qquad
 a=\begin{cases}1/6&(P),\\1/8&(S).
 \end{cases}
\end{equation}
Its stationary distribution is $(1/2,1/2)$. Define $D_x=x\partial_x$ and $D_y=y\partial_y$. For example, the conditional mean of coordinate $r$ in state $i$ is
\begin{equation}
 d_{i,r}=\left.\sum_jD_rK_{ij}\right|_{x=y=1}.
\end{equation}
The means satisfy $d_{\oState}=-d_{\e}$. Solving
$d_i+\sum_jP_{ij}(h_j-h_i)=0$ gives a bounded corrector with $h_{\oState}=-h_{\e}$:
\begin{equation}\label{cps:eq:drifts}
 \begin{array}{c|cc}
 &d_{\e}&h_{\e}\\ \hline
 P&(1/6,1/6)&(1/10,1/10)\\
 S&(-1/8,-1/8)&(-1/14,-1/14).
 \end{array}
\end{equation}
Consequently $M_t=X_t+h_{J_t}$ is a martingale.

The conditional covariance matrix of its increment in state $i$ has entries
\begin{align}\label{cps:eq:Cformula}
 (C_i)_{rs}={}&\sum_j\bigl[ D_rD_sK_{ij}
 +(h_{j,r}-h_{i,r})D_sK_{ij}
 +(h_{j,s}-h_{i,s})D_rK_{ij}\notag\\
 &\hspace{45mm}+(h_{j,r}-h_{i,r})(h_{j,s}-h_{i,s})K_{ij}\bigr]_{x=y=1}.
\end{align}
Evaluation gives
\begin{align}
 C^P_{\e}&=\begin{pmatrix}12/5&-9/5\\-9/5&4\end{pmatrix},&
 C^P_{\oState}&=\begin{pmatrix}4&-9/5\\-9/5&12/5\end{pmatrix},\label{cps:eq:PC}\\
 C^S_{\e}&=\begin{pmatrix}6&-88/21\\-88/21&30/7\end{pmatrix},&
 C^S_{\oState}&=\begin{pmatrix}30/7&-88/21\\-88/21&6\end{pmatrix}.\label{cps:eq:SC}
\end{align}
Averaging over the stationary internal chain gives the covariance per unit time:
\begin{equation}\label{cps:eq:Sigma}
 \Sigma_P=\begin{pmatrix}16/5&-9/5\\-9/5&16/5\end{pmatrix},\qquad
 \Sigma_S=\begin{pmatrix}36/7&-88/21\\-88/21&36/7\end{pmatrix}.
\end{equation}
Both are positive definite. The matrices displayed in \cite[Section 4.2]{FNS} are $\mu_P\Sigma_P$ and $\mu_S\Sigma_S$, respectively; the stochastic per-counted-step normalization here differs by these scalar factors. The correlation coefficients and conjectured powers are unchanged. The unequal matrices in \eqref{cps:eq:PC} and \eqref{cps:eq:SC} are one reason to avoid imposing a state-independent conditional covariance hypothesis on the original chain.

\subsection{First returns to a parity state}
Fix the initial parity $i$ and set $T_0=0$ and $T_{k+1}=\inf\{t>T_k:J_t=i\}$. By the Markov property and translation invariance within a parity class, the pairs
\begin{equation}
 (Y_k,\ell_k)=(X_{T_k}-X_{T_{k-1}},\ T_k-T_{k-1})
\end{equation}
are independent and identically distributed. With $j\ne i$ their joint generating function is
\begin{equation}\label{cps:eq:return}
 R_i(u,x,y)=uK_{ii}(x,y)+\frac{u^2K_{ij}(x,y)K_{ji}(x,y)}{1-uK_{jj}(x,y)}.
\end{equation}
The geometric term enumerates a first transition to $j$, any number of stays in $j$, and a last transition to $i$. The denominator is nonzero at $(1,1,1)$, and all coefficients are nonnegative. Analyticity there gives a joint exponential moment of $(\ell,Y)$.

Differentiating \eqref{cps:eq:return}, for either choice of $i$, yields
\begin{equation}\label{cps:eq:cyclemoments}
 \EE\ell=2,\qquad \EE Y=0,\qquad \Cov(Y)=2\Sigma,
 \qquad
 \Var(\ell)=\begin{cases}2/5&(P),\\2/7&(S).
 \end{cases}
\end{equation}
For instance, $\EE\ell=\partial_u R_i(1,1,1)$ and
$\EE(Y_rY_s)=D_rD_sR_i(1,1,1)$. For returns to the even state,
\begin{equation}
 \Cov(\ell,Y)=\begin{cases}(-2/5,-2/5)&(P),\\(2/7,2/7)&(S).
 \end{cases}
\end{equation}
Thus cycle displacement and duration are genuinely dependent. No argument below requires their independence or conditions on an exact total cycle duration.

\subsection{Time reversal and cone angle}
Write $K_{ij}(z)$ for a coefficient of the step matrix. The reverse transition probabilities are
\begin{equation}\label{cps:eq:reverse}
 K^*_{ij}(z)=K_{ji}(-z).
\end{equation}
For a fixed endpoint in $L$, the total incoming mass is the relevant column sum of $P$, which is one. Counting measure on $L$ is therefore invariant, and $K^*$ is stochastic. For either killed domain,
\begin{equation}\label{cps:eq:killedreverse}
 q_n(a,b)=q_n^*(b,a).
\end{equation}
The matrices $K_P$ and $K_S$ are symmetric as matrices of Laurent series, so $K^*(x,y)=K(x^{-1},y^{-1})$. The reverse corrector is $h^*=-h$, and its corrected conditional covariance matrices are the same as \eqref{cps:eq:PC}--\eqref{cps:eq:SC}. Reversing a first-return path preserves its probability and duration and negates its displacement. The reverse cycle law consequently has the same covariance as the forward one.

Whitening covariance $2\Sigma$ transforms the positive quadrant to a wedge with angle
\begin{equation}\label{cps:eq:angles}
 \theta=\arccos\!\left(-\frac{\Sigma_{12}}{\sqrt{\Sigma_{11}\Sigma_{22}}}\right),
 \qquad p=\frac{\pi}{\theta},
 \qquad
 \theta_P=\arccos(9/16),\quad \theta_S=\arccos(22/27).
\end{equation}
Thus $\alpha=1+p$ in each model. Multiplying the covariance by a scalar does not change this angle.

\section{A fixed buffer for regeneration cycles}\label{cps:sec:buffer}
Call an increment northwest if its first coordinate is nonpositive and its second coordinate is nonnegative. For $P$, every parity-preserving step is northwest, and every parity-changing step is either northwest or $\SE$. The same assertion holds for $S$ aggregates: any nonempty face string following the initial southeast step offsets its increase of the first coordinate and decrease of the second, since each face step has first coordinate at most $-1$ and second coordinate at least $1$. An empty aggregate is exactly $\SE$, and cannot preserve parity.

\begin{lemma}[One-unit cycle bound]\label{cps:lem:buffer}
For a first return to either parity in either model, with endpoints $a,b$, every intermediate position $z$ at the model's time scale satisfies
\begin{equation}\label{cps:eq:buffer}
 z_r\ge\min(a_r,b_r)-1\qquad(r=1,2).
\end{equation}
The same statement holds for reversed cycles.
\end{lemma}
\begin{proof}
A one-step return is northwest. Otherwise the cycle consists of one parity-changing step, a string of parity-preserving northwest steps, and one parity-changing step. The first and last steps may each be northwest or $\SE$. A drop of the second coordinate below the smaller endpoint height is possible only at the first step, when that step is $\SE$, and its size is one. A drop of the first coordinate below the smaller endpoint first coordinate is possible only immediately before a final $\SE$, also by one. All intervening changes have the northwest monotonicity. Reversal preserves the set of intermediate positions and exchanges the endpoints, so it preserves the bound.
\end{proof}
If all even-parity cycle endpoints have both coordinates at least two, every intervening position has both coordinates at least one. This ensures survival in $Q_P$ and $Q_S$. In the $S$ model Lemma~\ref{cps:lem:Sboundary} then ensures survival at every underlying physical step as well. The buffer is fixed, irrespective of cycle duration or displacement.

\subsection{Accessibility and finite entrance paths}\label{cps:sec:entrances}
The even-return cycle laws have explicit increments that make the required cone starting points accessible. In $P$, $(0,2)$ is a one-step even return, and $(2,-2)$ is realized by two southeast steps. A cycle with displacement $(-2,0)$ consists of
\begin{equation}
 (-1,1),\quad(-2,0),\quad(1,-1),
\end{equation}
with states $\e,\oState,\oState,\e$. The reverse cycle law therefore contains $(2,0)$ and $(-2,2)$. By raising the second coordinate before adding $(2,-2)$, the forward skeleton reaches arbitrarily deep diagonal points without exiting the translated cone. The analogous reverse construction first raises the first coordinate and then uses $(-2,2)$.

In $S$, two empty aggregates give $(2,-2)$. A cycle of aggregate displacements
\begin{equation}\label{cps:eq:Saccess}
 (1,-1),\quad(0,2),\quad(-1,1)
\end{equation}
gives $(0,2)$. The middle aggregate is a southeast step followed by the face step $(-1,3)$; the last is a southeast step followed by $(-2,2)$. A single aggregate consisting of a southeast step and $(-3,1)$ gives the even return $(-2,0)$. Hence the same forward and reverse accessibility argument applies.

Finally, there are positive-probability entrance paths to $(2,2)$. For $P$ they are
\begin{equation}\label{cps:eq:Pentrances}
 (0,0)\longrightarrow(0,4)\longrightarrow(1,3)\longrightarrow(2,2),
 \qquad
 (4,0)\xrightarrow{\text{reverse}}(3,1)\longrightarrow(2,2).
\end{equation}
For $S$ the forward aggregate path is
\begin{equation}\label{cps:eq:Sentrances}
 (0,2)\xrightarrow{\SE}(1,1)
 \xrightarrow{\SE+(-1,3)}(1,3)
 \xrightarrow{\SE}(2,2).
\end{equation}
Its middle aggregate passes through $(2,0)$, which is permitted physically. The reverse entrance is
\begin{equation}
 (1,1)\xrightarrow{\text{reverse }\SE}(0,2)
 \xrightarrow{(2,0)}(2,2).
\end{equation}
The last reversed aggregate is the reversal of the valid forward path
$(2,2)\to(3,1)\to(0,2)$. All displayed entrances respect their killed domains.

\section{Independent-increment cone input}\label{cps:sec:iid}
We use the following consequences of Denisov--Wachtel \cite[Theorems 2 and 3]{DWrevisited}. Let $Z_k$ be an iid centered random walk in a planar wedge, after a nonsingular linear normalization of its covariance to the identity. Assume a nondegenerate covariance and an exponential moment. Let $\tau$ be its first exit from the open wedge and let $p=\pi/\theta$, where $\theta$ is the wedge angle. Assume $p\ge1$, as holds for both wedges here. Then
\begin{equation}\label{cps:eq:iidupper}
 \PP_z(\tau>k)\le C(1+\norm z)^p k^{-p/2}.
\end{equation}
For a fixed starting point with positive killed harmonic function and any fixed ball $B$ compactly contained in the wedge,
\begin{equation}\label{cps:eq:iidlower}
 \PP_z(\tau>k,\ Z_k\in\sqrt{k}B)\sim c_{z,B}k^{-p/2},
 \qquad c_{z,B}>0.
\end{equation}
Indeed, Theorem 3(a) bounds survival by a constant times the wedge harmonic function at a fixed inward shift, divided by $k^{p/2}$. The harmonic function is homogeneous of degree $p$ and is bounded above by a constant times $\norm z^p$, giving \eqref{cps:eq:iidupper}. Theorem 3(b) gives a positive integral over the ball in \eqref{cps:eq:iidlower}; Theorem 2 and accessibility give positivity of the starting harmonic function.

All hypotheses can be checked for the present skeletons. The moment and covariance conditions follow from \eqref{cps:eq:cyclemoments}. A planar wedge is Lipschitz and star-like. In polar coordinates its positive Brownian harmonic function is
\begin{equation}
 u(r,\varphi)=r^p\sin(p\varphi),\qquad0<\varphi<\theta,
\end{equation}
which is comparable to $r^{p-1}$ times distance to the boundary. The accessibility constructions in Section~\ref{cps:sec:entrances} supply the positive starting points used in the lower bound.

The lattice boundary convention causes no difficulty, but must be fixed before whitening. On the even lattice, translating $(a,b)$ to $(a+1,b+1)$ makes $a,b\ge0$ equivalent to positivity of both translated coordinates. Translating to $(a-1,b-1)$ does the same for $a,b\ge2$. For $Q_S$, translation by $(1,0)$ makes $a\ge0,b\ge1$ equivalent to an open quadrant condition on either parity class. Each is a fixed shift. After diffusive rescaling, the induced displacement of endpoint balls tends to zero, and nested fixed balls suffice in every application of \eqref{cps:eq:iidlower}.

\section{Transfer to exact model time}\label{cps:sec:time}
This section applies to both models and their reversals. For $P$, time is physical time; for $S$, it is aggregate time. Put
\begin{equation}
 S_n(a)=\sum_{b\in Q}q_n(a,b)=\PP_a(X_t\in Q\text{ for }0\le t\le n).
\end{equation}

\begin{proposition}[Survival upper bound]\label{cps:prop:survival}
For the starting parity being used for regeneration,
\begin{equation}\label{cps:eq:survival}
 S_n(a)\le C(1+\norm a)^p n^{-p/2}+Ce^{-cn}.
\end{equation}
The constants can be chosen for all starting points in the relevant domain and either of the two parity classes.
\end{proposition}
\begin{proof}
Take $k=\lfloor n/4\rfloor$. Since the iid durations have mean two and an exponential moment, a Chernoff bound gives $\PP(T_k>n)\le Ce^{-cn}$. If the model-time path survives through $n$ and $T_k\le n$, then all of the first $k$ skeleton endpoints lie in $Q$. Drop all within-cycle restrictions and apply \eqref{cps:eq:iidupper}, with the appropriate fixed translation. This is an event inclusion, not an identification of skeleton survival with full path survival. There are only two parities, so their constants can be combined.
\end{proof}

\begin{proposition}[Interior mass at every exact time]\label{cps:prop:mass}
For each ball $B$ with closure in the positive quadrant, a constant $c_B>0$ exists such that
\begin{equation}\label{cps:eq:mass}
 \PP_{(2,2)}(X_t\in Q\ (0\le t\le n),\ X_n\in\sqrt nB)
 \ge c_Bn^{-p/2}
\end{equation}
for all sufficiently large integers $n$. This is true for the forward and reverse chains in both models.
\end{proposition}
\begin{proof}
Choose a smaller ball $B'$ with $\overline{B'}\subset B$ and set
\begin{equation}\label{cps:eq:kchoice}
 k=\left\lceil n/2+n^{3/4}\right\rceil.
\end{equation}
Constrain the even skeleton endpoints to have both coordinates at least two. Equation~\eqref{cps:eq:iidlower}, after the fixed translation, gives
\begin{equation}\label{cps:eq:skeletonmass}
 \PP\bigl(X_{T_j}\in[2,\infty)^2\ (0\le j\le k),\
 X_{T_k}\in\sqrt nB'\bigr)\ge cn^{-p/2}.
\end{equation}
To justify the varying normalization, note that $k/n\to1/2$. A sufficiently small fixed ball whose closure lies in $\sqrt2B'$ is contained in the rescaled target for all large $n$, even after the fixed lattice shift. Thus only a fixed-ball application of \eqref{cps:eq:iidlower} is needed.

The mean of $T_k$ is $2k=n+2n^{3/4}+O(1)$. Exponential moments of the durations give
\begin{equation}\label{cps:eq:clock}
 \PP\{T_k\notin[n,n+4n^{3/4}]\}\le Ce^{-cn^{1/2}}.
\end{equation}
For completeness, the usual moment-generating-function bound for a centered duration is $\EE e^{\lambda(\ell-2)}\le e^{C\lambda^2}$ for sufficiently small $|\lambda|$. Choosing $\lambda$ proportional to $n^{-1/4}$ in a deviation of order $n^{3/4}$ over $k=O(n)$ terms proves \eqref{cps:eq:clock}.

For the corrected model-time martingale, the conditional exponential moments are uniformly bounded over the two states. Since each increment has conditional mean zero, Taylor expansion gives
\begin{equation}
 \EE\bigl(e^{\lambda\Delta M_t^{(r)}}\mid\mathcal F_t\bigr)
 \le e^{C\lambda^2}\qquad(|\lambda|\le\lambda_0,\ r=1,2).
\end{equation}
The exponential supermartingale maximal inequality, applied to both signs and both coordinates, implies for fixed $\delta>0$
\begin{equation}\label{cps:eq:window}
 \PP\!\left(\max_{0\le r\le\lceil4n^{3/4}\rceil+2}
 \norm{X_{n+r}-X_n}>\delta\sqrt n\right)
 \le Ce^{-c_\delta n^{1/4}}.
\end{equation}
One can take $\lambda$ proportional to $\sqrt n/n^{3/4}=n^{-1/4}$; the bounded corrector changes thresholds by only a constant. This estimate is unconditional and uniform in the state at time $n$. The reverse corrector verified above gives the identical argument in reverse time.

On the event in \eqref{cps:eq:skeletonmass}, Lemma~\ref{cps:lem:buffer} gives survival through $T_k$. Remove the bad events in \eqref{cps:eq:clock} and \eqref{cps:eq:window}. If $\delta$ is smaller than the separation of $\overline{B'}$ from $B^c$, then $X_n\in\sqrt nB$ and the path survives through the exact time $n$. Both subtracted errors are exponentially smaller than $n^{-p/2}$. This proves \eqref{cps:eq:mass} without any conditioning of clock or displacement estimates on survival.
\end{proof}
The entrance paths in Section~\ref{cps:sec:entrances} extend \eqref{cps:eq:mass} to the fixed endpoints needed below. A smaller target ball absorbs their fixed time lengths. Explicitly, for $P$ the forward starting point is $(0,0)$ and the reverse starting point is $(4,0)$; for $S$ they are $(0,2)$ and $(1,1)$, respectively. The resulting endpoint masses are each bounded below by a positive constant times $n^{-p/2}$.

\section{Free local bounds and lattice aperiodicity}\label{cps:sec:Fourier}
Represent every $z\in L$ uniquely as
\begin{equation}
 z=(2a+j,2b+j),\qquad(a,b)\in\ZZ^2,\quad j\in\{0,1\}.
\end{equation}
This is a two-state translation-cell representation. Let $F(t)$ be the Fourier matrix of its cell increments, with $t\in[-\pi,\pi]^2$. If
$D(t)=\diag(1,e^{i(t_1+t_2)/2})$, then
\begin{equation}\label{cps:eq:cellFourier}
 F(t)=D(t)K(e^{it_1/2},e^{it_2/2})D(t)^{-1}.
\end{equation}
Indeed, a transition from internal state $i$ to $j$ with physical displacement $z$ has cell displacement $(z-(j-i)(1,1))/2$. The direct cell formula also makes $F$ $2\pi$-periodic in each coordinate.

\begin{lemma}[No secondary Fourier points]\label{cps:lem:aperiodic}
The only point of the cell torus where $F(t)$ has spectral radius one is $t=0$. At zero the sole unit-modulus eigenvalue is the simple eigenvalue one.
\end{lemma}
\begin{proof}
Suppose $F(t)v=\lambda v$ with $|\lambda|=1$, and normalize the largest component modulus of $v$ to one. Equality in the triangle inequality, together with irreducibility of the internal stochastic matrix, forces both component moduli to equal one and all supported transition phases to align. In particular, all $\e\to\e$ cell-increment phases are equal. For $P$ their support includes every $(-i,j)$ with $i\ge0,j\ge1$. Comparing $(0,1)$ and $(0,2)$ forces $t_2=0$ modulo $2\pi$; comparing $(0,1)$ and $(-1,1)$ forces $t_1=0$. For $S$, the support includes every $(-i,j)$ with $i\ge1,j\ge0$. Comparing $(-1,0),(-2,0),(-1,1)$ gives the same conclusion. At zero, the matrix in \eqref{cps:eq:internal} is primitive; its other eigenvalue is $2a-1$, strictly inside the unit circle.
\end{proof}

\begin{proposition}[Free local estimates]\label{cps:prop:local}
For either chain, with no killing,
\begin{equation}\label{cps:eq:freeupper}
 \sup_{a,b\in L}\PP_a(X_n=b)\le Cn^{-1}.
\end{equation}
For every fixed $A<\infty$, uniformly for $a,b\in L$ with $\norm{b-a}\le A\sqrt n$,
\begin{equation}\label{cps:eq:LLT}
 \PP_a(X_n=b)=\frac{2}{2\pi n\sqrt{\det\Sigma}}
 \exp\!\left(-\frac{(b-a)^T\Sigma^{-1}(b-a)}{2n}\right)+o(n^{-1}).
\end{equation}
The statements hold for all four initial and final parity pairs, and for the reverse chains.
\end{proposition}
\begin{proof}
Near zero the principal eigenvalue and spectral projection of $F(t)$ are analytic. Centering and \eqref{cps:eq:Sigma} give
\begin{equation}\label{cps:eq:lambda}
 \lambda(t)=1-\tfrac12t^T(\Sigma/4)t+O(\norm t^3),
 \qquad\Pi(t)\longrightarrow\one(1/2,1/2).
\end{equation}
To identify the Hessian without presuming an iid walk at model time, write $X_n=M_n-h_{J_n}$. Martingale increment orthogonality and convergence of the finite internal chain give $n^{-1}\Cov(M_n)\to\Sigma$; the bounded corrector changes the covariance divided by $n$ by $o(1)$. Analytic differentiation of the finite matrix moment-generating function identifies the same limiting covariance with the Hessian of the log principal eigenvalue. Passing to cells divides it by four. The mean is zero at linear scale for the same reason.

In a small neighborhood of zero, $|\lambda(t)|\le e^{-c\norm t^2}$ and the other spectral part decays exponentially in $n$. On the compact complement, Lemma~\ref{cps:lem:aperiodic} gives a uniform strict spectral-radius bound, and finite-dimensional matrix-power estimates give exponential decay. Fourier inversion therefore proves \eqref{cps:eq:freeupper}. Rescaling $t$ by $n^{-1/2}$ in the neighborhood of zero gives the Gaussian local limit, uniformly on bounded diffusive displacement ranges. The cell covariance contributes a factor four when expressed in physical coordinates; the terminal parity mass is $1/2$. Their product is the factor two in \eqref{cps:eq:LLT}, also the covolume of $L$ in $\ZZ^2$. Bounded parity-coordinate offsets change the exponent by $o(1)$. Reversal preserves the same argument.
\end{proof}
This Fourier calculation supplies the exact-time aperiodicity that is not supplied merely by an invariance principle. There is no leftover restriction to even times.

\section{A killed interior estimate and gluing}\label{cps:sec:gluing}
\begin{lemma}[Interior killed local lower bound]\label{cps:lem:killed}
Fix $\eta>0$. For a sufficiently large fixed $R$, set
\begin{equation}
 B=B((R,R),\eta),\qquad B_n=\sqrt nB\cap L.
\end{equation}
There is a constant $c>0$ such that, for all sufficiently large $n$ and every integer $m\in[n/4,n/2]$,
\begin{equation}\label{cps:eq:killedlocal}
 q_m(a,b)\ge c/n\qquad(a,b\in B_n).
\end{equation}
\end{lemma}
\begin{proof}
Equation~\eqref{cps:eq:LLT} gives $\PP_a(X_m=b)\ge c_0/n$ uniformly in the specified range. Here $c_0$ is independent of $R$, since the Gaussian depends on $b-a$, whose norm is at most $2\eta\sqrt n$.

Let $\tau$ be the first exit from $Q$. For paths that exit in the first half of the interval, the strong Markov property at $\tau$ and \eqref{cps:eq:freeupper} give
\begin{equation}
 \PP_a(X_m=b,\ \tau\le\lfloor m/2\rfloor)
 \le\frac Cn\PP_a(\tau\le\lfloor m/2\rfloor).
\end{equation}
Every starting point in $B_n$ is at coordinate distance at least $(R-\eta)\sqrt n-O(1)$ from the complement of $Q$. The bounded corrector and Doob's quadratic martingale maximal inequality yield
\begin{equation}
 \PP_a(\tau\le\lfloor m/2\rfloor)\le\frac C{(R-\eta)^2}.
\end{equation}
The expected martingale quadratic variation is $O(m)$ by the finite matrices \eqref{cps:eq:PC}--\eqref{cps:eq:SC}. For a path that first exits in the second half, reverse the full unrestricted path from its fixed endpoint $b$. Some reversed position outside $Q$ is then visited in the first half of reversed time. Equation~\eqref{cps:eq:reverse} and the same argument bound this mass by $C/[n(R-\eta)^2]$. Taking a union bound covers all paths with any exit. Choosing $R$ large makes this smaller than $c_0/(2n)$; subtracting from the free lower bound proves the claim. The fixed shift of $Q_S$ changes only the harmless $O(1)$ boundary distance.
\end{proof}
No cone-killed local limit theorem or conditioned bridge theorem is used in this step.

\subsection{Upper bounds}
Split $n=m+\ell+m$, where $m=\lfloor n/3\rfloor$ and $\ell=n-2m$. The Markov property, the unrestricted middle-kernel upper bound, and reversal give
\begin{align}\label{cps:eq:threeupper}
 q_n(a,b)
 &=\sum_{u,v\in Q}q_m(a,u)q_\ell(u,v)q_m(v,b)\notag\\
 &\le\frac Cn S_m(a)S_m^*(b).
\end{align}
For $P$, Proposition~\ref{cps:prop:survival} gives
\begin{equation}\label{cps:eq:Pupper}
 q_n^P((0,0),(2i,0))\le C(1+i)^{p_P}n^{-p_P-1}.
\end{equation}
Exponential clock errors are absorbed in this bound uniformly in $i$. Summation in \eqref{cps:eq:Pcount} is legitimate because $\sum_{i\ge1}3^{-i}(1+i)^{p_P}<\infty$. This proves the upper half of the $P$ estimate.

For $S$, apply \eqref{cps:eq:threeupper} to the fixed endpoints $a=(0,2)$ and $b=(1,1)$. Regenerate at even parity from $a$ and at odd parity for the reverse chain from $b$. Equations~\eqref{cps:eq:cyclemoments} and \eqref{cps:eq:survival} apply in both cases, so
\begin{equation}\label{cps:eq:Supper}
 q_N^S((0,2),(1,1))\le CN^{-p_S-1}.
\end{equation}

\subsection{Lower bounds}
Use the same three-piece split and restrict both joining points $u,v$ to $B_n$ from Lemma~\ref{cps:lem:killed}. The entrance paths and Proposition~\ref{cps:prop:mass} give forward and reverse masses at least $cn^{-p/2}$ in $B_n$. To handle the floor in $m$, choose a fixed ball $C$ with closure inside $\sqrt3B$. Since $\sqrt{n/m}\to\sqrt3$, one has $\sqrt mC\subset\sqrt nB$ for every sufficiently large $n$. A fixed-ball use of Proposition~\ref{cps:prop:mass} at time $m$ supplies the required mass. The fixed entrance times are handled by one further nested ball, if needed.

The middle length $\ell$ lies in $[n/4,n/2]$ for all large $n$. Thus
\begin{equation}\label{cps:eq:threelower}
 q_n(a,b)\ge
 \sum_{u,v\in B_n}q_m(a,u)\frac cn q_m^*(b,v)
 \ge c n^{-p/2}n^{-1}n^{-p/2}=cn^{-p-1}
\end{equation}
for the endpoint pairs $(a,b)=((0,0),(4,0))$ in $P$ and $((0,2),(1,1))$ in $S$.

Keeping only $i=2$ in \eqref{cps:eq:Pcount} proves
$p_n\ge c\mu_P^n n^{-p_P-1}$. For $S$, \eqref{cps:eq:Scount} gives
$s'_n=\Theta(\mu_S^n n^{-p_S-1})$; then \eqref{cps:eq:Sconversion}, with nonnegative summands and fixed shifts of $n$, gives the same estimate for $s_n$. This completes the proof of Theorem~\ref{cps:thm:main}.

\section{Consequences}\label{cps:sec:consequences}
\subsection{Irrational exponents and non-D-finiteness}
\begin{lemma}\label{cps:lem:irrational}
Both $\alpha_P$ and $\alpha_S$ are irrational.
\end{lemma}
\begin{proof}
Let $r\in\{9/16,22/27\}$ and suppose $1+\pi/\arccos r$ is rational. Then $\arccos r$ is a rational multiple of $\pi$, so $\zeta=e^{i\arccos r}$ is a root of unity. Therefore $2r=\zeta+\zeta^{-1}$ is an algebraic integer. As a rational algebraic integer it must be an integer, contradicting $2r=9/8$ or $44/27$.
\end{proof}

\begin{lemma}[A two-sided coefficient criterion]\label{cps:lem:nonD}
Suppose $a_n$ are nonnegative integers, $\mu>1$, $\alpha$ is irrational, and $a_n=\Theta(\mu^n n^{-\alpha})$. Then $A(z)=\sum_{n\ge0}a_nz^n$ is not D-finite.
\end{lemma}
\begin{proof}
Put $\rho=\mu^{-1}$ and choose a nonnegative integer $k>\alpha-1$. The coefficient bounds and nonnegativity imply, as $t\downarrow0$,
\begin{equation}\label{cps:eq:derivativegrowth}
 A^{(k)}(\rho(1-t))
 =\Theta\!\left(\sum_{n\ge k}n^{k-\alpha}(1-t)^n\right)
 =\Theta(t^{\alpha-k-1}).
\end{equation}
Factors $\rho^{-k}(1-t)^{-k}$ are bounded and positive, and the falling factorial $n(n-1)\cdots(n-k+1)$ is comparable to $n^k$ for large $n$. With $\beta=k-\alpha>-1$, the last estimate follows from
\begin{equation}
 \sum_{n\ge1}n^\beta e^{-sn}\sim\Gamma(\beta+1)s^{-\beta-1},\qquad s\downarrow0,
\end{equation}
by a Riemann-sum comparison and $s=-\log(1-t)\sim t$. Initial terms do not affect a divergent derivative.

If $A$ were D-finite, its integer coefficients and exponential upper bound would make it a $G$-function: its coefficient conjugates have exponential growth and its coefficient denominators are all one. The Andr\'e--Chudnovsky--Katz regular-singularity and rational-exponent theorem implies that its minimal differential equation has only regular singularities with rational local exponents. The finite rational-power and logarithmic form used below is stated explicitly in \cite[Theorem 6, Section 4.1]{FR}; see also \cite[Section 2.2, proof of Theorem 2]{BRS}.

Consequently, a local solution near $z=\rho$, and every fixed derivative of it, has a finite Frobenius expansion in rational powers and nonnegative integer powers of a logarithm, with analytic coefficient functions. Along the positive real approach, a nonzero divergent such function has a rational logarithmic growth exponent: combine terms with equal powers, choose the smallest surviving rational power, and then the highest nonzero logarithmic power at that order. In contrast, \eqref{cps:eq:derivativegrowth} gives
\begin{equation}
 \lim_{t\downarrow0}\frac{\log A^{(k)}(\rho(1-t))}{\log(1/t)}=k+1-\alpha,
\end{equation}
which is irrational. This contradiction proves the result.
\end{proof}
\begin{remark}
The coefficient criterion stated in \cite[Theorem 2]{BRS} assumes an asymptotic equivalent. It is not directly applicable to a $\Theta$ estimate. The positive-derivative argument in \eqref{cps:eq:derivativegrowth} is the additional step used here. For a rational-coefficient series, D-finiteness over $\mathbb C$ is equivalent to D-finiteness over $\mathbb Q$: a relation of fixed order and polynomial degree imposes rational linear equations on finitely many operator coefficients, and a rational row-space basis gives a nonzero rational solution whenever a nonzero complex solution exists.
\end{remark}

\begin{writenotelog}
The two-sided bounds are more than the argument needs: with a derivative order $k>\alpha-1$
and an $\varepsilon$ of room in the exponent, the same proof works for
$a_n=\mu^nn^{-\alpha+o(1)}$. This form is run in Part~\ref{cps:part:logcone}
(Section~\ref{cps:log:sec:nonD}) and stated in general in Remark~\ref{cps:log:rem:criterion}, and
also in Remark 10.1 (\texttt{osr:rem:criterion}) of the report
\repo{a348351-one-sided-rectangulations} in this collection (bundle Report 111, batch 103).
There the irrationality certificate differs from Lemma~\ref{cps:lem:irrational}:
$2\cos\theta=(-29+7\sqrt{17})/2$ is an algebraic integer, and the certificate is its conjugate,
which lies below $-2$.
\end{writenotelog}

\begin{corollary}\label{cps:cor:nonD}
The ordinary generating functions of A377922, A377920, and A377921 are not D-finite.
\end{corollary}
\begin{proof}
The first two cases follow from Theorem~\ref{cps:thm:main} and Lemmas~\ref{cps:lem:irrational}--\ref{cps:lem:nonD}. For A377921, let $\widetilde s_n$ be the rigid-orthogonal-surface sequence, and put
$S(z)=\sum_{n\ge2}s_nz^n$, $T(z)=\sum_{n\ge2}\widetilde s_nz^n$. Remark 20 of \cite{FNS} gives
\begin{equation}\label{cps:eq:rigid}
 S(z)=(1+z)^3T(z)+3z^2+z^3.
\end{equation}
D-finite series are closed under multiplication by rational functions and addition of polynomials. Equation~\eqref{cps:eq:rigid} therefore makes $S$ D-finite if and only if $T$ is D-finite.
\end{proof}
No coefficientwise $\Theta$ estimate for A377921 follows from this argument; the statement about that sequence concerns its generating function only.

\begin{writenote}
A coefficientwise statement for A377921 is now proved in Part~\ref{cps:part:amplitudes}:
equation~\eqref{cps:amp:eq:rigidmain} gives
$\widetilde s_n\sim(16/19)^3\kappa_S\mu_S^nn^{-\alpha_S}$, by a dominated signed convolution that
uses the full equivalent for $s_n$. As Part~\ref{cps:part:amplitudes} says, the $\Theta$ bounds
of Theorem~\ref{cps:thm:main} alone do not give it.
\end{writenote}

\begin{writenotelog}
A coefficientwise $\Theta$ estimate for A377921 does follow from Theorem~\ref{cps:thm:main}
together with the sleeve injection of Part~\ref{cps:part:logcone}, which gives
$s_{n-6}\le\widetilde s_n\le s_n$ for $n\ge8$ (Proposition~\ref{cps:log:prop:sleeve},
Remark~\ref{cps:log:rem:theta}). Part~\ref{cps:part:logcone} also proves non-D-finiteness from
the logarithmic law alone (Section~\ref{cps:log:sec:nonD}, Remark~\ref{cps:log:rem:criterion}), a
weaker hypothesis than that of Lemma~\ref{cps:lem:nonD}.
\end{writenotelog}

\subsection{First-threshold inverse asymptotics}
\begin{lemma}\label{cps:lem:inverse}
Suppose $a_n=\Theta(\mu^n n^{-\alpha})$ with $\mu>1$ and $\alpha\in\RR$, and define
$N(Y)=\min\{n\ge0:a_n\ge Y\}$. Then
\begin{equation}\label{cps:eq:inversegeneral}
 N(Y)=\frac{\log Y+\alpha\log\log Y}{\log\mu}+O(1)
 \qquad(Y\to\infty).
\end{equation}
The sequence $a_n$ need not be monotone.
\end{lemma}
\begin{proof}
Put $L=\log Y$ and $v(L)=(L+\alpha\log L)/\log\mu$. Uniformly for $h$ in a fixed bounded interval,
\begin{equation}
 \log\!\left(\mu^{v(L)+h}(v(L)+h)^{-\alpha}\right)
 =L+h\log\mu+\alpha\log\log\mu+o(1).
\end{equation}
Choose a constant $H$ sufficiently large that the upper coefficient comparison is less than $Y$ at $v(L)-H$, while the lower comparison is at least $Y$ at $v(L)+H$, allowing another fixed unit for integer rounding. The comparison function $\mu^n n^{-\alpha}$ is eventually increasing. Hence all sufficiently large integers below the lower bracket have $a_n<Y$, and the finitely many earlier values also do so for large $Y$. The upper bracket supplies at least one index above threshold. This proves \eqref{cps:eq:inversegeneral}.
\end{proof}
In particular,
\begin{align}
 N_P(Y)&=\frac{\log Y+\alpha_P\log\log Y}{\log(9/2)}+O(1),\label{cps:eq:Pinverse}\\
 N_S(Y)&=\frac{\log Y+\alpha_S\log\log Y}{\log(16/3)}+O(1).\label{cps:eq:Sinverse}
\end{align}
These estimates do not specify an exact rounding rule or a next additive constant.

\begin{writenote}
Part~\ref{cps:part:amplitudes}, Theorem~\ref{cps:amp:thm:inverse}, sharpens
\eqref{cps:eq:Pinverse}--\eqref{cps:eq:Sinverse} under the full equivalents: Lambert-$W_{-1}$
centering with a next additive constant $-(\alpha\log\log\mu+\log\gamma)/\log\mu$ and a
vanishing-width two-ceiling bracket. It still gives no exact rounding rule. Both inversions are
instances of the repository's transseries volume
(\repo{Analysis/Transseries/docs/series-and-transseries/Transseries_And_Inversion/}): the
model equation $\lambda r-\alpha\log r=\log(Y/\gamma)$ is \texttt{p0:thm:lambert-core} with
$a=\log\mu$, $b=-\alpha$ (branch $W_{-1}$), its expansion is the first step of
\texttt{p0:thm:lambert-centered}, and the integer bracket is the separation situation of
\texttt{p0:thm:staircase}. No novelty is claimed for the inversion mechanics.
\end{writenote}

\section{Verification and unresolved questions}\label{cps:sec:verification}
The reproducibility package contains editable \LaTeX\ source, the compiled PDF, exact symbolic and enumeration programs, their recorded outputs, source identifiers, and mathematical audit receipts. Appendix~\ref{cps:app:checks} specifies the finite checks. These checks corroborate normalizations and identities; they are not a formal verification of the probabilistic proof or an external peer review.

\begin{writenote}
In this repository the delivered PDF and the checksum ledger of this package are not
shipped; the report's \texttt{article.pdf} is a build of this merged text. The programs,
outputs, source excerpts and receipts are shipped byte-identical under the prefix
\texttt{45-cone-} (programs in \texttt{code/}, outputs and data in \texttt{data/}, the two audit
notes at the report root). The report README lists them and gives a rerun recipe.
\end{writenote}

Theorem~\ref{cps:thm:main} does not prove either of the limits
\begin{equation}\label{cps:eq:amplitudes}
 \lim_{n\to\infty}p_n\mu_P^{-n}n^{\alpha_P},\qquad
 \lim_{n\to\infty}s_n\mu_S^{-n}n^{\alpha_S}.
\end{equation}
It proves that both normalized sequences are eventually bounded above and bounded away from zero. Convergence to positive constants, the full asymptotic equivalents of \cite[Conjecture 25]{FNS}, remains to be established. Computing those constants and correction terms would require sharper information than the inequalities used in the clock transfer and gluing arguments. Other open directions include a direct coefficientwise estimate for A377921 and identifying other parity-dependent tandem models for which a fixed cycle buffer exists.

\begin{writenote}
Status after the merge. Convergence to positive constants and the full equivalents of
\cite[Conjecture 25]{FNS} are proved in Part~\ref{cps:part:amplitudes}
(Theorem~\ref{cps:amp:thm:main}), and so is a direct coefficientwise equivalent for A377921
(equation~\eqref{cps:amp:eq:rigidmain}). Still open: the numerical values of $\kappa_P$ and
$\kappa_S$, a rate of convergence, correction terms, an exact rounding rule for the threshold
inverses, and other parity-dependent tandem models with a fixed cycle buffer.
\end{writenote}

\begin{writenotelog}
For the last item, Part~\ref{cps:part:logcone}'s Theorem~\ref{cps:log:thm:cone} gives the
logarithmic-scale exponent for every walk with a finite internal state that satisfies
(H1)--(H5), with or without a fixed cycle buffer; bundle Report 111
(\repo{a348351-one-sided-rectangulations}, same collection directory) is one such chain
(Section~\ref{cps:log:sec:siblings}). Models with a buffer still matter for amplitudes and
$\Theta$ bounds, which the general theorem does not give.
\end{writenotelog}

This research note presents a mathematical proof with reproducible exact checks. It has not been externally peer reviewed or machine-formalized, and makes no publication-priority claim. The source credits distinguish the established bijections and iid cone input from the argument presented here.


\clearpage
\part{Limiting amplitudes (source 46)}\label{cps:part:amplitudes}
\begin{center}\large Limiting amplitudes for corner polyhedra and Schnyder labelings\\[2pt]
\normalsize Research addendum, 2 October 2026\end{center}
\begin{sourceabstract}{source 46}
We strengthen the two-sided estimates for the parity-dependent quadrant walks counting corner polyhedra and 3-connected Schnyder labelings to full asymptotic equivalents with finite positive amplitudes. The key step is a harmonic comparison for a cycle kernel trapped between an ordinary iid cone kernel and its fixed inward translate. It yields exact survival constants and a common endpoint limit law for valid regeneration cycles. A two-sided random-clock argument transfers these limits to every exact counted time. A killed local limit, proved from the free local theorem and the exit decomposition, then gives the fixed-endpoint amplitudes by a fixed-thirds semigroup limit. Geometrically weighted endpoint summation and a signed rational convolution give the enumerative equivalents, including the rigid-surface sequence. The amplitude constants have positive harmonic and heat-kernel expressions; no numerical digits are certified. We also give Lambert-$W$ inverse centering and a vanishing-width ceiling bracket, without asserting an unconditional rounding formula.
\end{sourceabstract}

\section{Results and the accompanying foundation}
This addendum is intended to be read with the included research note \emph{Polynomial exponents for corner polyhedra and Schnyder labelings}, cited as the \emph{Foundation} (in this report: Part~\ref{cps:part:exponents}). The Foundation and its exact-check package are included unchanged. They establish the walk models, covariance calculations, fixed-buffer geometry, free local theorem and two-sided estimates used below. This addendum supplies the limiting-amplitude argument, rather than assuming a cone local theorem for the original parity-dependent chain. Statements in the Foundation that amplitude convergence was unresolved describe the scope of that earlier note; the present addendum establishes those limits.

\begin{writenote}
In this report the Foundation is Part~\ref{cps:part:exponents} (Sections 1--10 and
Appendix~\ref{cps:app:checks}), printed with source 45's numbering, so that ``Foundation
Lemma~2.1'' below is Lemma~\ref{cps:lem:Sboundary} here; the Foundation references are
hyperlinks. The unchanged 30-file copy of the Foundation package that source 46 embedded is not
shipped: the same files are shipped once, from source 45, with the prefix \texttt{45-cone-}
(the copy's PDF and checksum ledger only in the arrival commit). The source's Section $n$ is
Section $n+10$ here.
\end{writenote}

Let $p_n$ count polyhedral orientations with $n$ inner vertices, equivalently corner-polyhedron types with $n+3$ flats. Let $s_n$ count 3-connected Schnyder labelings with $n$ inner faces, and let $\widetilde s_n$ be the rigid-orthogonal-surface sequence of Fusy--Narmanli--Schaeffer \cite{FNS}. These are A377922, A377920 and A377921, respectively. Set
\begin{equation}\label{cps:amp:eq:parameters}
 \mu_P=\frac92,\quad \mu_S=\frac{16}3,\qquad
 \alpha_P=1+\frac{\pi}{\arccos(9/16)},\quad
 \alpha_S=1+\frac{\pi}{\arccos(22/27)}.
\end{equation}

\begin{theorem}[Full enumerative equivalents]\label{cps:amp:thm:main}
There are finite strictly positive constants $\kappa_P$ and $\kappa_S$ such that
\begin{align}
 p_n&\sim\kappa_P\mu_P^n n^{-\alpha_P},\label{cps:amp:eq:Pmain}\\
 s_n&\sim\kappa_S\mu_S^n n^{-\alpha_S},\label{cps:amp:eq:Smain}\\
 \widetilde s_n&\sim\left(\frac{16}{19}\right)^3\kappa_S\mu_S^n n^{-\alpha_S}.
 \label{cps:amp:eq:rigidmain}
\end{align}
The limits hold on the full sequence of integers. Positive expressions for the amplitudes are given in Section~\ref{cps:amp:sec:enumeration}.
\end{theorem}
The first two conclusions are the positive-amplitude equivalents stated in \cite[Conjecture 25]{FNS}. The third follows here by a dominated signed-convolution argument, which uses the full equivalent and is not a consequence of the earlier $\Theta$ bounds alone.

\begin{center}
\begin{tabular}{@{}lccc@{}}
\toprule
Sequence & Rate & Power & Amplitude\\
\midrule
A377922 & $9/2$ & $\alpha_P$ & $\kappa_P>0$\\
A377920 & $16/3$ & $\alpha_S$ & $\kappa_S>0$\\
A377921 & $16/3$ & $\alpha_S$ & $(16/19)^3\kappa_S$\\
\bottomrule
\end{tabular}
\end{center}
No numerical evaluation of these constants, convergence rate, or further coefficient correction is claimed. This is a mathematical research addendum with exact reproducibility checks, not an externally peer-reviewed publication or a machine-formalized proof. No publication-priority claim is made.

\section{Model notation and established inputs}\label{cps:amp:sec:inputs}

\begin{writenote}
This section restates inputs proved or quoted in Part~\ref{cps:part:exponents}, in
source 46's notation; it is kept because the rest of Part~\ref{cps:part:amplitudes} cites its
displays. $D_P$ and $H_S$ here are Part~\ref{cps:part:exponents}'s $D$ and $H$
(\eqref{cps:eq:PK}, \eqref{cps:eq:SM}); the letter $D$ alone is reserved in this Part for the
valid-cycle kernel of Section~\ref{cps:amp:sec:kernels}. Table~\ref{cps:tab:notation} lists every
clash between the two Parts.
\end{writenote}
Let
\begin{equation}
 L=\{(a,b)\in\ZZ^2:a\equiv b\pmod2\},\qquad
 C=(0,\infty)^2.
\end{equation}
For $P$, one unit of counted time is one physical tandem step and the killed domain is $Q_P=L\cap[0,\infty)^2$. For $S$, one unit is an aggregate consisting of a southeast step $(1,-1)$ followed by zero or more weighted face steps, and the killed domain is $Q_S=\{(a,b)\in L:a\ge0,b\ge1\}$. In the latter model, restricting aggregate endpoints to $Q_S$ controls all physical intermediate steps, as proved in Foundation Lemma~\ref{cps:lem:Sboundary}.

Write $q_n(a,b)$ for the counted-time transition probability killed outside the model domain. The model is specified by a $2\times2$ probability-generating matrix; its two internal states are the parity of the ordinate. In the notation
\begin{equation}
 D_P=\left(1-\frac1{3x^2}\right)\left(1-\frac{y^2}3\right),\qquad
 H_S=1-x^{-2}-y^2-2x^{-2}y^2,
\end{equation}
the matrices are
\begin{align}
 K_P(x,y)&=\begin{pmatrix}
 2y^2/(27D_P)&2x/(3y)+2y/(27xD_P)\\
 2x/(3y)+2y/(27xD_P)&2/(27x^2D_P)
 \end{pmatrix},\label{cps:amp:eq:PK}\\
 K_S(x,y)&=\frac3{16}M_S(2x,y/2),\qquad
 M_S(x,y)=\begin{pmatrix}
 x^{-2}/H_S&x/y+x^{-1}y/H_S\\
 x/y+x^{-1}y/H_S&y^2/H_S
 \end{pmatrix}.\label{cps:amp:eq:SK}
\end{align}
The covariance matrices per unit counted time are
\begin{equation}\label{cps:amp:eq:Sigma}
 \Sigma_P=\begin{pmatrix}16/5&-9/5\\-9/5&16/5\end{pmatrix},\qquad
 \Sigma_S=\begin{pmatrix}36/7&-88/21\\-88/21&36/7\end{pmatrix}.
\end{equation}
Counting measure on $L$ is invariant. The reversed chain has
$K^*_{ij}(z)=K_{ji}(-z)$ and satisfies
\begin{equation}\label{cps:amp:eq:reverse}
 q_n(a,b)=q_n^*(b,a).
\end{equation}
The symmetric internal matrices have stationary distribution $(1/2,1/2)$.

For either initial parity, let $T_k$ be successive first-return times to that parity. The cycle pairs $(Y_k,\ell_k)$ of displacement and duration are iid, with
\begin{equation}\label{cps:amp:eq:cycles}
 \EE Y=0,\quad\Cov(Y)=2\Sigma,\quad\EE\ell=2,
 \quad\EE e^{\varepsilon(\norm Y+\ell)}<\infty
\end{equation}
for some $\varepsilon>0$. The reversed cycles have displacement law $-Y$ and the same duration law. Dependence between displacement and duration is permitted and retained.

The following previously established estimates will be used with their stated scopes:
\begin{enumerate}
\item \emph{Cycle buffer.} Every intermediate counted position in a first-return cycle with endpoints $a,b$ has coordinate $r$ at least $\min(a_r,b_r)-1$. This holds forward and backward and at either parity (Foundation Lemma~\ref{cps:lem:buffer}).
\item \emph{Bounded corrector.} $M_n=X_n+h_{J_n}$ is a martingale; its increments have uniformly bounded conditional exponential moments. The reversed chain has corrector $-h$ and the same averaged covariance (Foundation Section~\ref{cps:sec:moments}).
\item \emph{Survival upper bound.} For either chain,
\begin{equation}\label{cps:amp:eq:survivalupper}
 \PP_a(\tau>n)\le C(1+\norm a)^p n^{-p/2}+Ce^{-cn},
\end{equation}
where $p=\alpha-1>1$ (Foundation Proposition~\ref{cps:prop:survival}).
\item \emph{Free local theorem.} On the unrestricted chain,
\begin{equation}\label{cps:amp:eq:freeupper}
 \sup_{a,b\in L}\PP_a(X_n=b)\le C/n,
\end{equation}
and, uniformly on compact diffusive displacement sets,
\begin{equation}\label{cps:amp:eq:freelocal}
 \PP_a(X_n=b)=\frac{2}{2\pi n\sqrt{\det\Sigma}}
 \exp\!\left(-\frac{(b-a)^T\Sigma^{-1}(b-a)}{2n}\right)+o(n^{-1}).
\end{equation}
All parity pairs and all sufficiently large integer times are included. The factor two is the covolume of $L$ (Foundation Proposition~\ref{cps:prop:local}).
\item \emph{Uniform endpoint estimate for $P$.}
\begin{equation}\label{cps:amp:eq:endpointupper}
 q_n^P((0,0),(2i,0))\le C(1+i)^{p_P}n^{-p_P-1},\qquad i\ge1
\end{equation}
(Foundation equation~\eqref{cps:eq:Pupper}).
\end{enumerate}
The normalization identities are
\begin{align}
 p_n&=\mu_P^n\sum_{i\ge1}3^{-i}q_n^P((0,0),(2i,0)),\label{cps:amp:eq:Pcount}\\
 s'_n&=\frac{\mu_S^{n+1}}4q_{n+1}^S((0,2),(1,1)),
 &s_n&=s'_n+2s'_{n-1}+s'_{n-2}\quad(n\ge4),\label{cps:amp:eq:Scount}
\end{align}
where $s'_n$ restricts the two outer $G$ vertices to be non-isolated. The full rigid-series relation is
\begin{equation}\label{cps:amp:eq:series}
 S(z)=(1+z)^3\widetilde S(z)+3z^2+z^3,
 \quad S(z)=\sum_{n\ge2}s_nz^n,\quad
 \widetilde S(z)=\sum_{n\ge2}\widetilde s_nz^n.
\end{equation}
These are the counting bijections and relations of \cite[Propositions 13--14, Remarks 9 and 20]{FNS}, in the normalized form verified in the Foundation.

\section{Ordinary and valid-cycle cone kernels}\label{cps:amp:sec:kernels}
Fix a model, an initial parity and a time direction. Let
\begin{equation}\label{cps:amp:eq:whitening}
 T=(2\Sigma)^{-1/2},\qquad K=TC.
\end{equation}
Whitening the cycle covariance makes its increment covariance the identity. The fixed translation of the original domain is $c_P=(1,1)$ or $c_S=(1,0)$, before applying $T$. A starting point $a$ thus becomes $x=T(a+c_M)$, for model $M$. The open wedge $K$ has angle
\begin{equation}
 \theta=\arccos\!\left(-\frac{\Sigma_{12}}{\sqrt{\Sigma_{11}\Sigma_{22}}}\right),
 \qquad p=\pi/\theta>1.
\end{equation}
In polar coordinates adapted to its boundary rays, normalize the positive Brownian harmonic function by
\begin{equation}\label{cps:amp:eq:u}
 u(r,\varphi)=r^p\sin(p\varphi),\qquad0<\varphi<\theta.
\end{equation}
All kernels below live on the corresponding translated and whitened lattice coset. Fixed translations need not preserve that coset to define the translated starting state of the same iid walk.

Let $P_0$ be the iid cycle-displacement kernel killed outside $K$. Choose
\begin{equation}
 b_P=T(2,2),\qquad b_S=T(2,1),\qquad K_b=b+K.
\end{equation}
On either parity lattice, the inner condition $T(a+c_M)\in K_b$ is exactly $a_1,a_2\ge2$. Let $P_b$ be the iid displacement kernel killed outside $K_b$, extended by zero for starting points outside $K_b$. Let $D$ be the true substochastic kernel of one complete cycle, rejecting a cycle if any of its counted steps leaves the model domain. By the cycle buffer,
\begin{equation}\label{cps:amp:eq:kernelorder}
 0\le P_b\le D\le P_0
\end{equation}
pointwise as kernels. In fact, if both endpoints lie in $K_b$, every marked cycle between them is safe, so $D$ agrees there with the unrestricted cycle kernel. The inequality concerns entire kernels, not only their row sums.

\subsection{The iid harmonic input}
The hypotheses of Denisov--Wachtel \cite[Theorems 2--3]{DWrevisited} hold for these iid increments: centering, identity covariance, exponential moments, the wedge geometry and $p\ge1$ were checked in the Foundation. Write $Z_m$ for the whitened iid walk, $\tau_0$ for its exit from $K$, and
\begin{equation}\label{cps:amp:eq:V}
 V(x)=\lim_{m\to\infty}\EE_x[u(Z_m);\tau_0>m].
\end{equation}
The theorem provides $P_0V=V$, uniform $V(z)/u(z)\to1$ as $\operatorname{dist}(z,\partial K)\to\infty$, a polynomial survival upper bound, and
\begin{align}
 m^{p/2}P_0^m\one(x)&\longrightarrow\beta V(x),\label{cps:amp:eq:iidsurvival}\\
 m^{p/2}\EE_x[f(Z_m/\sqrt m);\tau_0>m]
 &\longrightarrow\beta V(x)\nu(f),\label{cps:amp:eq:iidYaglom}
\end{align}
for bounded continuous nonnegative $f$. Here $\beta>0$ uses the normalization \eqref{cps:amp:eq:u}, and $\nu$ is the normalized Brownian cone-meander endpoint law in $K$. The corresponding translated kernel has harmonic function
\begin{equation}\label{cps:amp:eq:Vb}
 V_b(x)=V(x-b)\one_{\{x\in K_b\}},\qquad P_bV_b=V_b,
\end{equation}
and the same constants $\beta,\nu$, since the shift is fixed. Coupling the same iid increments started at $x$ and $x-b$ makes survival in the translated inner wedge imply survival in the outer wedge. Taking the limits in \eqref{cps:amp:eq:iidsurvival} gives
\begin{equation}\label{cps:amp:eq:Vorder}
 0\le V_b(x)\le V(x).
\end{equation}

\section{The boundary-shift comparison}\label{cps:amp:sec:shift}
\begin{lemma}\label{cps:amp:lem:shift}
For every fixed outer starting point $x$,
\begin{equation}\label{cps:amp:eq:shiftlimit}
 P_0^m(V-V_b)(x)\longrightarrow0.
\end{equation}
\end{lemma}
\begin{proof}
Since $P_0^mV(x)=V(x)$, it suffices to prove $P_0^mV_b(x)\to V(x)$. Put $d(z)=\operatorname{dist}(z,\partial K)$. We first show that for every fixed $R$,
\begin{equation}\label{cps:amp:eq:strip}
 \EE_x[u(Z_m);\tau_0>m,\ d(Z_m)\le R]\longrightarrow0.
\end{equation}
On $\norm{Z_m}\le\sqrt m\log m$, the wedge estimate
$u(z)\le C\norm z^{p-1}d(z)$ and the survival upper bound give
\begin{equation}
 \EE_x[u(Z_m);\tau_0>m,d(Z_m)\le R,\norm{Z_m}\le\sqrt m\log m]
 \le C_xR m^{-1/2}(\log m)^{p-1}=o(1).
\end{equation}
For the complementary region no killing estimate is needed. Centered iid exponential moments imply, for $t$ sufficiently larger than the fixed starting norm,
\begin{equation}\label{cps:amp:eq:tail}
 \PP_x(\norm{Z_m}>t)\le C_x\exp[-c\min(t^2/m,t)].
\end{equation}
For $L_m=\sqrt m\log m$, integration of this tail bound yields
\begin{align}
 \EE_x[\norm{Z_m}^p;\norm{Z_m}>L_m]
 &\le L_m^p\PP_x(\norm{Z_m}>L_m)
   +\int_{L_m}^{\infty}pt^{p-1}\PP_x(\norm{Z_m}>t)\,dt\notag\\
 &=O\!\left(m^{p/2}e^{-c'(\log m)^2}\right)+O(e^{-c'm})=o(1),\label{cps:amp:eq:tailmoment}
\end{align}
after absorbing powers of $\log m$ into a smaller positive exponential constant. Because $u(z)\le C\norm z^p$, this proves \eqref{cps:amp:eq:strip}.

Given $\varepsilon>0$, choose $R$ large. If $d(z)>R$, then $z-b\in K$ and the uniform asymptotic for $V$ gives $V(z-b)/u(z-b)$ close to one. Also
\begin{equation}\label{cps:amp:eq:ushift}
 \frac{u(z-b)}{u(z)}\longrightarrow1
 \quad\text{uniformly as }d(z)\longrightarrow\infty.
\end{equation}
To see the uniformity even when $\norm z$ grows independently, use the explicit wedge function to obtain $\norm{\nabla\log u(w)}\le C/d(w)$. Along $w=z-tb$, $0\le t\le1$, the boundary distance is at least $d(z)-\norm b$; integrate the logarithmic derivative along this bounded segment. Consequently, for sufficiently large $R$,
\begin{equation}
 V_b(z)\ge(1-\varepsilon)u(z)\qquad(d(z)>R).
\end{equation}
Combining this inequality with \eqref{cps:amp:eq:V} and \eqref{cps:amp:eq:strip} gives
$\liminf_mP_0^mV_b(x)\ge(1-\varepsilon)V(x)$. By \eqref{cps:amp:eq:Vorder}, the corresponding upper bound is $P_0^mV_b(x)\le V(x)$. Let $\varepsilon\downarrow0$.
\end{proof}
The argument needs only the ordinary iid harmonic asymptotic and survival bound. It does not assume a conditioned invariance principle for the valid-cycle kernel.

\section{Harmonic and endpoint limits for valid cycles}\label{cps:amp:sec:validcycles}
\begin{theorem}\label{cps:amp:thm:cyclelimits}
The limit
\begin{equation}\label{cps:amp:eq:H}
 \mathcal H(x)=\lim_{m\to\infty}D^mV(x)
\end{equation}
exists, is finite, is $D$-harmonic, and is positive at every starting point that can reach a point with $V_b>0$ by a valid finite cycle path. Moreover,
\begin{align}
 n^{p/2}D^n\one(x)&\longrightarrow\beta\mathcal H(x),\label{cps:amp:eq:Dsurvival}\\
 n^{p/2}\EE_x[f(Z_n/\sqrt n);\text{the first }n\text{ cycles are valid}]
 &\longrightarrow\beta\mathcal H(x)\nu(f)\label{cps:amp:eq:DYaglom}
\end{align}
for every bounded continuous nonnegative $f$.
\end{theorem}
\begin{proof}
The inequality $DV\le P_0V=V$ makes $D^mV$ nonincreasing and nonnegative, proving existence and finiteness. Since $DV_b\ge P_bV_b=V_b$ and $V\ge V_b$, one has $D^mV\ge V_b$. Dominated convergence against $D(x,dz)$, using $DV(x)\le V(x)<\infty$, gives $D\mathcal H=\mathcal H$. Thus $\mathcal H\ge V_b$, and harmonicity propagates strict positivity backwards along every valid finite path to a point where $V_b>0$.

Fix $m$ and write $n=m+r$. Kernel order gives
\begin{equation}\label{cps:amp:eq:sandwich}
 D^mP_b^r\one(x)\le D^n\one(x)\le D^mP_0^r\one(x).
\end{equation}
Multiply by $n^{p/2}$ and let $n\to\infty$ with $m$ fixed. The ordinary survival theorem and its uniform polynomial upper bound permit dominated convergence against $D^m(x,dz)$. Indeed, $D^m$ is bounded by the fixed-time iid displacement kernel, so this measure has all polynomial moments. It follows that
\begin{equation}\label{cps:amp:eq:limitbracket}
 \beta D^mV_b(x)
 \le\liminf_n n^{p/2}D^n\one(x)
 \le\limsup_n n^{p/2}D^n\one(x)
 \le\beta D^mV(x).
\end{equation}
The width is at most
\begin{equation}
 \beta D^m(V-V_b)(x)\le\beta P_0^m(V-V_b)(x)\longrightarrow0
\end{equation}
by Lemma~\ref{cps:amp:lem:shift}. The upper endpoint tends to $\beta\mathcal H(x)$, proving \eqref{cps:amp:eq:Dsurvival}.

For \eqref{cps:amp:eq:DYaglom}, insert the nonnegative terminal test function $f(z/\sqrt n)$ in \eqref{cps:amp:eq:sandwich}. For each fixed intermediate starting point, the ordinary endpoint theorem gives the same factor $\nu(f)$ in the inner and outer kernels: the shift is fixed and $r/n\to1$. Domination follows from $\norm f_\infty$ times the same survival bound. Thus the analogue of \eqref{cps:amp:eq:limitbracket} has both endpoints multiplied by $\nu(f)$, and the same vanishing-width argument proves the result. Taking $f=1$ gives convergence of total mass; together the statements give the normalized weak conditional endpoint limit.
\end{proof}

\section{Exact counted-time constants}\label{cps:amp:sec:clock}
Let $R_n$ be survival of the original counted-step chain through time $n$. For an initial parity $i$, use its complete-cycle kernel $D_i$ and the corresponding harmonic function $\mathcal H_i$; forward and reverse kernels may have different harmonic functions. Define
\begin{equation}\label{cps:amp:eq:A}
 A(a)=2^{p/2}\beta\mathcal H_{i(a)}(T(a+c_M)),\qquad
 A^*(b)=2^{p/2}\beta^*\mathcal H_{i(b)}^*(T(b+c_M)).
\end{equation}
Here $i(a)$ and $i(b)$ denote the respective parities of the two endpoints. The starred harmonic normalization uses the same wedge function \eqref{cps:amp:eq:u}; writing $\beta^*$ separately allows the ordinary reverse iid theorem to be applied without identifying its constant in advance. Its endpoint law is the same $\nu$, since the covariance and wedge agree.

\begin{theorem}\label{cps:amp:thm:exacttime}
At every relevant forward and reverse starting point,
\begin{equation}\label{cps:amp:eq:exacttime}
 \PP_a(R_n)\sim A(a)n^{-p/2},\qquad
 \mathcal L_a(X_n/\sqrt n\mid R_n)\Longrightarrow\nu_\Sigma,
\end{equation}
with the analogous reversed statement and amplitude $A^*(b)$. The probability law $\nu_\Sigma$ is the endpoint law at time one of the Brownian meander in $C$ with covariance $\Sigma$.
\end{theorem}
\begin{proof}
Put
\begin{equation}
 k_- =\lfloor n/2-n^{3/4}\rfloor,\qquad
 k_+ =\lceil n/2+n^{3/4}\rceil.
\end{equation}
Chernoff bounds for the iid durations, with mean two, show that outside an event of probability $O(e^{-c\sqrt n})$,
\begin{equation}\label{cps:amp:eq:clockwindow}
 n-4n^{3/4}\le T_{k_-}\le n\le T_{k_+}\le n+4n^{3/4}.
\end{equation}
The corrected-martingale exponential maximal inequality from the Foundation, applied on the deterministic interval beginning at $\lfloor n-4n^{3/4}\rfloor$ and of length $O(n^{3/4})$, gives total positional oscillation at most $\delta\sqrt n$ outside an event of probability $O(e^{-c_\delta n^{1/4}})$. The bounded corrector contributes only $O(1)$. Both estimates are unconditional; no independence between duration and displacement and no conditioning on valid cycles is used.

For a continuity set $B$ of the limiting endpoint law, let $B^{-\delta}$ and $B^{+\delta}$ be its inner and outer $\delta$ neighborhoods. The two complete-cycle events give
\begin{align}\label{cps:amp:eq:eventbracket}
 &\PP\{D\text{-survival through }k_+,\ X_{T_{k_+}}/\sqrt n\in B^{-\delta}\}-\varepsilon_n\notag\\
 &\hspace{9mm}\le\PP_a(R_n,X_n/\sqrt n\in B)\notag\\
 &\hspace{9mm}\le\PP\{D\text{-survival through }k_-,\ X_{T_{k_-}}/\sqrt n\in B^{+\delta}\}+\varepsilon_n,
\end{align}
where $\varepsilon_n$ is exponentially small compared with $n^{-p/2}$. First choose $\delta$ so the neighborhood boundaries are continuity sets; then let $\delta\downarrow0$. Since $k_\pm/n\to1/2$, Theorem~\ref{cps:amp:thm:cyclelimits} yields a common amplitude factor $2^{p/2}\beta\mathcal H_i$. The spatial transformation is
\begin{equation}
 \frac1{\sqrt2}T^{-1}=(\Sigma)^{1/2},
\end{equation}
so the pushed-forward endpoint law is precisely $\nu_\Sigma$. The fixed domain translation vanishes under scaling. Using $B=\RR^2$ gives the survival equivalent, and general continuity sets give the conditional weak limit.
\end{proof}

\subsection{Positivity at the enumeration endpoints}
The valid entrances to the even inner cone were exhibited in Foundation Section~\ref{cps:sec:entrances}. For completeness, a complete odd-return path for the reversed $S$ chain from its odd endpoint is
\begin{equation}\label{cps:amp:eq:oddentrance}
 (1,1)\longrightarrow(0,2)\longrightarrow(2,2)
 \longrightarrow(4,2)\longrightarrow(3,3).
\end{equation}
The first and last steps are reversed southeast steps. Each middle aggregate of displacement $(2,0)$ reverses a southeast step followed by the face step $(-3,1)$. Its internal physical points remain nonnegative. The final odd point lies in the inner domain; the middle internal states are even, so this is a complete first return to odd parity.

For reversed $P$ from $(2i,0)$, $i\ge2$, two reversed southeast steps reach $(2i-2,2)$. For $i=1$, the valid complete even-return path is
\begin{equation}
 (2,0)\longrightarrow(1,1)\longrightarrow(3,1)\longrightarrow(2,2),
\end{equation}
whose middle step is the reversed odd-to-odd face displacement $(2,0)$. From these points the cycle-support accessibility constructions reach arbitrarily deep points where $V_b>0$. Thus every endpoint in the $P$ sum has a positive reversed amplitude, and all the fixed endpoints in \eqref{cps:amp:eq:Scount} have positive amplitudes as well.

\section{Uniform interior killed local convergence}\label{cps:amp:sec:killed}
Let $g_t^\Sigma(x,y)$ be the unrestricted Brownian transition density with covariance $\Sigma$, and let $k_t^{C,\Sigma}(x,y)$ be its density killed on leaving the open quadrant $C$. The distinction between the physical quadrant $C$ and the whitened wedge $K=TC$ is maintained throughout this section.

\begin{lemma}\label{cps:amp:lem:killed}
Fix $t>0$ and integers $m_n$ with $m_n/n\to t$. For compact sets $E,F\subset C$,
\begin{equation}\label{cps:amp:eq:killedlocal}
 \sup_{\substack{a,b\in L:\ a/\sqrt n\in E,\ b/\sqrt n\in F}}
 \left|nq_{m_n}(a,b)-2k_t^{C,\Sigma}(a/\sqrt n,b/\sqrt n)\right|
 \longrightarrow0.
\end{equation}
\end{lemma}
\begin{proof}
We give the exit-decomposition proof. First, the unrestricted counted process satisfies a functional central limit theorem with covariance $\Sigma$. This follows from the bounded-corrector martingale: its predictable quadratic variation, divided by $n$, converges to $t\Sigma$ uniformly on compact time intervals by ergodicity of the finite internal chain, while the exponential moment implies the martingale Lindeberg condition. The bounded corrector disappears at scale $\sqrt n$. Equivalently, one may use the iid cycle functional limit theorem and the renewal law $T_k/k\to2$; the dependent clock still converges to a deterministic time change. Exponential tails bound both maximum cycle length and maximum individual jump over $O(n)$ steps by $O(\log n)$ in probability, leaving an unfinished-cycle error $O(\log^2n)=o(\sqrt n)$.

It is enough first to describe $m_n=n$. The free exit decomposition is
\begin{equation}\label{cps:amp:eq:Hunt}
 q_n(a,b)=\PP_a(X_n=b)
 -\EE_a[\PP_{X_\tau}(X_{n-\tau}=b);\tau\le n],
\end{equation}
where $\tau$ is the first exit from the discrete model domain. Fix $\varepsilon>0$ and restrict the expectation to $\tau\le(1-\varepsilon)n$. The unrestricted functional limit theorem gives convergence of the rescaled exit time and position on this truncated event. The continuous-mapping assertion uses familiar path properties of nondegenerate Brownian motion in a quadrant: it crosses a regular face immediately upon first hitting it, has zero probability of exit at the vertex, and its exit time has no atom at a fixed positive truncation time. The maximum jump divided by $\sqrt n$ tends to zero, so the discrete overshoot vanishes. The fixed displacement of the $S$ boundary also vanishes.

On a compact set of rescaled exit positions, the free local theorem \eqref{cps:amp:eq:freelocal} applies uniformly, with remaining durations in $[\varepsilon n,n]$. After multiplication by $n$, its limit is twice the appropriate free Brownian density at the remaining time. Uniformity for the varying time ratios follows by compactness of $[\varepsilon,1]$ and the local theorem, or by taking convergent subsequences of these ratios. Outside a compact set of exit positions, \eqref{cps:amp:eq:freeupper} bounds the rescaled integrand by $C/\varepsilon$, and functional-limit tightness of the path maximum removes that tail. Thus the truncated expectation converges to the corresponding Brownian exit integral.

It remains to control late exits. Reverse every path that ends at $b$ and exits after time $(1-\varepsilon)n$. Its reversal visits the complement within its first $\varepsilon n$ steps. Strong Markov at the first reversed exit, followed by the global free upper bound for at least $(1-\varepsilon)n$ remaining steps, gives
\begin{equation}\label{cps:amp:eq:late}
 \PP_a(X_n=b,(1-\varepsilon)n<\tau\le n)
 \le\frac Cn\PP_b^*(\tau^*\le\varepsilon n).
\end{equation}
If $b/\sqrt n\in F$ and $F$ has boundary distance at least $\delta>0$, the reversed corrected-martingale exponential inequality gives, for each fixed $\varepsilon$ and sufficiently large $n$,
\begin{equation}\label{cps:amp:eq:earlyexit}
 \PP_b^*(\tau^*\le\varepsilon n)
 \le C\exp[-c\delta^2/\varepsilon].
\end{equation}
Indeed, an exit requires a displacement of at least $\delta\sqrt n-O(1)$, and the optimized moment-generating parameter is of order $1/(\varepsilon\sqrt n)$, within its allowed range for large $n$. Hence the late-exit mass times $n$ tends uniformly to zero as $\varepsilon\downarrow0$. The Brownian analogue removes the truncation in the limiting integral.

Subtracting the exit integral from the limiting free density in \eqref{cps:amp:eq:Hunt} is the Brownian Hunt formula and gives $2k_1^{C,\Sigma}$. The argument with $m_n/n\to t$ uses the identical truncation, with remaining time ratios bounded away from zero; it gives $2k_t^{C,\Sigma}$. Finally, compact uniformity follows by the subsequence principle: any sequence of endpoint pairs in $E\times F$ has a convergent subsequence, and a further subsequence fixes the finite parity pair. All the bounds above are uniform on those compact sets, and the limiting heat kernel is continuous there.
\end{proof}
Only an unrestricted functional limit theorem and the already verified free local theorem enter this proof. No cone-killed local theorem for the two-state chain is assumed.

\section{Fixed-thirds convergence and positive bridge amplitudes}\label{cps:amp:sec:bridge}
Fix relevant endpoints $a,b$, put $m=\lfloor n/3\rfloor$ and $\ell=n-2m$, and define finite measures
\begin{equation}\label{cps:amp:eq:measures}
 \mu_{n,a}=n^{p/2}\sum_uq_m(a,u)\delta_{u/\sqrt n},\qquad
 \mu_{n,b}^*=n^{p/2}\sum_vq_m^*(b,v)\delta_{v/\sqrt n}.
\end{equation}
Let $\nu_{1/3,\Sigma}$ be the image of $\nu_\Sigma$ under $x\mapsto x/\sqrt3$. By Theorem~\ref{cps:amp:thm:exacttime}, including convergence of total mass,
\begin{equation}\label{cps:amp:eq:measurelimit}
 \mu_{n,a}\Longrightarrow3^{p/2}A(a)\nu_{1/3,\Sigma},\qquad
 \mu_{n,b}^*\Longrightarrow3^{p/2}A^*(b)\nu_{1/3,\Sigma}.
\end{equation}
For this model define the positive continuum constant
\begin{equation}\label{cps:amp:eq:B}
 \mathcal B_\Sigma
 =2\,3^p\iint_{C\times C}k_{1/3}^{C,\Sigma}(x,y)
       \,\nu_{1/3,\Sigma}(dx)\nu_{1/3,\Sigma}(dy).
\end{equation}
It is finite by the unrestricted Gaussian bound and strictly positive because the killed heat kernel and the meander density are positive in the interior.

\begin{theorem}[Fixed-endpoint amplitude]\label{cps:amp:thm:bridge}
For each relevant pair of fixed endpoints,
\begin{equation}\label{cps:amp:eq:bridgeamp}
 q_n(a,b)\sim\mathcal C(a,b)n^{-p-1},\qquad
 \mathcal C(a,b)=\mathcal B_\Sigma A(a)A^*(b)>0.
\end{equation}
\end{theorem}
\begin{proof}
The exact semigroup identity and reversal give
\begin{equation}\label{cps:amp:eq:semigroup}
 n^{p+1}q_n(a,b)
 =\iint nq_\ell(u,v)\,
 \mu_{n,a}(d(u/\sqrt n))\mu_{n,b}^*(d(v/\sqrt n)).
\end{equation}
On compact interior sets, Lemma~\ref{cps:amp:lem:killed} makes the kernel converge uniformly to $2k_{1/3}^{C,\Sigma}$. The global bound $nq_\ell\le C$ holds everywhere by \eqref{cps:amp:eq:freeupper}. The finite measures in \eqref{cps:amp:eq:measurelimit} are tight and have convergent total masses. Their limits put no mass on the boundary and none at infinity. Therefore compact interior sets, chosen with limiting boundary measure zero, leave arbitrarily small mass in the complement, uniformly for large $n$. The contribution to \eqref{cps:amp:eq:semigroup} from either complementary coordinate is bounded by that small mass times $C$ and the bounded total mass of the other factor. On the compact product, weak product-measure convergence and uniform kernel convergence give the integral in \eqref{cps:amp:eq:B}, multiplied by $A(a)A^*(b)$. This proves \eqref{cps:amp:eq:bridgeamp}.
\end{proof}
The argument uses a fixed positive fraction of the time in all three pieces. It needs neither a marked-clock local theorem in three dimensions nor a shrinking-time gluing limit.

\begin{writenotelog}
This is the fixed-endpoint killed local limit that Report 112 (Part~\ref{cps:part:logcone},
Section~\ref{cps:log:sec:literature}) names as the stronger result its method cannot give, here
proved for the two models.
\end{writenotelog}

\subsection{Evaluation of the universal Brownian factor}\label{cps:amp:sec:Brownian}
The positive integral \eqref{cps:amp:eq:B} is the primary definition of the bridge factor. Its universal part can also be evaluated explicitly, while the discrete harmonic values remain unevaluated. All Brownian densities below use generator $\tfrac12\Sigma:\nabla^2$ and physical planar Lebesgue measure. Write
\begin{equation}
 h(x)=u(\Sigma^{-1/2}x),\qquad
 C_B=\frac{2^{1-p}}{\theta\Gamma(p+1)\sqrt{\det\Sigma}},\qquad
 \beta_B=\frac{2^{1-p/2}\Gamma(p/2)}{\pi\Gamma(p)}.
\end{equation}
The covariance-identity wedge kernel, obtained from the spectral/Bessel formula of B\~anuelos--Smits \cite[Lemma 1, equation (2.2)]{BS}, is
\begin{equation}\label{cps:amp:eq:wedgekernel}
 k_t^K((r,\varphi),(s,\psi))=
 \frac2{\theta t}e^{-(r^2+s^2)/(2t)}
 \sum_{j\ge1}I_{jp}(rs/t)\sin(jp\varphi)\sin(jp\psi).
\end{equation}
Its leading term uses $I_p(z)\sim(z/2)^p/\Gamma(p+1)$. After the physical coordinate change, the small-start entrance density, divided by the starting harmonic value, is
\begin{equation}\label{cps:amp:eq:entrance}
 g_t(y)=C_Bt^{-p-1}h(y)
       e^{-y^T\Sigma^{-1}y/(2t)}.
\end{equation}
The angular and radial integrals are
\begin{equation}
 \int_0^\theta\sin(p\varphi)\,d\varphi=\frac2p,\qquad
 \int_0^\infty r^{p+1}e^{-r^2/(2t)}\,dr
 =2^{p/2}t^{p/2+1}\Gamma(p/2+1).
\end{equation}
Since $\theta p=\pi$, integration of \eqref{cps:amp:eq:entrance} gives
$\int_Cg_t(y)\,dy=\beta_Bt^{-p/2}$. Thus
\begin{equation}\label{cps:amp:eq:meanderdensity}
 \nu_{t,\Sigma}(dy)=\frac{g_t(y)\,dy}{\beta_Bt^{-p/2}}.
\end{equation}
With the harmonic normalization \eqref{cps:amp:eq:u}, the ordinary iid survival constant is this Brownian constant, so $\beta=\beta^*=\beta_B$.

Define the physical harmonic functions, with the parity of their argument selecting the relevant cycle law, by
\begin{equation}\label{cps:amp:eq:U}
 U_i(a)=2^{p/2}\mathcal H_i((2\Sigma)^{-1/2}(a+c_M)),\qquad
 U_i^*(a)=2^{p/2}\mathcal H_i^*((2\Sigma)^{-1/2}(a+c_M)).
\end{equation}
Then $A=\beta_BU$ and $A^*=\beta_BU^*$. They are harmonic for the respective killed counted-time kernels: apply the one-step survival recursion, multiply by $n^{p/2}$, and pass to the limit using \eqref{cps:amp:eq:survivalupper} and the exponential one-step moment as an integrable majorant. The bounds $V_b\le\mathcal H\le V$ and the fixed-shift comparison give $U(a)\sim h(a)$ as the boundary distance tends to infinity. No uniqueness assertion is required.

The entrance densities satisfy $g_sK_t=g_{s+t}$, by the killed semigroup and its small-start limit. This also follows from \eqref{cps:amp:eq:wedgekernel}, angular orthogonality and the termwise Bessel identity
\begin{equation}
 \int_0^\infty r^{p+1}e^{-ar^2}I_p(br)\,dr
 =\frac{b^p e^{b^2/(4a)}}{(2a)^{p+1}},\qquad a>0.
\end{equation}
For $s=1/3$, each limiting measure in \eqref{cps:amp:eq:measurelimit} is $U(a)g_s(x)\,dx$. The limiting bridge integral therefore reduces to
\begin{equation}
 2U(a)U^*(b)\int_C g_{2s}(y)g_s(y)\,dy.
\end{equation}
A polar integration, using $\int_0^\theta\sin^2(p\varphi)d\varphi=\theta/2$, gives
\begin{align}
 \int_Cg_{2s}(y)g_s(y)\,dy
 &=C_B^2\sqrt{\det\Sigma}\frac{\theta\Gamma(p+1)}4
   (2s)^{-p-1}s^{-p-1}(4s/3)^{p+1}\notag\\
 &=C_B(3s)^{-p-1}=C_B.
\end{align}
Consequently the equivalent explicit form of \eqref{cps:amp:eq:bridgeamp} is
\begin{align}\label{cps:amp:eq:explicitbridge}
 \mathcal C(a,b)&=2C_BU(a)U^*(b)\notag\\
 &=\frac{4}{\theta\Gamma(p+1)\sqrt{\det\Sigma}}
 \mathcal H_i(T(a+c_M))\mathcal H_j^*(T(b+c_M)).
\end{align}
The factor two in the first line is the discrete lattice covolume; it is not part of the Brownian density. Here $i,j$ are the initial and terminal parities. In particular, the $S$ bridge uses the even forward and odd reversed harmonic functions. The determinants are
\begin{equation}
 \det\Sigma_P=7,\qquad \det\Sigma_S=80/9.
\end{equation}
As normalization checks, for identity covariance and $p=1$ the constants are $C_B=1/\pi$, $\beta_B=\sqrt{2/\pi}$; for $p=2$, they are $C_B=1/(2\pi)$, $\beta_B=1/\pi$, with $u(x,y)=2xy$. These formulas evaluate only the universal Brownian factors. The discrete harmonic values, and the weighted endpoint sum for $P$, still determine the counting amplitudes.

\section{Enumeration amplitudes and the rigid sequence}\label{cps:amp:sec:enumeration}
Attach subscripts $P,S$ to the constants for the two models.

\subsection{Corner polyhedra}
By \eqref{cps:amp:eq:Pcount}, Theorem~\ref{cps:amp:thm:bridge}, and the uniform bound \eqref{cps:amp:eq:endpointupper}, dominated convergence in the geometrically weighted endpoint sum gives
\begin{equation}\label{cps:amp:eq:kappaP}
 \kappa_P=\sum_{i\ge1}3^{-i}\mathcal C_P((0,0),(2i,0))
 =\mathcal B_{\Sigma_P}A_P((0,0))
   \sum_{i\ge1}3^{-i}A_P^*((2i,0)).
\end{equation}
The dominating series $\sum_{i\ge1}3^{-i}(1+i)^{p_P}$ converges. Positivity follows already from the single endpoint $(4,0)$. This proves \eqref{cps:amp:eq:Pmain}.

\subsection{Schnyder labelings}
The fixed-endpoint identity \eqref{cps:amp:eq:Scount} gives
\begin{equation}\label{cps:amp:eq:kappaSprime}
 s'_n\sim\kappa'_S\mu_S^nn^{-\alpha_S},\qquad
 \kappa'_S=\frac{\mu_S}{4}\mathcal C_S((0,2),(1,1)).
\end{equation}
Its positive finite combination then gives
\begin{equation}\label{cps:amp:eq:kappaS}
 \kappa_S=(1+\mu_S^{-1})^2\kappa'_S
 =\frac{361}{192}\mathcal B_{\Sigma_S}A_S((0,2))A_S^*((1,1)).
\end{equation}
Here $\mu_S/4=4/3$ and $(1+\mu_S^{-1})^2=(19/16)^2$. This proves \eqref{cps:amp:eq:Smain}.

\subsection{A dominated signed convolution}
Write
\begin{equation}
 (1+z)^{-3}=\sum_{j\ge0}b_jz^j,\qquad
 b_j=(-1)^j\binom{j+2}{2}.
\end{equation}
From \eqref{cps:amp:eq:series}, $\widetilde s_n$ is the signed convolution of $s_n$ with $b_j$, minus the coefficient of $(3z^2+z^3)/(1+z)^3$. Enlarge a constant so that
\begin{equation}
 0\le s_m\le C\mu_S^m(m+1)^{-\alpha_S}\qquad(m\ge0).
\end{equation}
For $0\le j\le n$,
\begin{equation}
 \frac{n}{n-j+1}\le j+1,
\end{equation}
so the absolute normalized convolution summand is bounded by
\begin{equation}\label{cps:amp:eq:convolutionmajorant}
 n^{\alpha_S}\mu_S^{-n}|b_j|s_{n-j}
 \le C|b_j|\mu_S^{-j}(j+1)^{\alpha_S}.
\end{equation}
The right side is summable over $j$, independently of $n$. For each fixed $j$, the summand has signed limit $\kappa_Sb_j\mu_S^{-j}$. Dominated convergence therefore gives
\begin{equation}
 \lim_{n\to\infty}n^{\alpha_S}\mu_S^{-n}
 \sum_{j=0}^nb_js_{n-j}
 =\kappa_S\sum_{j\ge0}b_j\mu_S^{-j}
 =\frac{\kappa_S}{(1+\mu_S^{-1})^3}.
\end{equation}
The polynomial correction divided by $(1+z)^3$ has coefficients $O(n^2)$, so it is exponentially negligible under this normalization. Since $(1+\mu_S^{-1})^{-3}=(16/19)^3$, \eqref{cps:amp:eq:rigidmain} follows. This completes the proof of Theorem~\ref{cps:amp:thm:main}.

\section{Sharper first-threshold inversion}\label{cps:amp:sec:inverse}
The following conclusion applies to each of the three counting sequences, with its rate $\mu$, exponent $\alpha$ and amplitude $\gamma$ from Theorem~\ref{cps:amp:thm:main}.

\begin{theorem}\label{cps:amp:thm:inverse}
Suppose $a_n\sim\gamma\mu^nn^{-\alpha}$ with $\gamma>0$, $\mu>1$ and $\alpha>0$, and define $N(Y)=\min\{n\ge0:a_n\ge Y\}$. Put $\lambda=\log\mu$, $T_Y=\log Y$, and
\begin{equation}\label{cps:amp:eq:W}
 r(Y)=-\frac\alpha\lambda W_{-1}\!\left(-\frac\lambda\alpha(\gamma/Y)^{1/\alpha}\right).
\end{equation}
Then
\begin{equation}\label{cps:amp:eq:inverseexpansion}
 r(Y)=\frac{T_Y+\alpha\log T_Y-\alpha\log\lambda-\log\gamma}{\lambda}+o(1).
\end{equation}
There is a positive function $\varepsilon_Y\to0$ such that, for all sufficiently large $Y$,
\begin{equation}\label{cps:amp:eq:ceiling}
 \lceil r(Y)-\varepsilon_Y\rceil\le N(Y)\le\lceil r(Y)+\varepsilon_Y\rceil.
\end{equation}
No separate monotonicity assumption on $a_n$ is needed.
\end{theorem}
\begin{proof}
The defining relation $W(z)e^{W(z)}=z$ gives, on the negative real branch $W_{-1}$,
\begin{equation}\label{cps:amp:eq:modelinverse}
 \gamma e^{\lambda r(Y)}r(Y)^{-\alpha}=Y.
\end{equation}
The branch selects the large positive solution tending to infinity; the other real branch selects the small solution. To obtain \eqref{cps:amp:eq:inverseexpansion}, first $r(Y)\sim T_Y/\lambda$ by taking logarithms in \eqref{cps:amp:eq:modelinverse}; substituting this into
$\lambda r=T_Y+\alpha\log r-\log\gamma$ gives the stated expansion, with $o(1)$ remainder after one substitution.

Write $\log a_n=g(n)+e_n$ for all sufficiently large $n$, where
\begin{equation}
 g(t)=\lambda t-\alpha\log t+\log\gamma,\qquad e_n\to0.
\end{equation}
The coefficients are eventually positive by the full equivalent. For large $t$, $g'(t)\ge\lambda/2$. Put
$\delta_Y=\sup_{n\ge r(Y)/2}|e_n|$, where the supremum is over integers in that range. Then $\delta_Y\to0$. Choose $\varepsilon_Y>2\delta_Y/\lambda$ with $\varepsilon_Y\to0$, for instance $3\delta_Y/\lambda+1/r(Y)$.

Every integer $n$ with $r(Y)/2\le n<r(Y)-\varepsilon_Y$ satisfies
$g(n)+e_n<g(r(Y))=\log Y$. Conversely the integer
$n=\lceil r(Y)+\varepsilon_Y\rceil$ satisfies $g(n)+e_n>\log Y$. Integers below $r(Y)/2$ remain below threshold for large $Y$ by the exponential-polynomial upper bound and eventual increase of its comparison envelope; finitely many initial values cause no exception for large $Y$. This proves the two ceiling bounds.
\end{proof}
The conclusion is deliberately a bracket. There is no unconditional assertion that $N(Y)=\lceil r(Y)\rceil$, because $r(Y)$ can approach an integer more closely than the unknown forward relative error permits. Nor do additional formal terms of the Lambert-$W$ model provide a justified error expansion for the actual counting sequence without a quantitative remainder in its forward asymptotic. No numerical amplitude estimate is substituted into \eqref{cps:amp:eq:W}.

\begin{writenote}
Theorem~\ref{cps:amp:thm:inverse} is an instance of the transseries volume's
\texttt{p0:thm:lambert-core} ($a=\lambda$, $b=-\alpha$, branch $W_{-1}$), with
\eqref{cps:amp:eq:inverseexpansion} the first step of \texttt{p0:thm:lambert-centered} and the
bracket \eqref{cps:amp:eq:ceiling} the separation situation of \texttt{p0:thm:staircase}
(\repo{Analysis/Transseries/docs/series-and-transseries/Transseries_And_Inversion/}). No
novelty is claimed for the inversion mechanics. The generic staircase separation lemma is
formalized in Lean as \texttt{Fabius.staircase\_separation}; nothing about $p_n$, $s_n$ or
$\widetilde s_n$ is formalized.
\end{writenote}

\section{Verification, reproducibility and limits of the result}
The accompanying package contains the addendum source and PDF, the unchanged 30-file Foundation package, portable build commands, exact symbolic checks of the new algebraic amplitude factors and inverse identities, source and mathematical-review receipts, and checksums. The supplied integrity program verifies every included Foundation file against its original manifest before any reproduction command is run.

\begin{writenote}
In this repository the addendum PDF, the embedded Foundation copy, the integrity
receipt \texttt{results/foundation-integrity.json} (a hash ledger of that copy) and the checksum
ledgers are not shipped, so \texttt{code/46-amp-verify\_foundation.py} and
\texttt{code/46-amp-run\_all.py} cannot run from the report directory. The report README says how
to rerun them from the archive in the arrival commit, and how to run the new algebra check
\texttt{code/46-amp-verify\_addendum.py} alone.
\end{writenote}

Finite computations verify algebraic identities, normalization factors and indexing. They do not establish Lemma~\ref{cps:amp:lem:shift}, the valid-cycle harmonic sandwich, the functional-limit and exit-decomposition argument, or the limiting interchange in the semigroup identity; those are the mathematical proofs given above. The positive amplitude formulas \eqref{cps:amp:eq:kappaP} and \eqref{cps:amp:eq:kappaS} characterize the constants through killed harmonic functions and Brownian heat kernels. They do not certify numerical amplitude digits, a rate of convergence, or any unproved correction term.


\clearpage
\part{A log-scale cone theorem for state-modulated walks (source 47)}\label{cps:part:logcone}
\begin{center}\large Logarithmic asymptotics for polyhedral orientations and Schnyder labelings\\[2pt]
\normalsize Research report 112, 2 October 2026\end{center}
\begin{sourceabstract}{source 47 (bundle Report 112)}
This report treats polyhedral orientations counted by OEIS A377922 and
3-connected Schnyder labelings counted by A377920. It gives a finite-phase
cone argument with exponentially decaying jump tails, then verifies the exact
walk models, counting time, boundary conditions, covariance and endpoint
normalizations needed for the two applications. The conclusions are
\[
 p_n=(9/2)^n n^{-\ap+o(1)},\qquad
 s_n=(16/3)^n n^{-\as+o(1)},
\]
where $\ap=1+\pi/\arccos(9/16)$ and
$\as=1+\pi/\arccos(22/27)$. The growth constants were already proved by
Fusy, Narmanli and Schaeffer in 2023. The advance established here is the
logarithmic polynomial power, with non-D-finiteness and two-term
first-threshold inverses. An explicit six-face sleeve injection transfers the all-index
logarithmic power, non-D-finiteness and inverse expansion to A377921. Neither a
leading amplitude nor a full multiplicative equivalent is obtained.
\end{sourceabstract}

\begin{writenotelog}
Report 112 was written without knowledge of Parts~\ref{cps:part:exponents}
and~\ref{cps:part:amplitudes}, which prove more about the three sequences: $\Theta$ bounds
(Theorem~\ref{cps:thm:main}), full equivalents with positive amplitudes
(Theorem~\ref{cps:amp:thm:main}) and inverses with $O(1)$ error and a next constant. Its
sentence ``the advance established here is the logarithmic polynomial power'' is therefore
true only of Report 112 alone. What Part~\ref{cps:part:logcone} adds to the report is listed in
Section~\ref{cps:log:sec:guide}: chiefly a general cone theorem, Theorem~\ref{cps:log:thm:cone},
and a six-face sleeve injection, Proposition~\ref{cps:log:prop:sleeve}.
\end{writenotelog}

\section{Guide to Part III}\label{cps:log:sec:guide}
This section was written when bundle Report 112 was added as Part~\ref{cps:part:logcone}
(5 October 2026). Everything after it, up to Appendix~\ref{cps:app:checks}, is Report 112 as
delivered, apart from labels and the changes listed below, with dated notes marked
\textbf{[write, 2026-10-05]}. In Part~\ref{cps:part:logcone}, ``this report'' and ``the
source'' are Report 112's words: the first means Report 112, the second the paper of
Fusy, Narmanli and Schaeffer \cite{FNS}. An independent check of the write, made afterwards, is
recorded at the end of Section~\ref{cps:log:sec:research}.

\subsection{What is new, and what is a second route}\label{cps:log:sec:new}
Part~\ref{cps:part:logcone} proves no asymptotic statement about A377922, A377920 or A377921
that Parts~\ref{cps:part:exponents}--\ref{cps:part:amplitudes} do not already prove in a
stronger form. Its statements about the three sequences are a second route to
Part~\ref{cps:part:exponents}, by a different method and with weaker conclusions (logarithmic
scale, $n^{o(1)}$ instead of $\Theta$ or $\sim$); Table~\ref{cps:log:tab:map} maps each of them to
its counterpart. New to the report are:
\begin{enumerate}[label=(\roman*)]
\item \emph{Theorem~\ref{cps:log:thm:cone}}, a cone theorem for a Markov-additive chain on
 $\ZZ^2\times E$ with a finite internal state (phase) set $E$ and unbounded steps with a uniform
 exponential moment, killed on leaving the closed quadrant: under hypotheses (H1)--(H5),
 $K_n((0,a),(0,b))=n^{-1-\pi/\theta+o(1)}$ at every integer time, and under (H1)--(H4) the
 survival bound $\PP(\tau>n)\le C_\epsilon n^{-\nu/2+\epsilon}$. Part~\ref{cps:part:exponents}
 states that it ``does not posit a new general cone theorem for arbitrary Markov-modulated
 walks''; Theorem~\ref{cps:log:thm:cone} is such a theorem, at the logarithmic scale only.
 Lemma~\ref{cps:log:lem:bridge}, its exact-time interior bridge bound, is the general form of
 Lemma~\ref{cps:lem:killed}.
\item \emph{Proposition~\ref{cps:log:prop:sleeve}}, an explicit injection from Schnyder
 labelings of size $n$ into those of size $n+6$ whose outer white vertices are non-isolated (a
 six-face annulus around the outer hexagon), hence $s_{n-6}\le b_n\le s_n$ for $n\ge8$. Neither
 Parts~\ref{cps:part:exponents}--\ref{cps:part:amplitudes} nor \cite{FNS} as cited contain it.
\item \emph{A non-D-finiteness criterion with a weaker hypothesis} (Section~\ref{cps:log:sec:nonD}
 and Remark~\ref{cps:log:rem:criterion}): nonnegative integer coefficients with
 $a_n=\Gamma^nn^{-\alpha+o(1)}$, $\Gamma>1$ and $\alpha$ irrational already exclude
 D-finiteness. Lemma~\ref{cps:lem:nonD} needs $a_n=\Theta(\mu^nn^{-\alpha})$.
\item \emph{The fixed-endpoint sandwich} $e_{n-1}\le p_n\le e_{n+1}$ of
 \eqref{cps:log:eq:Psandwich}, with $e_L$ the number of legal corner paths of length $L$ from
 $(0,0)$ to $(1,1)$, by two explicit injections. It replaces the tilted endpoint sum
 \eqref{cps:eq:Pcount} by one fixed endpoint of tilt one.
\item \emph{An explicit exponential-moment certificate} for a whole Schnyder aggregate,
 \eqref{cps:log:eq:Smoment}: at $e^t=6/5$ the moment of the total elementary length is $189/13$.
 Part~\ref{cps:part:exponents} obtains the moment from analyticity of the return generating
 function \eqref{cps:eq:return}. (The delivered checks also certify the corner quotient
 moment $524/405$ at the same point; see Section~\ref{cps:log:sec:repro}.)
\item \emph{Exact bounds} $4<\ap<5$ and $6<\as<7$ (Section~\ref{cps:log:sec:nonD});
 Parts~\ref{cps:part:exponents}--\ref{cps:part:amplitudes} give decimals only.
\item \emph{Three references} not cited in Parts~\ref{cps:part:exponents}--\ref{cps:part:amplitudes}:
 Denisov--Zhang \cite{DZ}, Grama--Lauvergnat--Le Page \cite{GLLP} and Pham--Peign\'e--Son
 \cite{PPS} (Section~\ref{cps:log:sec:literature}), together with Whitt's martingale functional
 limit theorem \cite{Whitt}, which Theorem~\ref{cps:log:thm:cone} uses.
\end{enumerate}
Also new are several derivations of known quantities: a probabilistic reconstruction of the
corner kernel, direct second-moment computations of both covariances, and a second proof of the
aperiodicity Lemma~\ref{cps:lem:aperiodic} with new seed paths. Report 112's further questions
(Section~\ref{cps:log:sec:research}) are mostly answered by Part~\ref{cps:part:amplitudes};
dated notes there say which.

\begin{table}[htbp]
\centering\footnotesize
\caption{Report 112 mapped to Parts~I--II. ``Second route'' means a different proof of a
conclusion that Parts~I--II already contain, here with the strength stated.}\label{cps:log:tab:map}
\begin{tabular}{@{}>{\raggedright}p{0.25\textwidth}>{\raggedright}p{0.14\textwidth}>{\raggedright}p{0.32\textwidth}>{\raggedright\arraybackslash}p{0.19\textwidth}@{}}
\toprule
Report 112 & Here & Parts I--II & Class\\
\midrule
Thm 1.1: $p_n$, $s_n$ on the log scale & Thm~\ref{cps:log:thm:main} & $\Theta$: Thm~\ref{cps:thm:main}, \eqref{cps:eq:loglimits}; $\sim$: Thm~\ref{cps:amp:thm:main} & second route, weaker\\
Thm 1.1, Cor 7.2: $b_n$ on the log scale & \eqref{cps:log:eq:mainb}, Cor~\ref{cps:log:cor:921} & $\widetilde s_n\sim(16/19)^3\kappa_S\mu_S^nn^{-\alpha_S}$, \eqref{cps:amp:eq:rigidmain} & second route, weaker; needs no amplitude\\
Thm 1.1: three GFs not D-finite & Thm~\ref{cps:log:thm:main}, Secs~\ref{cps:log:sec:nonD}, \ref{cps:log:sec:921} & Cor~\ref{cps:cor:nonD} & same theorem; criterion new (iii)\\
Thm 1.1, Sec 6: inverse, $o(\log\log y)$ & \eqref{cps:log:eq:inverse-main}, Sec~\ref{cps:log:sec:inverse} & $O(1)$: Lemma~\ref{cps:lem:inverse}; next constant: Thm~\ref{cps:amp:thm:inverse} & weaker\\
Thm 2.1, general cone theorem & Thm~\ref{cps:log:thm:cone} & disclaimed (Part I, Sec.~1) & new (i)\\
Lemma 2.2, interior bridge & Lemma~\ref{cps:log:lem:bridge} & Lemma~\ref{cps:lem:killed} (two models) & general form\\
Secs 2.2, 2.4--2.5: FCLT, LLT, annular crossings, first-hit gluing & Secs~\ref{cps:log:sec:limits}, \ref{cps:log:sec:crossings}--\ref{cps:log:sec:lower} & Prop~\ref{cps:prop:local}; Lemma~\ref{cps:amp:lem:killed} uses an FCLT; Part I glues by regeneration with a cycle buffer & general; crossing method new to the report\\
Quotient kernels $P(u,v)$, $G(u,v)$; corner reconstruction & \eqref{cps:log:eq:Pkernel}, \eqref{cps:log:eq:Gkernel} & $K_P$, $K_S$: \eqref{cps:eq:PK}, \eqref{cps:eq:SM} & same kernels in half-scale coordinates; new derivation\\
$\Sigma_P$, $\Sigma_S$, drifts, correctors, direct moments & \eqref{cps:log:eq:Psigma}, \eqref{cps:log:eq:Ssigma} & \eqref{cps:eq:Sigma}, \eqref{cps:eq:drifts} & same, divided by 4 (Remark~\ref{cps:log:rem:scale})\\
(H4) aperiodicity, (H5) seeds & Secs~\ref{cps:log:sec:corner}, \ref{cps:log:sec:schnyder} & Lemma~\ref{cps:lem:aperiodic}; Sec.~\ref{cps:sec:entrances} & second proof; different seeds\\
$e_{n-1}\le p_n\le e_{n+1}$ & \eqref{cps:log:eq:Psandwich} & tilted sum \eqref{cps:eq:Pcount} & new (iv)\\
Aggregates, shifted quadrant, endpoint factor 4 & Sec.~\ref{cps:log:sec:schnyder}, \eqref{cps:log:eq:Sendpoint} & Lemma~\ref{cps:lem:Sboundary}, \eqref{cps:eq:Scount} & same, second proof\\
Moment $189/13$ & \eqref{cps:log:eq:Smoment} & analyticity of \eqref{cps:eq:return} & new (v)\\
Irrationality; $4<\ap<5$, $6<\as<7$ & Sec.~\ref{cps:log:sec:nonD} & Lemma~\ref{cps:lem:irrational}; decimals & same proof; bounds new (vi)\\
Prop 7.1, sleeve injection & Prop~\ref{cps:log:prop:sleeve} & none & new (ii)\\
Sec 7.2: convolution relation & \eqref{cps:log:eq:921gf}, \eqref{cps:log:eq:921conv} & \eqref{cps:eq:rigid} & same identity (Remark~\ref{cps:log:rem:rigid})\\
Sec 8: phase covariances $V^P_c$, $V^S_c$ & Sec.~\ref{cps:log:sec:literature} & \eqref{cps:eq:PC}, \eqref{cps:eq:SC} and the remark after \eqref{cps:eq:Sigma} & same, divided by 4\\
Sec 8: literature & Sec.~\ref{cps:log:sec:literature} & \cite{DWrevisited} only & three references new (vii)\\
Sec 9: questions & Sec.~\ref{cps:log:sec:research} & answered in part by Part II & dated notes\\
\bottomrule
\end{tabular}
\end{table}

\subsection{Provenance}\label{cps:log:sec:provenance}
Part~\ref{cps:part:logcone} is bundle Report 112. Its archive is
\texttt{Geometric\_OEIS\_\allowbreak Sequences\_Cone\_\allowbreak Exponents\_and\_\allowbreak Inverses\_Source.zip}
(wrapper \texttt{report112/}, main file \texttt{report112.tex}, a 19-page PDF). It arrived in
commit \texttt{60f54ea06} and was placed as source 47 of this report, manuscript 01 of
cluster 103-A377922 of batch 103, in \texttt{9c995cefe}. Its author line and PDF metadata say
``Research report 112''; it names no tool and no person. It is dated 2 October 2026 and names no
ProveIt commit, path or search: there is no pin. It was written about eight hours after
sources 45 and 46 (their literature snapshots are 01:22--01:39 UTC on 2 October; the archive was
created at 09:26 UTC) in a different session, and it shows no knowledge of them. Its text
carries no ``prepared for private review'' line and no personal data. Its programs, fixtures and
recorded outputs are shipped byte-identical under the prefix \texttt{47-logcone-}; the
manuscript, its PDF, its README and its checksum inventory \texttt{manifest.json} are not
shipped (all in the arrival commit). The report README lists the files and gives a rerun
route.

\subsection{How Report 112 is printed}\label{cps:log:sec:printing}
\begin{itemize}
\item Report 112 is printed in full, with its abstract. Its Section $k$ is Section $k+22$ here,
 and its Theorem, Lemma, Proposition or Corollary $k.m$ is $(k+22).m$; so its Theorem 2.1 is
 Theorem~\ref{cps:log:thm:cone} and its Proposition 7.1 is
 Proposition~\ref{cps:log:prop:sleeve}. Its equations, numbered consecutively as delivered, are
 numbered within sections here.
\item Its statements about the three sequences are printed, not replaced by pointers, because
 they are where Theorem~\ref{cps:log:thm:cone}'s hypotheses are verified for the two models, and
 because the new items (iv) and (v) sit inside them. The dated notes and
 Table~\ref{cps:log:tab:map} mark them as second routes.
\item Labels carry the prefix \texttt{cps:log:}. Report 112's macros are mapped to this report's
 ($\mathbb E$, $\mathbb R$, $\mathbb Z$ by \verb|\EE|, \verb|\RR|, \verb|\ZZ|); its title block,
 table of contents and running heads are omitted. No statement, proof, symbol or number was
 changed.
\item Bibliography: Report 112's entries for Fusy--Narmanli--Schaeffer and Fischler--Rivoal are
 the same papers as \cite{FNS} and \cite{FR}; its \texttt{DW} is \cite{DWrevisited}. Its entries
 for Whitt, Denisov--Zhang, Grama--Lauvergnat--Le Page and Pham--Peign\'e--Son are added. Its
 uncited OEIS entry (A377922, A377920, A377921, A377923) duplicates \cite{OEISP,OEISS}, except
 for the excluded A377923, and is omitted.
\item Report 112 cites \cite[Proposition 19]{FNS} for the shift relation
 \eqref{cps:log:eq:Sshift}, which Part~\ref{cps:part:exponents} cites as Proposition 14 and
 Remark 9. Both are right: in arXiv:2202.09172v3 Remark 9 gives the refined relation by numbers
 of white and black vertices, and Proposition 19 the relation \eqref{cps:eq:Sconversion} for
 $n\ge4$ (checked 5 October 2026). Report 112's form \eqref{cps:log:eq:921gf} of Remark 20 is the
 one printed there.
\end{itemize}

\subsection{Notation}\label{cps:log:sec:notation}
Report 112 works in \emph{half-scale} (quotient) coordinates: a point of the parity lattice
$L$ is written $x=2U+c$, $y=2V+c$ for corner polyhedra \eqref{cps:log:eq:cornerquot} and
$x=2U+c$, $y=2V+2-c$ for Schnyder labelings \eqref{cps:log:eq:Squot}, with phase
$c\in\{0,1\}$ equal to the parity of $y$. Phase $0$ is Parts~I--II's even state $\e$, phase $1$
the odd state $\oState$. For corner polyhedra $(U,V)$ are Part~\ref{cps:part:exponents}'s cell
coordinates $(a,b)$ of Section~\ref{cps:sec:Fourier}; for Schnyder labelings the cell coordinates
are $(U,V+1-c)$, because Report 112 shifts the quadrant $Q_S$ to $U,V\ge0$. Time is the same as in Parts~I--II: physical steps for corner
polyhedra, aggregates for Schnyder labelings. The letters are kept; Table~\ref{cps:log:tab:notation}
lists the clashes.

\begin{remark}[write, 2026-10-05: the factor 4]\label{cps:log:rem:scale}
Part~\ref{cps:part:logcone}'s covariances are a quarter of Parts~I--II's:
$\Sigma_P^{\mathrm{III}}=\Sigma_P/4$, $\Sigma_S^{\mathrm{III}}=\Sigma_S/4$, and the phase-conditional
covariances of Section~\ref{cps:log:sec:literature} are $V^P_0=C^P_{\e}/4$, $V^P_1=C^P_{\oState}/4$,
$V^S_0=C^S_{\e}/4$, $V^S_1=C^S_{\oState}/4$ (compare \eqref{cps:log:eq:Psigma},
\eqref{cps:log:eq:Ssigma} with \eqref{cps:eq:Sigma}, \eqref{cps:eq:PC}, \eqref{cps:eq:SC}). The
drifts and correctors correspond as
\[
 d_c=2m_c+(P_{c,1-c})(1-2c)(1,\pm1),\qquad
 h^{\mathrm{I}}_c=2h_c-c(1,\pm1)+\tfrac12(1,\pm1),
\]
with the upper sign for corner polyhedra and the lower for Schnyder labelings, where
$h^{\mathrm{I}}$ is Part~\ref{cps:part:exponents}'s corrector of \eqref{cps:eq:drifts}.
Correlations, the angle $\theta$, $\nu=p=\pi/\theta$ and the exponents are the same.
\end{remark}
\begin{proof}
Write the physical position as $X_n=2S_n+\phi(C_n)$, with $\phi(c)=(c,c)$ for corner polyhedra
and $\phi(c)=(c,2-c)$ for Schnyder labelings, and $S_n=(U_n,V_n)$. A step from phase $c$ to phase
$d$ has physical displacement $2\Delta S+\phi(d)-\phi(c)$, and $\phi(d)-\phi(c)=(d-c)(1,\pm1)$.
Taking the conditional mean, $\EE_c[d-c]=P_{c,1-c}(1-2c)$, which gives the first identity.
If $Z_n=S_n+h(C_n)$ is the martingale of Section~\ref{cps:log:sec:limits}, then
$2Z_n=X_n-\phi(C_n)+2h(C_n)$, so $X_n+h^{\mathrm{I}}(C_n)$ is a martingale for
$h^{\mathrm{I}}=2h-\phi+k$ with any constant $k$; Part~\ref{cps:part:exponents}'s normalization
$h^{\mathrm{I}}_{\oState}=-h^{\mathrm{I}}_{\e}$ fixes $k$, and since $\phi(c)=c(1,\pm1)+(0,1\mp1)$ and
$h_1=-h_0$, this gives the second identity. Its increments are twice the increments of $Z$, so
every conditional and every averaged covariance of the physical corrected increment is four
times the corresponding quarter-scale one. With the numbers: corner polyhedra,
$m_0=(-1/3,-1/3)$, $P_{01}=5/6$, $d_{\e}=(-2/3+5/6)(1,1)=(1/6,1/6)$, $h_0=(-1/5,-1/5)$,
$h^{\mathrm{I}}_{\e}=(-2/5+1/2)(1,1)=(1/10,1/10)$; Schnyder labelings, $m_0=(-1/2,3/8)$,
$P_{01}=7/8$, $d_{\e}=(-1+7/8,\ 3/4-7/8)=(-1/8,-1/8)$, $h_0=(-2/7,3/14)$,
$h^{\mathrm{I}}_{\e}=(-4/7+1/2,\ 3/7-1/2)=(-1/14,-1/14)$, as in \eqref{cps:eq:drifts}. Scaling a
covariance does not change a correlation, hence not $\theta$.
\end{proof}

\begin{table}[htbp]
\centering\footnotesize
\caption{Letters of Part~III that clash with Parts~I--II. ``Reading to avoid'' is the tempting
false identification.}\label{cps:log:tab:notation}
\begin{tabular}{@{}>{$}l<{$}>{\raggedright}p{0.32\textwidth}>{\raggedright}p{0.3\textwidth}>{\raggedright\arraybackslash}p{0.2\textwidth}@{}}
\toprule
\text{Letter} & Part III (source 47) & Parts I--II & Reading to avoid\\
\midrule
\Sigma_P,\Sigma_S & covariance per unit time in half-scale coordinates & the same per physical unit & they are equal; Part III's is a quarter (Remark~\ref{cps:log:rem:scale})\\
V^P_c,V^S_c & phase-conditional covariances (Section~\ref{cps:log:sec:literature}) & $C^P_i$, $C^S_i$, four times larger; Part II's $V$ is a harmonic function & \\
m,h & half-scale drift and corrector, $(I-P)h=m$ \textup{[}corrected at the independent check\textup{]} & $d$, $h$ physical & $h_0=h^{\mathrm I}_{\e}$\\
\Gamma_P,\Gamma_S & growth rates $9/2$, $16/3$ & $\mu_P$, $\mu_S$; Part II's $\Gamma$ is the Gamma function & \\
\nu & $\pi/\theta$ & $p$; Part II's $\nu$ is the meander endpoint law & $\nu$ a measure\\
c,d,a,b & phases $0,1$ ($0=\e$, $1=\oState$); $a,b$ endpoint phases in Thm~\ref{cps:log:thm:cone} & $a,b$ points & \\
e_L & legal corner paths $(0,0)\to(1,1)$ of length $L$ & $\e$ the even state & \\
P & phase matrix $P_{cd}$; corner kernel $P(u,v)$ & $P$ the model; $P=K(1,1)$ & $P(u,v)$ is $K_P$ in half-scale variables, not $K_P(u,v)$\\
p_{cd}(v) & transition probabilities & $p=\pi/\theta$ & \\
D & $(1-z/u)(1-zv)$, Part I's $D$ at $u=x^2$, $v=y^2$; a constant in $L_n=D\log(n+2)$ & Part II: valid-cycle kernel & \\
A & aggregate matrix ($=M_S$); $A(t)=S(z)$; tail constant in \eqref{cps:log:eq:tail} & $A(z)$ in Lemma~\ref{cps:lem:nonD}; Part II amplitudes $A(a)$ & \\
F & face matrix ($=W$); Fourier matrix; $F(z)$ a generating function & Fourier matrix $F(t)$ of cells & \\
J & SE matrix ($=E$); angular interval in (H5); a stage index & the state $J_t$ & \\
W & whitening $\Sigma^{-1/2}$ & Part I: face matrix; Part II whitens with $T=(2\Sigma)^{-1/2}$ & \\
T & corrected increment; time constant $T(q)$ & Part I: $T(z)$; Part II: whitening & \\
G,C & Schnyder half-scale kernel and its scalar $C(u,v)$ \textup{[}corrected at the independent check\textup{]} & $K_S$; Part II's $C$ the open quadrant & \\
K_n & killed kernel & $q_n$; $K_P$, $K_S$ step matrices & \\
b_n,B(t) & A377921 and its series & $\widetilde s_n$, $T(z)=\widetilde S(z)$; Part II's $b_j$, $b_P$, $b_S$ & \\
H & $H(t)=F(Rt)$; Dirichlet solution $H_f$ & Schnyder denominator; Part II's $\mathcal H$ & \\
\kappa & exponential tail rate & Part II amplitudes $\kappa_P,\kappa_S$ & \\
E & finite phase set & SE matrix $E(x,y)$ & \\
L_n,L,R & jump cutoff; face sums; $R=1/\Gamma$, radii & lattice $L$; return function $R_i$ & \\
N(y) & first threshold at level $y$ & $N(Y)$ & \\
\bottomrule
\end{tabular}
\end{table}

\subsection{Reports that use the same method}\label{cps:log:sec:siblings}
Report 112 is one of three bundle reports of 2 October 2026 (archives created 08:51, 09:26
and 14:23 UTC) that run one argument: a Perron tilt to a centred walk with a finite phase, a
bounded Poisson corrector and the martingale functional limit theorem, a Fourier-matrix local
limit theorem, an exact-time interior bridge, Brownian wedge harmonic measure over geometric
annuli, two first-hit prefixes glued at the exact time, then non-D-finiteness by a positive
derivative and the $G$-function theorem, and an $o(\log\log y)$ inverse. Report 111 proves it for
one-sided rectangulations, A348351, with bounded steps; it is the report \repo{a348351-one-sided-rectangulations} in the same collection directory as this one,
placed in \texttt{9c995cefe}. Report 123 proves it for the inversion-sequence class A279571; it
arrived in \texttt{60f54ea06} and is Part~II of the report \repo{a279571-inversion-cone-walk} in the
same collection directory (placed in \texttt{612787fb4}, written in batch~104), beside Report~125 as
Part~I. Reports 111 and 123 argue
model by model; Report 112 alone states the argument as a theorem,
Theorem~\ref{cps:log:thm:cone}. Report 111 proves
$\log a(n)=n\log\Gamma-\alpha\log n+o(\log n)$ for A348351, with $\Gamma=(7+\sqrt{17})/2$, for a
four-colour quadrant walk with bounded steps; it does not cite Report 112, nor Report 112 it.
That report's Remark 1.2 (\texttt{osr:rem:cone}) checks that its walk satisfies (H1)--(H5) of
Theorem~\ref{cps:log:thm:cone} at the phases $(c,\mathsf W)$, $c\in\{\mathsf B,\mathsf R,\mathsf G\}$
(primitive colour matrix, bounded steps, zero drift with positive definite covariance, its own
aperiodicity section, explicit seeds from both endpoints), so its exponent is also an instance of
Theorem~\ref{cps:log:thm:cone}; Report 111 is the earlier, model-specific proof (bundle index
08:51 against 09:26 UTC). Its non-D-finiteness argument is the one of
Section~\ref{cps:log:sec:nonD}, and its Remark 10.1 (\texttt{osr:rem:criterion}) states the
criterion in general, as Remark~\ref{cps:log:rem:criterion} does here. Its irrationality
certificate differs from Lemma~\ref{cps:lem:irrational}: there $2\cos\theta=(-29+7\sqrt{17})/2$
is an algebraic integer, and the certificate is its conjugate, which lies below $-2$. No text is
shared, and the reports stay separate.

\noindent(Dated note, 5~October 2026, batch~104: the sentence on Report~123 above first read ``it
arrived in \texttt{60f54ea06} and has not been placed (5 October 2026)''. Remark~26.1 of
\repo{a279571-inversion-cone-walk} (\texttt{icw:log:rem:seeds}) checks Report~123's walk against
Theorem~\ref{cps:log:thm:cone}. (H1)--(H4) hold, with $\theta=\arccos(2/\sqrt7)$ and
$\nu=\pi/\theta\in(4,5)$, but (H5) fails at the origin endpoints that the theorem states: its forward
half holds there, but the stationary dual has no legal excursion from either phase at the origin, and
$K_n(o,((0,0),b))=0$ for all $n\ge1$ and both phases $b$. Report~123's exponent is therefore an
instance of the theorem's \emph{proof}, run at the endpoint $((1,0),P)$, not of its statement.
Proposition~26.2 there (\texttt{icw:log:prop:endpoints}) writes this out: under (H1)--(H4) the proof of
Theorem~\ref{cps:log:thm:cone} gives survival and the kernel upper bound uniformly over starts and
endpoints at distance $O(\log n)$ from the origin, and, with seeds from arbitrary fixed endpoint
states, the matching lower bound; its Corollary~26.3 (\texttt{icw:log:cor:exponent}) derives
Report~123's exponent from it. The independent check of that report (commit \texttt{1c7e526c7}) read
the proof of Theorem~\ref{cps:log:thm:cone} in full and confirmed that Proposition~26.2 follows from it
unchanged: the start enters only through the first radius $r_0(n)$, and the endpoints only through the
seeds, their largest radii, the reversal ratio of stationary masses and first-hit uniqueness. The same
walk is an example for Question~\ref{cps:log:q:seeds}, and that report's Part~I bears on
Question~\ref{cps:log:q:amplitude}; see the notes there.)

\section{Results and their precise scope}
Write $p_n=\mathrm{A377922}(n)$ and $s_n=\mathrm{A377920}(n)$, both with
index starting at zero. The former counts polyhedral orientations with $n$
inner vertices; the latter counts all 3-connected Schnyder labelings with
$n$ inner faces in the source's $(6,4)$-dissection model. The established
bijections and conventions are those of \cite{FNS}, Propositions 13--14,
19 and 23. These definitions matter: A377923 counts corner-polyhedron
\emph{graphs}, which are different objects and are excluded here.

For a positive sequence $a_n$, the notation
\[
 a_n=\Gamma^n n^{-\alpha+o(1)}
\]
means
\[
 \lim_{n\to\infty}\frac{\log(a_n/\Gamma^n)}{\log n}=-\alpha.
\]
It does not mean that $a_n/(\Gamma^n n^{-\alpha})$ converges.

\begin{theorem}[The three counting sequences]\label{cps:log:thm:main}
Let $b_n=\mathrm{A377921}(n)$ and
\[
 \gp=\frac92,\quad \ap=1+\frac\pi{\arccos(9/16)},\qquad
 \gs=\frac{16}3,\quad \as=1+\frac\pi{\arccos(22/27)}.
\]
Then
\begin{align}
 p_n&=\gp^n n^{-\ap+o(1)},\label{cps:log:eq:mainp}\\
 s_n&=\gs^n n^{-\as+o(1)},\label{cps:log:eq:mains}\\
 b_n&=\gs^n n^{-\as+o(1)}.\label{cps:log:eq:mainb}
\end{align}
All three sequences are eventually positive at every integer index. Their
ordinary generating functions are not D-finite over $\mathbb C(z)$.
For each triple $(a_n,\Gamma,\alpha)$ equal to
$(p_n,\gp,\ap)$, $(s_n,\gs,\as)$ or $(b_n,\gs,\as)$, and the first threshold
$N_a(y)=\min\{n\ge0:a_n\ge y\}$,
\begin{equation}\label{cps:log:eq:inverse-main}
 N_a(y)=\frac{\log y}{\log\Gamma}
       +\frac{\alpha}{\log\Gamma}\log\log y
       +o(\log\log y)\qquad(y\to\infty).
\end{equation}
No monotonicity assumption is needed.
The $b_n$ conclusions use the six-face injection in Section~\ref{cps:log:sec:921}.
\end{theorem}

\begin{writenotelog}
Every conclusion of this theorem is proved in a stronger form in
Parts~\ref{cps:part:exponents}--\ref{cps:part:amplitudes}: \eqref{cps:log:eq:mainp} and
\eqref{cps:log:eq:mains} are \eqref{cps:eq:loglimits}, consequences of the $\Theta$ bounds of
Theorem~\ref{cps:thm:main} and of the equivalents of Theorem~\ref{cps:amp:thm:main};
\eqref{cps:log:eq:mainb} follows from the equivalent \eqref{cps:amp:eq:rigidmain}; non-D-finiteness
of the three generating functions is Corollary~\ref{cps:cor:nonD}; and Lemma~\ref{cps:lem:inverse}
gives \eqref{cps:log:eq:inverse-main} with error $O(1)$, Theorem~\ref{cps:amp:thm:inverse} with the
next constant. Theorem~\ref{cps:log:thm:main} is therefore a second route, by a different method
(Table~\ref{cps:log:tab:map}).
\end{writenotelog}

The proof separates into a general probabilistic result
(Section~\ref{cps:log:sec:cone}), exact model checks
(Sections~\ref{cps:log:sec:corner}--\ref{cps:log:sec:schnyder}), and arithmetic and inverse
arguments (Sections~\ref{cps:log:sec:nonD}--\ref{cps:log:sec:inverse}). Section~\ref{cps:log:sec:921}
proves the explicit six-face injection and resulting all-index theorem
for A377921.

\paragraph{What was already known.}
Theorem 24 of \cite{FNS} proves the exponential growth constants $9/2$ and
$16/3$. Its Conjecture 25 proposes stronger full coefficient equivalents
with the powers appearing above. Equations~\eqref{cps:log:eq:mainp}--\eqref{cps:log:eq:mains}
establish the exponents on the logarithmic scale only. They leave the
positive multiplicative constants and the full equivalents unresolved.
The literature review in Section~\ref{cps:log:sec:literature} is scoped; it is not
a certificate of priority or an exhaustive novelty claim.

\paragraph{Numerical orientation.}
The exponents are approximately $4.23$ and $6.08$, respectively. All proofs
and all exact checks use the symbolic expressions, not these decimals.
The stronger inverse precision $O(1)$, further inverse terms and
coefficient transseries are not consequences of this report.

\begin{writenotelog}
Both paragraphs describe Report 112 alone. The positive constants and the full equivalents left
unresolved here are proved in Part~\ref{cps:part:amplitudes} (Theorem~\ref{cps:amp:thm:main}); the
inverse with error $O(1)$ is Lemma~\ref{cps:lem:inverse}, and its next additive constant is in
Theorem~\ref{cps:amp:thm:inverse}. No coefficient transseries is proved anywhere in this report.
\end{writenotelog}

\section{A cone theorem with finite phases and exponential tails}\label{cps:log:sec:cone}
\subsection{Hypotheses and statement}
Consider a translation-invariant Markov-additive chain $(S_n,C_n)$ on
$\ZZ^2\times E$, with $E$ finite and
\[
 p_{cd}(v)=\PP(S_{n+1}-S_n=v,C_{n+1}=d\mid C_n=c),\qquad
 \sum_{d,v}p_{cd}(v)=1.
\]
Let $P_{cd}=\sum_v p_{cd}(v)$. The spatial coordinates need not have
independent or bounded increments. Assume the following.
\begin{enumerate}[label=\textup{(H\arabic*)}]
\item $P$ is primitive. Its stationary distribution $\pi$ has all entries
strictly positive.
\item For some $\lambda>0$, $\sup_c\EE_c e^{\lambda|S_1-S_0|}<\infty$.
Any fixed norm may be used.
\item The stationary drift is zero and the effective covariance
$\Sigma$ is positive definite.
\item For the Fourier matrix
$F(t)_{cd}=\sum_v p_{cd}(v)e^{it\cdot v}$,
$\rad F(t)<1$ at every nonzero point of
$\RR^2/(2\pi\ZZ)^2$. At zero the only unit eigenvalue is the simple
Perron eigenvalue $1$.
\item Fix endpoint phases $a,b$ at spatial origin. Put
$W=\Sigma^{-1/2}$ and let $\mathcal K=W\RR_+^2$ be the whitened wedge,
with angle $\theta\in(0,\pi)$. There is a compact angular interval of positive length
$J\subset(0,\theta)$ such that, for every $R>0$, a finite
positive-probability legal path from $(0,a)$ reaches a point $x$ with
$|Wx|\ge R$ and $\argu(Wx)\in J$. The same holds from $(0,b)$ for the
stationary dual, defined by
\begin{equation}\label{cps:log:eq:dual}
 p^*_{dc}(-v)=\frac{\pi_c}{\pi_d}p_{cd}(v).
\end{equation}
\end{enumerate}
Legal means that every visited lattice point lies in the \emph{closed}
quadrant. Seed paths in (H5) have no prescribed length, no probability
uniform in $R$, and no requirement to stay inside their terminal radius.
For each tolerance in the proof a single sufficiently deep seed is fixed
before $n$ tends to infinity.

Let $K_n((x,c),(y,d))$ denote the killed transition kernel, and let
$\tau=\inf\{n\ge0:S_n\notin\RR_+^2\}$. Write
\[
 \rho=\frac{\Sigma_{12}}{\sqrt{\Sigma_{11}\Sigma_{22}}},\qquad
 \theta=\arccos(-\rho),\qquad \nu=\frac\pi\theta.
\]
The angle formula follows by computing the angle between $We_1$ and
$We_2$ from $W^TW=\Sigma^{-1}$.

\begin{theorem}[Finite-phase logarithmic cone excursions]\label{cps:log:thm:cone}
Under \textup{(H1)--(H5)}, at the specified phases,
\begin{equation}\label{cps:log:eq:cone}
 K_n((0,a),(0,b))=n^{-1-\nu+o(1)}.
\end{equation}
Under \textup{(H1)--(H4)} alone, for each fixed starting phase and every
$\epsilon>0$,
\begin{equation}\label{cps:log:eq:survival}
 \PP_{(0,c)}(\tau>n)\le C_\epsilon n^{-\nu/2+\epsilon}.
\end{equation}
The same upper bound holds for the dual chain.
\end{theorem}

The seed condition is substantive. An unrestricted centered chain with
nondegenerate covariance can have a boundary phase from which no legal
interior excursion is possible. The theorem makes no claim for such an
endpoint. It also does not assert a constant equivalent in
\eqref{cps:log:eq:cone}.

\begin{writenotelog}
Part~\ref{cps:part:exponents} (Section 1) says that it does not posit a general cone theorem for
Markov-modulated walks and that its argument uses a geometric property of the two models, the
one-unit cycle buffer of Lemma~\ref{cps:lem:buffer}. Theorem~\ref{cps:log:thm:cone} is a general
theorem of that kind at the logarithmic scale: it needs no regeneration, no cycle buffer and no
iid cone input, and it allows unbounded steps. Its price is the factor $n^{o(1)}$. For the two
models of this report Parts~\ref{cps:part:exponents}--\ref{cps:part:amplitudes} give $\Theta$
bounds and equivalents; whether (H1)--(H5) alone give either is open
(Question~\ref{cps:log:q:amplitude}). Section~\ref{cps:log:sec:siblings} names two other chains
of the same kind.
\end{writenotelog}

\subsection{Corrector and unrestricted limits}\label{cps:log:sec:limits}
Let $m(c)=\EE_c(S_1-S_0)$. Stationary centering solves the finite Poisson
equation
\[
 (I-P)h=m,\qquad \sum_c\pi_ch(c)=0.
\]
The process
\[
 Z_n=S_n+h(C_n)-h(C_0)
\]
is a martingale. Its increments have a uniform exponential moment. If
$V(c)$ is their conditional covariance, then
$\Sigma=\sum_c\pi_cV(c)$. The finite-state ergodic theorem gives the
rescaled predictable quadratic variation $t\Sigma$. Its jumps are
uniformly bounded before the $n^{-1}$ scaling, while
\[
 \EE\max_{j\le n}|\Delta Z_j|^2=O(\log^2(n+2)).
\]
The last estimate follows by integrating the union bound for exponential
tails. Thus the maximum rescaled squared jump tends to zero and the
martingale functional central limit theorem applies; see
\cite[Theorem 2.1(ii)]{Whitt}. Since $h$ is bounded,
\begin{equation}\label{cps:log:eq:fclt}
 R^{-1}(S_{\lfloor R^2t\rfloor}-S_0)\ \Longrightarrow\ B_\Sigma(t).
\end{equation}
The convergence is uniform over the finite phase set. Translation
invariance allows varying initial positions whose rescaled positions
converge. Polygonal and step interpolation differ by $o_{\PP}(1)$,
since $R^2\PP(|\Delta S|>\epsilon R)\to0$.

The exponential moment makes $F$ analytic near the real Fourier torus.
Its simple eigenvalue near zero and spectral projection satisfy
\[
 \lambda(t)=1-\tfrac12t^T\Sigma t+O(|t|^3),\qquad
 \Pi(t)\longrightarrow\one\pi.
\]
For example, the Hessian of the log Perron eigenvalue agrees with the
martingale covariance because the bounded corrector does not change
asymptotic variance. Away from zero, (H4), compactness and finite-matrix
spectral estimates bound powers exponentially. Near zero, the positive
definite quadratic term gives a Gaussian bound. Fourier inversion,
rescaling $t=w/\sqrt n$, and dominated convergence therefore give
\begin{equation}\label{cps:log:eq:llt}
 \sup_{z,c,d}\left|n\PP_c(S_n-S_0=z,C_n=d)
           -\pi_d g_\Sigma(z/\sqrt n)\right|\longrightarrow0,
\end{equation}
where $g_\Sigma$ is the centered normal density with covariance $\Sigma$.
Uniformity in $z$ follows because the Fourier factors have modulus one.
In particular the unrestricted kernel is at most $C/n$ everywhere and
at least $c/n$ on compact sets of scaled displacements.

The dual preserves centering, exponential tails and effective covariance.
Its Fourier matrix is $\diag(\pi)^{-1}F(-t)^T\diag(\pi)$, so it preserves
(H4). Reversing an individual path gives
\begin{equation}\label{cps:log:eq:reversal}
 \pi_cK_j((x,c),(y,d))=\pi_dK_j^*((y,d),(x,c)).
\end{equation}

\subsection{Exact-time bridges in the interior}\label{cps:log:sec:bridge}
\begin{lemma}\label{cps:log:lem:bridge}
For every compact $A\subset(0,\infty)^2$ and $0<u<v<\infty$, there
exists $c_A>0$ such that, for all sufficiently large $n$,
\begin{equation}\label{cps:log:eq:bridge}
 K_m((x,c),(y,d))\ge c_A/n
\end{equation}
uniformly for $un\le m\le vn$, $x/\sqrt n,y/\sqrt n\in A$ and all phases.
\end{lemma}
\begin{proof}
Condition the unrestricted chain on its exact endpoint. Through time
$k=\lfloor tm\rfloor$, for fixed $t<1$, its density relative to the
unconditioned chain is the remaining endpoint kernel divided by the full
endpoint kernel. Equation~\eqref{cps:log:eq:llt} bounds this density uniformly
and makes it converge on compact scaled positions to the Gaussian
ratio defining a Brownian bridge. The FCLT, tightness and compact
truncation give the bridge limit on $[0,t]$. Reversal gives the analogous
terminal portion through a dual bridge. Tightness on $[0,2/3]$ and
$[1/3,1]$ yields tightness on the whole interval, as every pair of
sufficiently close times lies in one of these intervals.

The segment between two points of $A$ stays a uniformly positive
distance from the sides of the quadrant. A Brownian bridge has a
uniformly positive probability of staying in a sufficiently thin open
tube around this segment, uniformly over the compact endpoint and time
parameters. Portmanteau, followed by the unrestricted endpoint lower
bound $c/n$, proves~\eqref{cps:log:eq:bridge}. Convexity makes polygonal
containment equivalent to containment of all visited vertices, even
when a jump is long. The bridge tightness argument already controls
rare large jumps under conditioning; no extra conditional tail
assumption is used.
\end{proof}

\subsection{Brownian crossing bounds and relative overshoot}\label{cps:log:sec:crossings}
Use whitened coordinates $X_n=WS_n$. Both chains have a common bound
\begin{equation}\label{cps:log:eq:tail}
 \PP_c(|\Delta X|>r)\le A e^{-\kappa r}.
\end{equation}
The FCLT also supplies constants $A_0,b_0,R_1>0$ such that, from every
point in the radius-$R$ disk and every phase,
\begin{equation}\label{cps:log:eq:confine}
 \PP_x\left(\max_{j\le t}|X_j|\le R\right)
 \le A_0 e^{-b_0t/R^2}\qquad(R\ge R_1).
\end{equation}
Indeed an unrestricted displacement larger than $3R$ in
$\lceil R^2\rceil$ steps forces an exit, with probability bounded below
by the FCLT uniformly in phase. Iterate the Markov property.

For planar Brownian motion in the wedge, the probability of reaching
radius $q$ before a side, starting at radius one, is at most
$C_Bq^{-\nu}$ uniformly in the starting angle. If the starting angle
lies in a fixed compact interval $J$ inside the wedge, the probability
of exiting through the interior of the outer arc with angle in $J$ is
at least $c_Bq^{-\nu}$ for all sufficiently large $q$. The constants
are independent of $q$. One direct verification is the sine expansion
of the Dirichlet solution with outer boundary data $f$:
\begin{equation}\label{cps:log:eq:brownian}
 H_f(r,\phi)=\sum_{k\ge1}b_k(r/q)^{k\nu}\sin(k\nu\phi),\qquad
 b_k=\frac2\theta\int_0^\theta f(\psi)\sin(k\nu\psi)\,d\psi.
\end{equation}
For $f=1$, the sum is $O(q^{-\nu})$ at $r=1$. For
$f=\one_J$, the first coefficient is positive and
$\min_{\phi\in J}\sin(\nu\phi)>0$, whereas the remaining sum is
$O(q^{-2\nu})$. Brownian disk confinement truncates the successful
exit time at $T(q)q^2$ while retaining a fixed fraction of its mass.

Let $\sigma_r=\inf\{j:|X_j|\ge r\}$. There are constants $c,C>0$
independent of sufficiently large fixed $q$, then choices $T=T(q)$ and
$R_0=R_0(q)$, such that every start with
$R\le|X_0|\le R+\sqrt R$, $R\ge R_0$, satisfies
\begin{equation}\label{cps:log:eq:cross-upper}
 \PP_x(\sigma_{qR}<\tau)\le Cq^{-\nu}.
\end{equation}
If additionally its angle lies in $J$, then
\begin{equation}\label{cps:log:eq:cross-lower}
 \PP_x\left(\begin{array}{l}
 \sigma_{qR}<\tau,\quad \sigma_{qR}\le T(qR)^2,\\
 \argu X_{\sigma_{qR}}\in J,\quad
 \max_{1\le j\le\sigma_{qR}}|\Delta X_j|\le\sqrt R
 \end{array}\right)\ge cq^{-\nu}.
\end{equation}
Here the maximum concerns increments \emph{after this stage begins}.

For completeness, the uniform transfer from Brownian motion is as
follows. The permitted starting radii divided by $R$ tend uniformly
to one. On each fixed scaled horizon the event of reaching the outer
circle while staying in the closed wedge until that time is closed.
At a limiting side start its Brownian probability is zero, since
Brownian motion immediately enters the exterior half-plane. At an
interior start the harmonic bound applies. Compact subsequences of
starting points and the finite phase set make the upper estimate
uniform. Hits after the chosen horizon are bounded by
\eqref{cps:log:eq:confine}. Choose $T(q)$ first and then $R_0(q)$ so each
additional error is at most a fixed multiple of $q^{-\nu}$.

For the lower estimate, the exit event through the interior of $J$
is an almost-sure continuity event: Brownian motion hits no corner or
arc endpoint and immediately crosses a smooth exit circle. On a
successful path it stays a positive distance from the sides up to and
just after exit. A strict time cutoff with no atom may be chosen.
The FCLT gives the lower estimate before imposing the jump restriction.
The extra excluded probability is at most
\begin{equation}\label{cps:log:eq:overshoot-error}
 \lceil Tq^2R^2\rceil A e^{-\kappa\sqrt R}=o_R(1).
\end{equation}
For the already fixed $q$, it is smaller than the allotted fraction
of $c_Bq^{-\nu}$ for all large $R$. The terminal radius then lies in
$[qR,qR+\sqrt R]\subset[qR,qR+\sqrt{qR}]$, precisely the allowance
required for the next stage. The probabilities above are always under
the original chain; its transition law is never conditioned or truncated.

\subsection{Upper survival from good past events}\label{cps:log:sec:upper}
Fix $\eta\in(0,\nu)$ and choose one sufficiently large $q\ge2$ so that
\begin{equation}\label{cps:log:eq:chooseq}
 Cq^{-\nu}\le q^{-\nu+\eta},\qquad
 cq^{-\nu}\ge q^{-\nu-\eta}.
\end{equation}
After $T(q),R_0(q)$ are fixed, enlarge $R_0$ so
$\sqrt R<(q-1)R$ for $R\ge R_0$.
Take $M>\nu+2$ and $D$ sufficiently large. Put
\[
 L_n=D\log(n+2),\qquad
 \mathcal G_n=\{|\Delta X_k|\le L_n\text{ for }1\le k\le n\}.
\]
By~\eqref{cps:log:eq:tail},
\begin{equation}\label{cps:log:eq:badjump}
 \PP(\mathcal G_n^c)\le nAe^{-\kappa L_n}=O(n^{-M}).
\end{equation}
For this horizon define
$r_0(n)=\max\{R_0,L_n^2\}$ and $r_j=q^jr_0(n)$, and let $\sigma_j$
be the first global hit of radius $r_j$. The relevant measurable
\emph{past} event is
\[
 A_j=\{\sigma_j\le n,\ \sigma_j<\tau,
             \ |\Delta X_k|\le L_n\text{ for all }k\le\sigma_j\}.
\]
On $A_j$ the overshoot is at most $L_n\le\sqrt{r_j}$.
A good jump cannot skip the next radius, so $A_{j+1}\subset A_j$.
At $\sigma_j$, drop all future jump restrictions and the time deadline.
The probability of the required next crossing before killing is at
most $q^{-\nu+\eta}$ by~\eqref{cps:log:eq:cross-upper}. Strong Markov gives
\begin{equation}\label{cps:log:eq:annularupper}
 \PP(A_{j+1})\le q^{-\nu+\eta}\PP(A_j),\qquad
 \PP(A_J)\le q^{-(\nu-\eta)J}.
\end{equation}
This step does not condition the process on the future event
$\mathcal G_n$. All future crossing estimates use its original law.

Take $R_n=\sqrt{n/(B\log(n+2))}$ with fixed $B$ so large that the
confinement bound is $O(n^{-M})$, and choose $J$ with
$r_J\le R_n<r_{J+1}$. A surviving path either has a bad jump, remains
inside radius $R_n$, or belongs to $A_J$. Therefore
\begin{align}
 \PP(\tau>n)
 &\le O(n^{-M})+C(R_n/r_0(n))^{-\nu+\eta}
                  +A_0e^{-b_0n/R_n^2}\notag\\
 &\le n^{-(\nu-\eta)/2}(\log n)^{O(1)}+O(n^{-M}).\label{cps:log:eq:survival-proof}
\end{align}
Choosing $\eta$ sufficiently small proves~\eqref{cps:log:eq:survival} for
any prescribed $\epsilon>0$. The same argument applies to the dual.
No eventual escape estimate from a fixed small disk is needed.

\subsection{Lower escapes and exact-time concatenation}\label{cps:log:sec:lower}
After $\eta,q,T,R_0$ have been fixed, choose one legal seed for each
chain whose terminal radius $r_0$ is at least $R_0$ and whose terminal
angle lies in $J$. The two seed radii may differ. Each seed length,
probability and largest visited radius is a fixed finite constant.
For each chain separately use radii $r_j=q^jr_0$ and successively
accept the events in~\eqref{cps:log:eq:cross-lower}. Their total mass through
radius $r_J$ is at least the seed probability times $q^{-(\nu+\eta)J}$.
The total elapsed time is at most
\begin{equation}\label{cps:log:eq:lower-time}
 t_{\rm seed}+T\sum_{j=1}^Jr_j^2
 \le t_{\rm seed}+\frac{Tr_J^2}{1-q^{-2}}.
\end{equation}
Choose fixed $\delta>0$ with
$T\delta^2/(1-q^{-2})\le1/8$, and for each seed use the largest $J$
with $r_J\le\delta\sqrt n$. Then $r_J>\delta\sqrt n/q$, the total
time is at most $n/4$ for large $n$, and endpoints divided by $\sqrt n$
lie in one fixed compact subset of the open quadrant. The forward and
dual masses $F_n,G_n$ obey
\begin{equation}\label{cps:log:eq:escapemass}
 F_n,G_n\ge c_\eta n^{-(\nu+\eta)/2}.
\end{equation}

For large $n$, the chosen final radius exceeds the largest radius
visited anywhere in the corresponding fixed seed. Hence an accepted
prefix ends at the \emph{first global hit} of its final radius, even
if the seed had previously overshot its own terminal radius. This is
why no monotone-radius condition was imposed on the seeds.

Join any accepted forward prefix of length $t\le n/4$ to the reversal
of any accepted dual prefix of length $u\le n/4$ with a killed bridge
of exact length $n-t-u$. Lemma~\ref{cps:log:lem:bridge} bounds its mass below
by $c/n$. Reversal~\eqref{cps:log:eq:reversal} changes the suffix weight by
$\pi_b/\pi_d$, bounded above and below. There is no overcounting:
the first forward hit of the prescribed forward radius and first
backward hit of the prescribed backward radius recover the two
prefixes uniquely. Their lengths cannot overlap. Consequently
\begin{equation}\label{cps:log:eq:lower-final}
 K_n((0,a),(0,b))\ge(c/n)F_nG_n
       \ge c_\eta n^{-1-\nu-\eta}.
\end{equation}
For the upper bound, split $n$ into three comparable integer pieces.
Bound the middle unrestricted kernel by $C/n$ using~\eqref{cps:log:eq:llt},
and bound the outer pieces by forward and dual survival using
\eqref{cps:log:eq:reversal}. Equation~\eqref{cps:log:eq:survival} yields
\begin{equation}\label{cps:log:eq:upper-final}
 K_n((0,a),(0,b))\le C_\epsilon n^{-1-\nu+\epsilon}.
\end{equation}
Together these bounds prove Theorem~\ref{cps:log:thm:cone}. They retain every
integer time and every intervening jump. In particular, the result
is not a random-time or subsequence estimate.

\section{Polyhedral orientations and A377922}\label{cps:log:sec:corner}
\subsection{Original walk and quotient coordinates}
For $n\ge1$, the exact walk correspondence in \cite[Proposition 13]{FNS} counts $p_n$
by length-$n$ walks from $(0,0)$ to the horizontal axis, remaining in
$\ZZ_{\ge0}^2$. An SE step is $(1,-1)$. A face step is $(-i,j)$ with
$i,j\ge0$ and $i+j$ even; when the current height is even it must have
$j\ge1$, and when the current height is odd it must have $i\ge1$.
These conditions keep $x\equiv y\pmod2$.

Encode this physical parity by
\begin{equation}\label{cps:log:eq:cornerquot}
 x=2U+c,\qquad y=2V+c,\qquad c\in\{0,1\}.
\end{equation}
This is a genuine two-phase chain on a full quotient lattice. The
physical inequalities $x,y\ge0$ are exactly $U,V\ge0$ for either
phase, including the boundary.

Set $z=1/3$, $\gp=9/2$ and $D=(1-z/u)(1-zv)$. The normalized transition
generating matrix is
\begin{equation}\label{cps:log:eq:Pkernel}
 P(u,v)=\gp^{-1}
 \begin{pmatrix}
 zv/D&z^{-1}v^{-1}+z/(uD)\\
 u/z+zv/D&z/(uD)
 \end{pmatrix}.
\end{equation}
A coefficient $u^rv^s$ records quotient displacement $(r,s)$.
The tilted physical weight of an elementary step is
$z^{(\Delta y-\Delta x)/2}/\gp$.
The factor $1/2$ in this exponent is required by the two-coordinate
quotient. In particular it must be retained in telescoping path weights.

An elementary probabilistic reconstruction makes positivity and every
phase explicit. Choose SE with probability $2/3$ or a face with
probability $1/3$. For a face, choose destination phase $d$ uniformly,
and independently choose $I,J\ge0$ with
$\PP(I=i)=\PP(J=i)=(2/3)(1/3)^i$. Then
\begin{align*}
 \text{SE: }&d=1-c,\quad\Delta=(c,c-1),\\
 \text{face: }&\Delta=(-I-d,J+1-d).
\end{align*}
This reproduces~\eqref{cps:log:eq:Pkernel} directly. Thus
\[
 P(1,1)=\begin{pmatrix}1/6&5/6\\5/6&1/6\end{pmatrix},\qquad
 \pi=(1/2,1/2).
\]

\subsection{Centering and covariance}\label{cps:log:sec:Pcov}
The conditional drifts and mean-zero corrector are
\[
 m_0=(-1/3,-1/3)=-m_1,\qquad
 h_0=(-1/5,-1/5)=-h_1.
\]
They satisfy $(I-P(1,1))h=m$. The covariance of the corrected
increment $T=\Delta+h_d-h_c$ under stationarity is
\begin{equation}\label{cps:log:eq:Psigma}
 \Sigma_P=\begin{pmatrix}4/5&-9/20\\-9/20&4/5\end{pmatrix},
 \qquad \rho_P=-9/16.
\end{equation}
Its eigenvalues are $7/20$ and $5/4$.

Here is a direct second-moment calculation independent of Perron
root differentiation. On a face,
$T=(-I-B,J+1-B)$ with $B=(3d+2c)/5$. Conditional on a face in
stationarity, $c,d$ are independent uniform bits,
$\EE B=1/2$ and $\Var B=13/100$. Hence the face mean is $(-1,1)$,
each diagonal covariance is $22/25$, and the off-diagonal covariance
is $13/100$. On SE,
$T=(2/5+c/5,-3/5+c/5)$, its mean is $(1/2,-1/2)$, and every covariance
entry is $1/100$. Mixing the \emph{second moments} with probabilities
$1/3$ and $2/3$ gives~\eqref{cps:log:eq:Psigma}; the total mean is zero.

The executable certificate also computes the Hessian at $(1,1)$ of
$\log\lambda_P(u,v)$ with logarithmic derivatives
$u\partial_u,v\partial_v$, where $\lambda_P$ is the Perron root.
This is a separate exact verification of the same effective covariance.

The physical asymptotic covariance is $4\Sigma_P$, because the phase
coordinate offsets are bounded. The matrix displayed in the source's
asymptotic calculation is $\gp$ times this physical stochastic
covariance. Differentiating an unnormalized Perron root rather than
its normalized logarithm produces that scale. Correlation, and thus
the cone angle, are unchanged; the scales must not be interchanged in
a stochastic normalization.

\begin{writenotelog}
Part~\ref{cps:part:exponents} makes the same observation after \eqref{cps:eq:Sigma}, in physical
units: its $\Sigma_P$ is the $4\Sigma_P$ of this paragraph, and \cite[Section 4.2]{FNS} displays
$\mu_P\Sigma_P$. The kernel \eqref{cps:log:eq:Pkernel} is $K_P$ of \eqref{cps:eq:PK} in
half-scale variables (both give the face with $c=d=0$, $I=i$, $J=j$ the weight
$(2/27)3^{-i-j}$); the drifts, correctors and covariances correspond as in
Remark~\ref{cps:log:rem:scale}.
\end{writenotelog}

\subsection{Exponential tails and exact aperiodicity}
The quotient increments are affine functions of geometric variables
and bounded phases. Thus $\sup_c\EE_c e^{t\|\Delta\|_1}<\infty$ for
$0<t<\log3$. The bounded corrector and stationary dual preserve a
positive exponential-moment neighborhood.

To check the full lattice and time condition (H4), suppose the Fourier
matrix has a unit-modulus eigenvalue $\lambda$ at frequency $t$.
Irreducibility and equality in the triangle inequality give unit
phases $a_c$ with
\begin{equation}\label{cps:log:eq:phase-equality}
 e^{it\cdot\delta}a_d=\lambda a_c
 \quad\text{on every supported edge }(c,\delta,d).
\end{equation}
The supported phase-$0$ loops $(0,1),(0,2)$ force
$e^{it_2}=e^{2it_2}=\lambda$, hence $t_2=0$ modulo $2\pi$ and
$\lambda=1$. The phase-$1$ loop $(-1,0)$ forces $t_1=0$.
This proves (H4) on the full quotient torus. An assertion of full
lattice aperiodicity in the redundant physical coordinates would be
incorrect; their parity is precisely what the phase representation
removes.

\subsection{Legal seeds and endpoint normalization}
For any integer $H\ge1$, a forward seed from quotient $(0,0,0)$ is:
take a phase-$0$ loop $(0,2H)$, then repeat
$0\to1:(0,-1)$ followed by $1\to0:(1,0)$, $H$ times. It reaches
$(H,H,0)$ legally.

The dual from $(0,0,1)$ takes $2H$ phase-$1$ east loops $(1,0)$,
then repeats $1\to0:(0,1)$ followed by $0\to1:(-1,0)$, $H$ times.
It reaches $(H,H,1)$ legally. Every listed edge has positive
probability. After whitening, the ray through $(H,H)$ has a fixed
strictly interior angle. This proves (H5), with an interval of positive
length surrounding that angle. The seeds may visit a greater radius
before their terminal point, which is permitted.

Let $e_L$ count legal original paths of length $L$ from $(0,0)$ to
$(1,1)$. The endpoint tilt equals one, and therefore
\begin{equation}\label{cps:log:eq:Pendpoint}
 e_L=\gp^L K_L((0,0,0),(0,0,1)).
\end{equation}
The ordered triples here mean $(U,V,c)$. For $n\ge1$,
\begin{equation}\label{cps:log:eq:Psandwich}
 e_{n-1}\le p_n\le e_{n+1}.
\end{equation}
For the lower injection append SE to a path ending at $(1,1)$;
its endpoint becomes $(2,0)$. For the upper injection, a nonempty
legal path ending on the horizontal axis cannot end at $(0,0)$:
an incoming SE would require $x=-1$, and an incoming face at even
height zero is forbidden. Its endpoint is $(2i+2,0)$ for some $i\ge0$.
Append the allowed face $(-2i-1,1)$ to reach $(1,1)$. Removing that
last step recovers the original path, proving injectivity.

Theorem~\ref{cps:log:thm:cone},~\eqref{cps:log:eq:Pendpoint} and the two fixed shifts
in~\eqref{cps:log:eq:Psandwich} prove~\eqref{cps:log:eq:mainp}. This argument avoids
summing boundary probabilities with unequal endpoint tilt weights.

\begin{writenotelog}
The sandwich \eqref{cps:log:eq:Psandwich} is new to this report.
Part~\ref{cps:part:exponents} uses the tilted endpoint sum \eqref{cps:eq:Pcount}, which needs the
uniform endpoint bound \eqref{cps:eq:Pupper}; Part~\ref{cps:part:amplitudes} sums the same series
with amplitudes \eqref{cps:amp:eq:kappaP}. Combined with Theorem~\ref{cps:amp:thm:bridge} at the
endpoints $(0,0)$ and $(1,1)$, the sandwich gives only $\Theta$ for $p_n$, because the two fixed
shifts change $\mu_P^n$ by the factors $\mu_P^{\pm1}$.
\end{writenotelog}

\section{Schnyder labelings and A377920}\label{cps:log:sec:schnyder}
\subsection{The original weighted walk and its time parameter}
The source's auxiliary subfamily $s'_n$ imposes nonisolation at the
two outer vertices labeled $G$. Its walk correspondence
\cite[Proposition 14]{FNS}, for genuine sizes $n\ge1$, uses paths
from physical $(0,2)$ to $(2,0)$ in the quadrant with exactly $n+2$
SE steps. Face multiplicities are important. With $\ell,r\ge0$,
a face $(-i,j)$ has weight $\binom{\ell+r}{r}$ for
\[
 (i,j)=(2\ell+2,2r+2),
\]
or for
\[
 (i,j)=
 \begin{cases}
 (2\ell+1,2r+3),&\text{current height even},\\
 (2\ell+3,2r+1),&\text{current height odd}.
 \end{cases}
\]
Every face strictly lowers $x$ and strictly raises $y$.

One aggregate is one SE step followed by all succeeding faces up to
the next SE. Its duration as a sequence of elementary steps is
variable, but it contains exactly one SE. This aggregate count is
the required deterministic size parameter. Treating a face
micro-decomposition or its random duration as the counting time would
change the enumerative problem.

For physical generating variables $a,b$ put
\begin{equation}\label{cps:log:eq:FJ}
 F(a,b)=\frac1{1-a^{-2}-b^2}
 \begin{pmatrix}a^{-2}b^2&a^{-1}b^3\\a^{-3}b&a^{-2}b^2\end{pmatrix},
 \qquad J(a,b)=\begin{pmatrix}0&a/b\\a/b&0\end{pmatrix}.
\end{equation}
Then the aggregate matrix is
\begin{equation}\label{cps:log:eq:aggregate}
 A(a,b)=J(a,b)(I-F(a,b))^{-1}.
\end{equation}
The matrix order is essential: SE changes parity before the following
faces. The nonnegative expansion $J\sum_{k\ge0}F^k$ retains all
face multiplicities. At $a=2,b=1/2$, every entry of $F$ is $1/8$,
its row sums are $1/4$, and the row sums of $A$ are $\gs=16/3$.

\subsection{Shifted quadrant and exact killing equivalence}
An original path must start with SE, since no face is legal at $x=0$,
and end with SE, since faces strictly raise $y$. Deleting its last
SE leaves $n+1$ aggregates from $(0,2)$ to $(1,1)$.
Aggregate endpoints lie in the \emph{shifted} physical quadrant
$x\ge0,y\ge1$. Introduce
\begin{equation}\label{cps:log:eq:Squot}
 x=2U+c,\qquad y=2V+2-c,\qquad c\in\{0,1\}.
\end{equation}
For both phases the shifted condition is exactly $U,V\ge0$.
The two endpoints become quotient $(0,0,0)$ and $(0,0,1)$.

Within an aggregate, SE raises $x$ and lowers $y$ by one, after which
faces only lower $x$ and raise $y$. Its minimum $x$ is at an endpoint;
its minimum $y$ is immediately after SE. Shifted-quadrant endpoints
therefore guarantee that every elementary point stays in the original
quadrant. Conversely, every nonfinal aggregate endpoint precedes a
legal SE and hence has $y\ge1$; the prescribed final endpoint has
$y=1$. This proves the equivalence in both directions, including every
intermediate face point. There is no hidden within-block boundary
crossing after aggregation.

\subsection{The rational two-phase kernel}
The normalized quotient kernel is
\begin{equation}\label{cps:log:eq:Gtransform}
 G_{cd}(u,v)=\gs^{-1}A_{cd}(2\sqrt u,\sqrt v/2)
                u^{(c-d)/2}v^{(d-c)/2}.
\end{equation}
The phase factors remove the apparent half-integral exponents. A
simpler exact rational form is obtained by setting
\[
 C(u,v)=\frac3{8(8u-2uv-v-2)}
       =\frac{3/(64u)}{1-v/4-1/(4u)-v/(8u)}.
\]
Then
\begin{equation}\label{cps:log:eq:Gkernel}
 G(u,v)=\begin{pmatrix}
 C&3/4+vC\\
 3u/(4v)+uC&uvC
 \end{pmatrix}.
\end{equation}
The displayed geometric expansion proves nonnegative Laurent
coefficients directly. In particular
\[
 G(1,1)=\begin{pmatrix}1/8&7/8\\7/8&1/8\end{pmatrix},\qquad
 \pi=(1/2,1/2).
\]
The conditional means and corrector are
\[
 m_0=(-1/2,3/8)=-m_1,\qquad
 h_0=(-2/7,3/14)=-h_1.
\]
The resulting effective covariance is
\begin{equation}\label{cps:log:eq:Ssigma}
 \Sigma_S=\begin{pmatrix}9/7&-22/21\\-22/21&9/7\end{pmatrix},
 \qquad\rho_S=-22/27.
\end{equation}
Its eigenvalues are $5/21$ and $7/3$. The physical covariance is
$4\Sigma_S$; the source's displayed matrix is $\gs$ times this
normalized stochastic matrix. Again the correlation is identical.

For an independent moment derivation, let $N$ be the number of faces
in one tilted aggregate:
$\PP(N=k)=(3/4)(1/4)^k$. Each face contributes $(\ell,r)$ with law
\[
 \PP(\ell,r)=\tfrac12\binom{\ell+r}{r}4^{-\ell-r}.
\]
Its marginal means are $1/2$, marginal variances $3/4$, and covariance
$1/4$. Let $L,R$ be the sums of these two coordinates over the $N$
faces. If $N=0$, then $d=1-c$; otherwise the terminal phase $d$ is a
uniform bit independent of $c,N,L,R$. Telescoping gives
\[
 \Delta=(c-N-L,\ N+R+d-1).
\]
With $X=N+L$, $Y=N+R$, $B=(3c+4d)/7$, the corrected increment is
$T=(-X+B,Y+B-1)$. Since $\EE N=1/3$ and $\Var N=4/9$,
\[
 \EE X=\EE Y=1/2,\quad \Var X=\Var Y=5/4,\quad\Cov(X,Y)=13/12.
\]
Also $\EE(B\mid N)=1/2$, $\Var B=1/28$ and
$\Cov(B,X)=\Cov(B,Y)=0$. These identities give~\eqref{cps:log:eq:Ssigma}.
The exact checker independently compares the result with the log
Perron-root Hessian and the Poisson-corrector calculation.

\begin{writenotelog}
$G$ is $K_S$ of \eqref{cps:eq:SM} in half-scale variables, and \eqref{cps:log:eq:Ssigma} is
\eqref{cps:eq:Sigma} divided by $4$ (Remark~\ref{cps:log:rem:scale}). The aggregate time, the
shifted quadrant and its equivalence with physical survival are Part~\ref{cps:part:exponents}'s
(Lemma~\ref{cps:lem:Sboundary}), proved here a second time.
\end{writenotelog}

\subsection{A moment bound that retains the whole aggregate}
Let $L_{\rm tot}$ be the sum of physical $\ell^1$ lengths of every
elementary step of one aggregate. At the critical tilt, the sum of
weighted face contributions in each row, with extra factor
$e^{t\,\text{length}}$, is
\[
 f(t)=\frac{e^{4t}}{8(1-e^{2t}/2)}.
\]
SE contributes $(3/4)e^{2t}$ after normalization. Summing the number
of faces gives the exact identity
\begin{equation}\label{cps:log:eq:Smoment}
 \EE_c e^{tL_{\rm tot}}=\frac{(3/4)e^{2t}}{1-f(t)},
 \qquad f(t)<1,\quad e^{2t}<2.
\end{equation}
For $t=\log(6/5)$, $e^{2t}=36/25$ and
$f(t)=162/175<1$; indeed the displayed moment equals $189/13$.
The aggregate displacement norm is bounded by $L_{\rm tot}$, and
quotient coordinates differ by division by two and a bounded phase
shift. Thus (H2) holds uniformly, including all within-aggregate
lengths rather than only a formal rational-kernel derivative.

\begin{writenotelog}
This explicit certificate is new to the report. Part~\ref{cps:part:exponents} obtains the
exponential moment of a return cycle from analyticity of \eqref{cps:eq:return} at $(1,1,1)$. The
arithmetic was rechecked at the write: $e^{2t}=36/25$, $f(t)=(1296/625)/(8\cdot7/25)=162/175$, and
$(3/4)(36/25)/(13/175)=189/13$.
\end{writenotelog}

\subsection{Aperiodicity and legal seeds}
Supported diagonal increments include phase-$0$ loops
$(-1,0),(-2,0)$ and a phase-$1$ loop $(0,1)$.
Equation~\eqref{cps:log:eq:phase-equality} first forces $t_1=0,\lambda=1$,
then $t_2=0$ modulo $2\pi$. Hence (H4) holds in exact aggregate time.

A forward seed from $(0,0,0)$ takes $0\to1:(0,0)$, then $2H$
phase-$1$ north loops, followed by $H$ repetitions of
$1\to0:(1,-1)$ and $0\to1:(0,0)$. It ends at $(H,H,1)$.
A dual seed from $(0,0,1)$ takes $1\to0:(0,0)$, then $2H$ phase-$0$
east loops, followed by $H$ repetitions of
$0\to1:(-1,1)$ and $1\to0:(0,0)$. It ends at $(H,H,0)$.
All edges have positive probability and all coordinates stay
nonnegative. The resulting fixed interior ray proves (H5).

\subsection{Endpoint factor and transfer to all labelings}
The truncated path has physical displacement $(1,-1)$, and its
telescoping tilt is $2^1(1/2)^{-1}=4$. For genuine sizes $n\ge1$,
\begin{equation}\label{cps:log:eq:Sendpoint}
 K_{n+1}((0,0,0),(0,0,1))=4\gs^{-n-1}s'_n.
\end{equation}
The positive fixed-shift relation from \cite[Proposition 19]{FNS} is
\begin{equation}\label{cps:log:eq:Sshift}
 s_n=s'_n+2s'_{n-1}+s'_{n-2}\qquad(n\ge4).
\end{equation}
Applying Theorem~\ref{cps:log:thm:cone} to~\eqref{cps:log:eq:Sendpoint} and then
using~\eqref{cps:log:eq:Sshift} proves~\eqref{cps:log:eq:mains}. Each positive term
has the same logarithmic power; no cancellation or random-time
replacement is involved.

\begin{writenotelog}
\eqref{cps:log:eq:Sendpoint} is \eqref{cps:eq:Scount}, and \eqref{cps:log:eq:Sshift} is
\eqref{cps:eq:Sconversion}; see Section~\ref{cps:log:sec:printing} for the two citations of
\cite{FNS}.
\end{writenotelog}

In finite enumeration a formal auxiliary walk value $s'_0=1$ is
convenient. It is an initialization outside the genuine size range
of the source bijection, not a count of a size-zero nonempty
dissection. The actual initial values are $s_0,s_1,s_2,s_3=0,0,3,2$;
\eqref{cps:log:eq:Sshift} is applied only from $n=4$.

\section{Non-D-finiteness from logarithmic powers}\label{cps:log:sec:nonD}
\subsection{Irrationality and the needed derivative order}
If $\theta/\pi$ were rational, $e^{i\theta}$ would be a root of
unity and $2\cos\theta$ would be an algebraic integer. In our two
cases the latter numbers are $9/8$ and $44/27$, rational nonintegers.
Thus both $\theta/\pi$, and hence both $\alpha=1+\pi/\theta$, are
irrational. Elementary angle comparisons give
\[
 4<\ap<5,\qquad 6<\as<7.
\]
For the second interval,
$\pi/6<\arccos(22/27)<\pi/5$: the nontrivial upper-angle inequality
uses $22/27>(1+\sqrt5)/4$, equivalent to
$61/27>\sqrt5$, and $3721>3645$ verifies its square. The first
interval follows from $\pi/4<\arccos(9/16)<\pi/3$.

\subsection{A positive derivative detects the real singular exponent}
Let $a_n$ denote either sequence and
$F(z)=\sum_{n\ge0}a_nz^n$. Put $R=\Gamma^{-1}$ and $H(t)=F(Rt)$.
Take $m=\lfloor\alpha\rfloor$, namely $m=4$ or $6$. The coefficients
of $H^{(m)}$ are positive for large indices and have size
$n^{m-\alpha+o(1)}$, with $m-\alpha\in(-1,0)$.
For any real $s>-1$,
\begin{equation}\label{cps:log:eq:powersum}
 \sum_{n\ge1}n^st^n\asymp(1-t)^{-s-1}\qquad(t\uparrow1).
\end{equation}
To see this, write $t=e^{-u}$. Restricting to
$n\in[1/u,2/u]$ gives the lower bound; partitioning into intervals
$[k/u,(k+1)/u]$ gives an upper bound from an exponentially convergent
sum and the integrability of $x^s$ near zero. Finally $u\sim1-t$.

For every sufficiently small $\epsilon>0$, positive termwise
comparison with exponents $m-\alpha\pm\epsilon$ yields
\[
 c_\epsilon(1-t)^{\alpha-m-1+\epsilon}
 \le H^{(m)}(t)
 \le C_\epsilon(1-t)^{\alpha-m-1-\epsilon}.
\]
Thus $H^{(m)}(t)\to+\infty$ and
\begin{equation}\label{cps:log:eq:irrational-limit}
 \lim_{z\uparrow R}\frac{\log F^{(m)}(z)}{\log(R-z)}
 =\alpha-m-1\notin\mathbb Q.
\end{equation}
The change from $H$ to $F$ only introduces a constant derivative
factor and a linear rescaling of distance to the singular point.

\subsection{Arithmetic local exponents give the contradiction}
Suppose $F$ were D-finite over $\mathbb Q(z)$. Its integer coefficients and
exponential upper bound make it a G-function: conjugates have the
same bound and common denominators are one. Its $m$th derivative
is also a G-function. The arithmetic local-structure theorem of
Andr\'e--Chudnovski--Katz, in the formulation
\cite[Theorem 6 and Corollary 1, pp.\ 328--329]{FR}, says that the
minimal operator is regular singular with rational exponents.
At the algebraic point $R$, on the branch approached along $z<R$,
its local expression is a finite sum of rational powers times
logarithmic powers and holomorphic factors. Combining equal
exponents and taking first nonzero Taylor coefficients gives
\[
 F^{(m)}(z)=c(R-z)^q(\log(R-z))^k(1+o(1)),\qquad
 q\in\mathbb Q,\quad k\ge0,\quad c\ne0.
\]
On the chosen real interval the leading sign is positive, because
$F^{(m)}$ is positive. Its logarithmic limit is therefore $q\in\mathbb Q$,
contradicting~\eqref{cps:log:eq:irrational-limit}. Possible operator
singularities inside the convergence disk do not matter: the original
power-series solution is analytic through them and determines this
branch at $R$.

The exclusion also holds over $\mathbb C(z)$. For a fixed order and degree,
the coefficients of a putative differential operator solve homogeneous
linear equations with rational entries, since $F$ has rational
coefficients. Their row space is finite dimensional and has a finite
rational basis. A nonzero complex null vector therefore entails a
nonzero rational null vector, reducing to the contradiction above.

This proof uses positivity to extract a real singular logarithmic
exponent. It does not invoke a full coefficient equivalent or assert
that arbitrary D-finite functions have only rational exponents.
The G-function arithmetic hypotheses are essential.

\begin{remark}[write, 2026-10-05: the criterion in general]\label{cps:log:rem:criterion}
\emph{Let $a_n\ge0$ be integers with $a_n=\Gamma^nn^{-\alpha+o(1)}$, $\Gamma>1$ and
$\alpha\in\RR\setminus\mathbb Q$. Then $\sum_na_nz^n$ is not D-finite over $\mathbb C(z)$.} This is what
the argument above proves; it is stated there for the two sequences, with $m=\lfloor\alpha\rfloor$.
In general take any integer $m\ge0$ with $m>\alpha-1$. The coefficient of $t^{n-m}$ in $H^{(m)}$ is
$n(n-1)\cdots(n-m+1)a_nR^{n}$, eventually positive and equal to $n^{m-\alpha+o(1)}$, and
$m-\alpha>-1$. For $0<\epsilon<m-\alpha+1$, the comparison \eqref{cps:log:eq:powersum} with
exponents $m-\alpha\pm\epsilon$ gives the two-sided bound above, so
$H^{(m)}(t)\to+\infty$ (the exponent $\alpha-m-1$ is negative) and the limit
\eqref{cps:log:eq:irrational-limit} equals $\alpha-m-1$, which is irrational. The $G$-function
step needs only integer coefficients and an exponential bound, both of which hold. By Pringsheim's
theorem $R=1/\Gamma$ is a singularity of $F$, hence a root of the leading coefficient of its
minimal operator, which lies in $\mathbb Q[z]$; so $R$ is algebraic and the local theorem applies
at $R$ (in particular $\Gamma$ is algebraic, as it must be) \textup{[}corrected at the independent check; see the end of
Section~\ref{cps:log:sec:research}\textup{]}.
The step then gives a rational limit, a contradiction; the reduction from $\mathbb C(z)$ to
$\mathbb Q(z)$ is the paragraph above. Lemma~\ref{cps:lem:nonD} is the special case $a_n=\Theta(\mu^nn^{-\alpha})$, and
Theorem~\ref{cps:thm:main} is not needed for the three sequences once
\eqref{cps:log:eq:mainp}--\eqref{cps:log:eq:mainb} are known; the irrationality of $\ap$ and $\as$ is
still used, and is proved at the start of this section by the argument of Lemma~\ref{cps:lem:irrational}
\textup{[}corrected at the independent check; see the end of
Section~\ref{cps:log:sec:research}\textup{]}. The bounds $4<\ap<5$ and
$6<\as<7$ above are rechecked: $\cos(\pi/3)=1/2<9/16<\cos(\pi/4)$, and
$\cos(\pi/5)=(1+\sqrt5)/4<22/27<\cos(\pi/6)$, the last because $(44/27)^2=1936/729<3$.
\end{remark}
\noindent(Dated note, 5~October 2026: Remark~\ref{cps:log:rem:criterion} first read ``The
coefficients of $H^{(m)}$ are $n(n-1)\cdots(n-m+1)a_nR^{n-m}$''. Since $H(t)=F(Rt)$, the
coefficient of $t^{n-m}$ in $H^{(m)}$ carries $R^{n}$; the factor $R^{n-m}$ belongs to the
coefficients of $F^{(m)}(Rt)$, which differs from $H^{(m)}(t)$ by the constant factor $R^{m}$.
The size $n^{m-\alpha+o(1)}$, the two-sided bound and the rest of the argument are unchanged. The
formula was corrected, and two sentences of the remark were clarified, at the independent check
recorded at the end of Section~\ref{cps:log:sec:research}.)

\section{First thresholds without monotonicity}\label{cps:log:sec:inverse}
Let $a_n=\Gamma^n n^{-\alpha+o(1)}$ with $\Gamma>1$ and put
$g=\log\Gamma$. The estimate gives eventual positivity and
unboundedness, so $N(y)=\min\{n:a_n\ge y\}$ is defined and tends
to infinity. First-threshold minimality itself implies
\[
 a_{N(y)-1}<y\le a_{N(y)}
\]
for large $y$. Since
\[
 \log a_n=gn-\alpha\log n+o(\log n),
\]
the two adjacent bounds first imply $N(y)\sim(\log y)/g$ and then
\[
 \log y=gN(y)-\alpha\log N(y)+o(\log N(y)).
\]
Substitution of the first-order estimate proves
\eqref{cps:log:eq:inverse-main}. In particular no estimate on successive
coefficient ratios, derivative of an error term, or eventual
monotonicity is assumed. The unknown $o(\log n)$ coefficient
remainder is exactly why the inverse remainder is only
$o(\log\log y)$, rather than $O(1)$.

\begin{writenotelog}
This is weaker than Lemma~\ref{cps:lem:inverse}, which gives the same two terms with error
$O(1)$ from the $\Theta$ bounds, and than Theorem~\ref{cps:amp:thm:inverse}. All three are
instances of \texttt{p0:thm:lambert-core} of the transseries volume named in the guide, here
at logarithmic precision; no novelty is claimed for the inversion mechanics.
\end{writenotelog}

\section{A six-face sleeve and the all-index A377921 result}\label{cps:log:sec:921}
Write $b_n=\mathrm{A377921}(n)$ for the subfamily in which all three
white outer vertices are nonisolated. Family inclusion gives $b_n\le s_n$.
The generating-function relation alone only controls a four-term
window. The following explicit injection supplies the missing
all-index lower bound.

\subsection{The source conventions}
In the source's definitions \cite[Section 2.1]{FNS}, a
$(6,4)$-dissection has a simple hexagonal outer boundary and all
inner faces quadrangular. Its root vertex is white. The outer labels
are $R,G,B,R,G,B$ counterclockwise from the root. Every inner edge
incident to an outer vertex has that vertex's prescribed label
(SL1). At each inner vertex, incident colors form three nonempty
clockwise blocks $R,G,B$ (SL2). The size is the number of inner faces.
Remark 3 gives the equivalent underlying condition of simplicity
and faciality of every four-cycle.

\begin{proposition}[Canonical sleeve injection]\label{cps:log:prop:sleeve}
For every genuine size there is an injection
\[
 \Phi:\mathcal S_n\hookrightarrow\widetilde{\mathcal S}_{n+6}
\]
on the rooted labeled map classes of the source. In particular,
\begin{equation}\label{cps:log:eq:sleeve-sandwich}
 s_{n-6}\le b_n\le s_n\qquad(n\ge8).
\end{equation}
\end{proposition}
\begin{proof}
Index the old outer vertices $v_0,\ldots,v_5$ counterclockwise, with
$v_0$ the white root. Indices below are modulo six; write
$c_i=R,G,B$ according to $i$ modulo three. Thus $v_i$ is white
exactly for even $i$ and its prescribed outer color is $c_i$.

Attach a topological annulus outside the old boundary. Its new
outer cycle is $w_0,\ldots,w_5$ in the same cyclic order, and its
six radial edges are $v_iw_i$. The six new quadrangular faces are
\[
 v_i v_{i+1}w_{i+1}w_i\qquad(0\le i<6).
\]
This construction is independent of the geometry of a particular
planar drawing. Give $w_i$ the bipartite color opposite to $v_i$,
keep all originally inner edge colors, and color the newly inner
edges by
\begin{equation}\label{cps:log:eq:sleevecolors}
 \operatorname{col}(v_iv_{i+1})=c_i,\qquad
 \operatorname{col}(v_iw_i)=c_{i+1}.
\end{equation}
Leave the new outer edges uncolored. Root the result at the unique
outer corner of $w_5$, which is white. In counterclockwise order
from this new root,
\[
 (w_5,w_0,w_1,w_2,w_3,w_4)\quad\hbox{have labels}\quad(R,G,B,R,G,B).
\]
Thus the label of $w_i$ is $c_{i+1}$. Its only inner edge is its
radial edge, of exactly that color. Each $w_i$ has degree three, so
all six new outer vertices are nonisolated and (SL1) holds.

Originally inner vertices have no change in their rotations or
colors. At a former boundary vertex $v_i$, the clockwise cyclic
order is the preceding boundary edge, all old inner edges, the
following boundary edge, and the radial edge. If $k_i\ge0$ old
inner edges meet $v_i$, their colors are all $c_i$ by old (SL1).
The new clockwise color list is
\[
 c_{i-1},\quad(c_i)^{k_i},\quad c_i,\quad c_{i+1}.
\]
These are exactly three nonempty blocks in cyclic $R,G,B$ order.
The following old boundary edge keeps the middle block nonempty
even if $k_i=0$. Thus (SL2) holds also at formerly isolated outer
vertices.

The new bipartition is consistent on boundary edges and radials,
the outer cycle is simple, and all old inner faces are unchanged.
There are no new multiple edges. One can also check irreducibility
directly. A four-cycle containing no new vertex is an old facial
cycle. A four-cycle using both old and new vertices crosses the
radial cut an even number of times. It cannot use four radials,
since each new vertex has only one radial edge, so it uses exactly
two. The two remaining edges must be one outer-ring edge and one
old edge. The radials consequently have consecutive indices, and
the old edge is the corresponding old boundary edge by simplicity.
The cycle is one of the six annular faces. A four-cycle entirely
among new vertices is impossible in their chordless hexagon.
Hence every four-cycle is facial.

Exactly six vertices, twelve edges and six inner faces were added.
The size increase is precisely six. There are three new vertices
of each bipartite color, so the refined size also changes by
$(a,b)\mapsto(a+3,b+3)$.

Finally the construction has an intrinsic rooted inverse on its
image. Each outer vertex of an output has a unique inner neighbor,
recovering the six old boundary vertices. The new boundary vertex
immediately counterclockwise after the new root is $w_0$; its
unique inner neighbor is the old root $v_0$. Delete the six new
outer vertices and all incident edges, expose the recovered old
hexagon, erase the colors on its boundary edges, and root at its
outer corner at $v_0$. Every old inner edge color is retained.
This recovers the input exactly. The procedure respects
orientation-preserving rooted isomorphisms, so symmetry creates
no multiplicity. This proves injectivity and the sandwich.
\end{proof}

\begin{corollary}[A377921 at every large index]\label{cps:log:cor:921}
The sequence $b_n=\mathrm{A377921}(n)$ satisfies
\begin{equation}\label{cps:log:eq:921alln}
 b_n=(16/3)^n n^{-\as+o(1)}.
\end{equation}
Its generating function is not D-finite over $\mathbb C(t)$, and
\[
 N_b(y)=\frac{\log y}{\log(16/3)}
       +\frac{\as}{\log(16/3)}\log\log y+o(\log\log y).
\]
\end{corollary}
\begin{proof}
The fixed six-unit shift obeys
\[
 \log s_{n-6}=n\log\gs-\as\log n+o(\log n).
\]
Taking logarithms in~\eqref{cps:log:eq:sleeve-sandwich} proves the all-index
estimate. The positive-series G-function and threshold arguments
in Sections~\ref{cps:log:sec:nonD}--\ref{cps:log:sec:inverse} apply directly.
\end{proof}

\begin{writenotelog}
Part~\ref{cps:part:amplitudes} proves the equivalent
$b_n=\widetilde s_n\sim(16/19)^3\kappa_S\mu_S^nn^{-\alpha_S}$, \eqref{cps:amp:eq:rigidmain}, by a
dominated signed convolution from the equivalent for $s_n$; the corollary is a second route on the
logarithmic scale that needs no amplitude. The two are consistent: at $n=30$,
$b_{30}/s_{30}=0.5429\ldots$ against the limit $(16/19)^3=0.5971\ldots$, and
$s_{24}\le b_{30}\le s_{30}$.
\end{writenotelog}

\begin{remark}[write, 2026-10-05: a $\Theta$ bound for A377921 from Part I]\label{cps:log:rem:theta}
\emph{$b_n=\Theta(\mu_S^nn^{-\alpha_S})$.} Indeed $b_n\le s_n$, and $b_n\ge s_{n-6}$ for $n\ge8$ by
Proposition~\ref{cps:log:prop:sleeve}. Theorem~\ref{cps:thm:main} gives positive $c,C$ with
$c\mu_S^nn^{-\alpha_S}\le s_n\le C\mu_S^nn^{-\alpha_S}$ for large $n$, hence
$b_n\ge c\mu_S^{-6}\mu_S^n(n-6)^{-\alpha_S}\ge c\mu_S^{-6}\mu_S^nn^{-\alpha_S}$. So the coefficientwise
$\Theta$ estimate that Part~\ref{cps:part:exponents} says does not follow from its argument
(after Corollary~\ref{cps:cor:nonD}) follows from it together with the sleeve, without
Part~\ref{cps:part:amplitudes}.
\end{remark}

\subsection{The independent positive-convolution consequences}
For comparison, set $A(t)=\sum_{n\ge2}s_nt^n$ and
$B(t)=\sum_{n\ge2}b_nt^n$. The source's exact relation
\cite[Remark 20]{FNS} is
\begin{equation}\label{cps:log:eq:921gf}
 A(t)=(1+t)^3(1+B(t))-1-3t.
\end{equation}
For $n\ge4$ this yields
\begin{equation}\label{cps:log:eq:921conv}
 s_n=b_n+3b_{n-1}+3b_{n-2}+b_{n-3}.
\end{equation}
Nonnegativity gives
\[
 s_n/8\le\max_{0\le j\le3}b_{n-j}
        \le\max_{0\le j\le3}s_{n-j}.
\]
Thus even without the sleeve one obtains the logarithmic estimate
for the four-term maximum, $\limsup b_n^{1/n}=\gs$, and
$N_s(y)\le N_b(y)\le N_s(8y)$ for large $y$. Also D-finiteness of
$B$ would force D-finiteness of $A$ by~\eqref{cps:log:eq:921gf}, giving a
short independent transfer of non-D-finiteness.

The distinction is important: a short positive convolution can
hide small coefficients in a parity class. Dividing by $(1+t)^3$
introduces alternating coefficients. Neither maneuver proves an
all-index lower bound; that bound comes from the explicitly rooted
sleeve injection. The injection does not supply a limiting ratio
$b_n/s_n$, a leading amplitude, or a full coefficient equivalent.

\begin{remark}[write, 2026-10-05: the same identity as Part I]\label{cps:log:rem:rigid}
\eqref{cps:log:eq:921gf} is \eqref{cps:eq:rigid}: with $A=S$, $B=T$ and $t=z$,
$(1+t)^3(1+B)-1-3t=(1+t)^3B+(1+3t+3t^2+t^3)-1-3t=(1+t)^3B+3t^2+t^3$. Report 112's caution is
borne out by Part~\ref{cps:part:amplitudes}: its proof of \eqref{cps:amp:eq:rigidmain} divides by
$(1+z)^3$, as here, and controls the alternating coefficients only by dominated convergence
against the full equivalent for $s_n$; a $\Theta$ or logarithmic bound would not suffice for that
step. The limiting ratio is $(16/19)^3$; Part~\ref{cps:part:amplitudes} gives no amplitude digits.
\end{remark}

\section{Why a full equivalent is still outside the result}\label{cps:log:sec:literature}
The logarithmic theorem loses arbitrarily small powers at its annular
crossings. Its proof consequently identifies the exponent but cannot
recover a stable amplitude. A sufficient stronger result would be a
fixed-endpoint killed local limit of the form
\[
 K_n((x,c),(y,d))\sim\kappa V_c(x)V_d^*(y)n^{-1-\pi/\theta},
\]
with positive forward and dual killed harmonic functions and the
correct lattice normalization. This is a different analytic task.

A scoped review of primary results found no theorem immediately
covering these exact two-dimensional lattice kernels with their
finite phases and unbounded edge-valued increments:
\begin{itemize}
\item Denisov--Wachtel's \emph{Random walks in cones revisited}
\cite{DWrevisited} assumes independent identically distributed increments.
Its cone harmonic and survival theory cannot simply be imported as
a finite-phase fixed-endpoint result.
\item Denisov--Zhang's \emph{Markov chains in the domain of attraction
of Brownian motion in cones} \cite{DZ} gives harmonic functions,
survival equivalents and a conditioned global endpoint limit under
conditional moment assumptions. It does not state the required
fixed-endpoint killed local limit, and its conditional covariance
assumptions fail for the exact kernels here.
\item Grama--Lauvergnat--Le Page \cite{GLLP} proves a conditioned
local limit for a scalar additive functional on a finite Markov
chain, killed at a half-line. Its Theorem 2.5 is not a theorem for a
two-dimensional wedge or these unbounded edge increments.
\item The March 2026 paper of Pham--Peign\'e--Son \cite{PPS} treats
nonlattice multidimensional random walks in cones; its setting is
not the present lattice internal-phase process.
\end{itemize}
This is a statement about the inspected hypotheses, not a claim that
no other applicable result exists.

The covariance obstruction in the second item is exact. After the
natural Poisson corrector, the phase-conditional covariance matrices
are
\begin{align*}
 V^P_0&=\begin{pmatrix}3/5&-9/20\\-9/20&1\end{pmatrix},&
 V^P_1&=\begin{pmatrix}1&-9/20\\-9/20&3/5\end{pmatrix},\\
 V^S_0&=\begin{pmatrix}3/2&-22/21\\-22/21&15/14\end{pmatrix},&
 V^S_1&=\begin{pmatrix}15/14&-22/21\\-22/21&3/2\end{pmatrix}.
\end{align*}
Their stationary averages are $\Sigma_P,\Sigma_S$, but the two
matrices in each pair are unequal. No common invertible whitening
can make both identity matrices. Their difference persists
arbitrarily deep in the cone because the transition law depends only
on phase. A martingale correction also shifts the boundary by a
bounded phase-dependent amount; it does not turn the exact killed
process into one on a single unchanged cone. The logarithmic proof
requires only the averaged FCLT covariance and handles the original
boundary throughout.

Regeneration is not an automatic remedy. A return block must retain
its duration and its within-block coordinate minima. Omitting the
duration changes the counting parameter; omitting the minima admits
paths that left the quadrant inside a block. A full-equivalent proof
would need to handle these marks or develop a phase-sensitive killed
harmonic and local-limit theory directly.

\begin{writenotelog}
For the two models of this report the full equivalent is no longer outside the report.
Part~\ref{cps:part:amplitudes} proves it by regeneration that keeps exactly the marks named
here: the cycle durations, by a two-sided random clock (Theorem~\ref{cps:amp:thm:exacttime}), and
the within-cycle minima, through the valid complete-cycle kernel $D$, squeezed as
$P_b\le D\le P_0$ between an iid cone kernel and its translate (\eqref{cps:amp:eq:kernelorder}),
which the one-unit buffer of Lemma~\ref{cps:lem:buffer} makes possible. It proves the killed
local limit from the free local theorem and an exit decomposition
(Lemma~\ref{cps:amp:lem:killed}), and the fixed-endpoint form asked for at the start of this
section is Theorem~\ref{cps:amp:thm:bridge}, $q_n(a,b)\sim\mathcal B_\Sigma A(a)A^*(b)n^{-p-1}$ with
positive harmonic $A$, $A^*$. Its proof uses Denisov--Wachtel \cite{DWrevisited} for the iid
cycle walk, not \cite{DZ}, \cite{GLLP} or \cite{PPS}. The review above remains an accurate account
of the papers it inspected; its sentence that no theorem covers these kernels was true of the
literature it checked, not of this report. In the generality of Theorem~\ref{cps:log:thm:cone}
the amplitude question is open (Question~\ref{cps:log:q:amplitude}). The phase-conditional
matrices displayed above are \eqref{cps:eq:PC} and \eqref{cps:eq:SC} divided by $4$, and
Part~\ref{cps:part:exponents} draws the same conclusion from them (after \eqref{cps:eq:Sigma}).
Report 112's bibliography cites \cite{PPS} from its abstract and paper record only.
\end{writenotelog}

\section{Further research and unresolved precision}\label{cps:log:sec:research}
The following are concrete next problems, not conclusions of the present
proof.
\begin{enumerate}
\item \textbf{Leading amplitudes and a killed local limit.}
Construct positive phase-indexed forward and dual harmonic functions for
the exact killed chains, and establish a fixed-endpoint equivalent with
its lattice normalization. This would strengthen the logarithmic theorem
rather than merely repeat the unrestricted local limit.

\begin{writenotelog}
Solved for the two models by Part~\ref{cps:part:amplitudes}: Theorem~\ref{cps:amp:thm:exacttime}
(positive harmonic functions $A$, $A^*$ at every relevant endpoint), Lemma~\ref{cps:amp:lem:killed}
and Theorem~\ref{cps:amp:thm:bridge} (the fixed-endpoint equivalent, with the lattice factor two).
Open in the generality of Theorem~\ref{cps:log:thm:cone}: Question~\ref{cps:log:q:amplitude}.
\end{writenotelog}
\item \textbf{A second phase correction.}
A possible route is to solve a finite-state Poisson equation for the
phase discrepancy $V(c)-\Sigma$ and couple the solution to the Hessian
of the Brownian wedge harmonic function. This may cancel the persistent
second-order phase oscillation. Boundary control, positivity, convergence
and the fixed-time local equivalent would still require separate proofs.

\begin{writenotelog}
Not pursued in Parts~\ref{cps:part:exponents}--\ref{cps:part:amplitudes}, which reach the
amplitudes through regeneration instead. The proposed route stays open as stated, with what is
missing listed by Report 112 itself (boundary control, positivity, convergence, the fixed-time
equivalent); it would matter for chains without a cycle buffer.
\end{writenotelog}
\item \textbf{Correction terms.}
Only after a leading equivalent is established should one determine
which further powers or logarithms occur, and whether the finite-phase
boundary contributes additional terms. No correction hierarchy or
transseries is inferred from $n^{o(1)}$ here.

\begin{writenotelog}
Still open. The leading equivalents now exist (Theorem~\ref{cps:amp:thm:main}); no rate and no
correction term is proved in any Part.
\end{writenotelog}
\item \textbf{The restricted-to-unrestricted ratio.}
The sleeve proves the exponent for $b_n$, but not a limit for $b_n/s_n$.
Determining that ratio, or its connection to possible leading amplitudes
through the exact generating-function identity, requires stronger
coefficient control that rules out alternating cancellation.

\begin{writenotelog}
Solved: the limit is $b_n/s_n\to(16/19)^3$, by \eqref{cps:amp:eq:Smain}--\eqref{cps:amp:eq:rigidmain}
and the dominated signed convolution of Part~\ref{cps:part:amplitudes}, the kind of coefficient
control asked for here. With the sleeve, Part~\ref{cps:part:exponents} alone already gives
$b_n=\Theta(\mu_S^nn^{-\alpha_S})$ (Remark~\ref{cps:log:rem:theta}).
\end{writenotelog}
\item \textbf{Sharper first thresholds.}
An $O(1)$ remainder in the inverse would require at least bounded
logarithmic coefficient error. Further constant-order inverse terms
would also need a precise amplitude and appropriate control of the
discrete first crossing. These data are currently absent.

\begin{writenotelog}
The $O(1)$ remainder is Lemma~\ref{cps:lem:inverse}, from the $\Theta$ bounds. The constant-order
term is Theorem~\ref{cps:amp:thm:inverse}: Lambert-$W_{-1}$ centering with the next constant
$-(\alpha\log\log\mu+\log\gamma)/\log\mu$ and a vanishing-width two-ceiling bracket. Still open: an
exact rounding rule $N=\lceil r\rceil$, which needs a quantitative remainder in the forward
asymptotic, and any numerical value of the amplitudes that enter it.
\end{writenotelog}

\subsection{Further questions recorded at the write}\label{cps:log:sec:questions}
Under Vladimir's standing rule of 4 October 2026, every claim of Report 112 that is not proved
in this report is recorded here or above with its source. Report 112 makes no claim that the
write found to be wrong.
\begin{enumerate}[label=\textup{Q\arabic*.},ref=Q\arabic*]
\item\label{cps:log:q:amplitude} \emph{An amplitude form of Theorem~\ref{cps:log:thm:cone}.} Under
 (H1)--(H5), possibly with a further hypothesis, is
 $K_n((0,a),(0,b))\sim\kappa\,V_a(0)V^*_b(0)\,n^{-1-\nu}$ with positive phase-indexed killed harmonic
 functions? Report 112, Sections~\ref{cps:log:sec:literature}--\ref{cps:log:sec:research}, poses
 it; Part~\ref{cps:part:amplitudes} answers it for the two models, using a fixed cycle buffer
 (Lemma~\ref{cps:lem:buffer}) that a general chain need not have. Missing: a harmonic comparison
 without a buffer, or a phase-sensitive killed local limit theory (Report 112's own
 suggestions); the crossing argument of Section~\ref{cps:log:sec:crossings} loses $q^{\pm\eta}$
 per annulus and cannot give a constant.
 \textup{[}Dated note, 5~October 2026: for one chain of this kind, the two-colour walk of A279571
 with unbounded steps, Part~I of \repo{a279571-inversion-cone-walk} (bundle Report~125) proves the
 amplitude form at every fixed pair of states,
 $K_n((s,c),(t,d))=(b_\theta\pi_d/\sqrt{\det\Sigma})\,V(s,c)V^*(t,d)\,n^{-1-\nu}+o(n^{-1-\nu})$
 (its display \texttt{icw:eq:killed-LLT}, where $\nu$ is written $p$), with no cycle buffer, by a
 phase-sensitive killed local limit theory. That walk is not an instance of this question ((H5)
 fails at the origin, its Remark~26.1), and the argument uses $4<\nu<5$ and $\nu>2$ (its
 Remark~26.4, \texttt{icw:log:rem:amplitude}); no general theorem is claimed there.\textup{]}
\item\label{cps:log:q:theta} \emph{A $\Theta$ form.} Do (H1)--(H5) give
 $K_n((0,a),(0,b))=\Theta(n^{-1-\nu})$? Part~\ref{cps:part:exponents} proves this for the two
 models by regeneration and the buffer. The proof of Theorem~\ref{cps:log:thm:cone} gives no rate
 for its $o(1)$ (Report 112 says so).
\item\label{cps:log:q:seeds} \emph{Endpoints without seeds.} Report 112 notes that a centred chain
 can have a boundary phase from which no legal interior excursion exists, and makes no claim
 there. What is the exponent, if any, at such endpoints?
 \textup{[}Dated note, 5~October 2026: A279571's walk is an example
 (\repo{a279571-inversion-cone-walk}, Remarks~26.1 and~26.4): at both phases of the origin the
 stationary dual has no seed, and the kernel from any state into the origin vanishes for every
 $n\ge2$ (from $o$, for every $n\ge1$), so there is no exponent there. The question remains open in
 general.\textup{]}
\item\label{cps:log:q:unchecked} \emph{Steps not independently rechecked at the write.} The uniform
 transfer from Brownian motion near the sides of the wedge (Section~\ref{cps:log:sec:crossings}),
 the bridge tightness in the proof of Lemma~\ref{cps:log:lem:bridge}, and the page locators of
 \cite[Theorem 2.1(ii)]{Whitt} and \cite[Theorem 6, Corollary 1]{FR} are given by Report 112 with
 the arguments printed above; the intake read them as sound but did not verify the cited pages.
 The locators of \cite{FR} were checked afterwards, against the journal version and the arXiv
 preprint; see the end of this section \textup{[}added at the independent check\textup{]}.
 The check that Report 111's chain satisfies (H1)--(H5)
 (Section~\ref{cps:log:sec:siblings}) is printed in that report, not here.
\item\label{cps:log:q:rates} \emph{Rates, corrections and rounding} (Report 112's items 3 and 5):
 open, as in the notes above.
\end{enumerate}
\end{enumerate}

\paragraph{Independent check of the write (dated note, 5~October 2026).}
An independent adversarial check of Part~\ref{cps:part:logcone}, made by the intake after the
write (\texttt{9b4001d29}), reread the four remarks of the write,
Remarks~\ref{cps:log:rem:scale}, \ref{cps:log:rem:criterion}, \ref{cps:log:rem:theta}
and~\ref{cps:log:rem:rigid}, and spot-checked three items of Report~112: the sleeve
Proposition~\ref{cps:log:prop:sleeve}, the moment certificate $189/13$ of
\eqref{cps:log:eq:Smoment} with the corner moment $524/405$ of
Section~\ref{cps:log:sec:repro}, and the bounds $4<\ap<5$, $6<\as<7$ of
Section~\ref{cps:log:sec:nonD}. It found all seven items valid, no counterexample, and no gap in
any proof chain. Four refinements were adopted where they stand. In
Table~\ref{cps:log:tab:notation} the row $m,h$ called the drift and corrector quarter-scale; they
live in half-scale coordinates, and only the covariances are a quarter of Parts~I--II's. The row
$G,C$, which called the Schnyder kernel quarter-scale, was corrected in the same way when the
check was applied. In Remark~\ref{cps:log:rem:criterion} the coefficient of $t^{n-m}$ in
$H^{(m)}$ now carries $R^{n}$, not $R^{n-m}$ (dated note after the remark; the two differ by the
constant factor $R^{m}$, so nothing else changes). The remark's sentence on what is not needed
now says that the irrationality of $\ap$ and $\as$ is still used, and where it is proved. Its
general statement now says why $R=1/\Gamma$ is algebraic, which the local step at $R$ uses:
by Pringsheim's theorem $R$ is a singularity of $F$, hence a root of the leading coefficient of
the minimal operator. No other statement or number changed.

The check also compared the locators of \cite{FR} with the paper. Part~\ref{cps:part:exponents}
cites Theorem~6 of Section~4.1, and Section~\ref{cps:log:sec:nonD} cites Theorem~6 and
Corollary~1 on pp.~328--329. They are right for the journal version that the bibliography names
(Theorem~6 on p.~328, Corollary~1 in Section~4.2 on p.~329, checked against the publisher's PDF
when the check was applied). In the preprint arXiv:1103.6022v2 the same results are Theorem~3
(Section~4.1) and Corollary~1 (Section~4.2). The bibliography entry now gives both numberings;
no citation in the text changed.

The check used neither the delivered programs nor the write's. In exact SymPy arithmetic it
recomputed both kernel sets. From Part~\ref{cps:part:logcone}'s kernels
\eqref{cps:log:eq:Pkernel} and \eqref{cps:log:eq:Gkernel}, solving $(I-P)h=m$, it recovered
$m_0$, $h_0$, $V^P_c$, $V^S_c$ and $\Sigma^{\mathrm{III}}_P$, $\Sigma^{\mathrm{III}}_S$ as
printed. From Part~\ref{cps:part:exponents}'s physical kernels \eqref{cps:eq:PK} and
\eqref{cps:eq:SM}, with Part~\ref{cps:part:exponents}'s own corrector equation, it recovered
$d_{\e}=(1/6,1/6)$ and $(-1/8,-1/8)$, $h^{\mathrm I}_{\e}=(1/10,1/10)$ and $(-1/14,-1/14)$, and
the matrices of \eqref{cps:eq:PC} and \eqref{cps:eq:SC}. The conversion formulas of
Remark~\ref{cps:log:rem:scale} reproduce $d$ and $h^{\mathrm I}$ in both phases exactly, and
$4V^P_c=C^P_c$, $4V^S_c=C^S_c$ for both phases, $\Sigma^{\mathrm{III}}_P=\Sigma_P/4$ and
$\Sigma^{\mathrm{III}}_S=\Sigma_S/4$. For Remark~\ref{cps:log:rem:criterion} it re-derived the
argument for every integer $m\ge0$ with $m>\alpha-1$: nonnegativity of all $a_n$ is used, for
the positivity of $F^{(m)}$, and the reduction from $\mathbb C(z)$ to $\mathbb Q(z)$ is correct.
It verified the bounds on $\ap$ and $\as$ exactly ($81/256<1/2$, $3645<3721$, $1936<2187$);
mpmath gives $\ap=4.2274760821\ldots$ and $\as=6.0803048461\ldots$.

For Remark~\ref{cps:log:rem:rigid} it read Remark~20 of \cite{FNS} in arXiv:2202.09172v3, which
is \eqref{cps:log:eq:921gf}, and confirmed symbolically that
$(1+t)^3(1+B)-1-3t=(1+t)^3B+3t^2+t^3$; both sums start at $n=2$. On the OEIS terms the identity
holds for all $31$ coefficients, and $s_n=b_n+3b_{n-1}+3b_{n-2}+b_{n-3}$ holds for
$n=4,\ldots,30$. For Proposition~\ref{cps:log:prop:sleeve} it compared the definitions used with
Section~2.1 of \cite{FNS} (white root; outer labels $R,G,B,R,G,B$ counterclockwise from the root;
the conditions on outer and inner vertices, the latter independent of the vertex colour;
``isolated'' meaning degree two), checked the rotation at $v_i$ on a convex model, the outer
labels with $w_5$ white, the four-cycle argument and the rooted inverse, and confirmed
$s_{n-6}\le b_n\le s_n$ on the OEIS terms for every $n=6,\ldots,30$ with no violation. The
smallest ratio $b_n/s_{n-6}$ is $14$, at $n=8$, and $b_{30}/s_{30}=0.542906$. The proposition is
stated for $n\ge8$; the inequality also holds at $n=6,7$. For Remark~\ref{cps:log:rem:theta} it
checked the direction of both inequalities and that, with Theorem~\ref{cps:thm:main} and
$\alpha>0$, they give a two-sided $\Theta$ for all large $n$; $b_n\mu_S^{-n}n^{\as}$ is $28.3$,
$60.6$ and $81.4$ at $n=10$, $20$ and $30$, bounded and slowly rising toward
$(16/19)^3\kappa_S$. For \eqref{cps:log:eq:Smoment} it substituted the critical tilt into the
face matrix, whose entries all become $x^4/(1-2x^2)$ with $x=e^t/2$; both row sums equal
$f(t)=162/175$ at $t=\log(6/5)$, SE contributes $\tfrac34e^{2t}=27/25$, the moment is $1$ at
$t=0$, and $(27/25)/(13/175)=189/13$ with $e^{2t}=36/25<2$. For $524/405$ it rebuilt the corner
quotient law (SE with $\ell^1$ length $1$ and probability $2/3$; a face with length $I+J+1$ and
probability $1/3$, with $I$, $J$ independent and geometric, $\PP(I=i)=\frac23 3^{-i}$), checked
that it reproduces the kernel weight $\frac2{27}3^{-i-j}$, and summed
$\frac23R+\frac4{27}R/(1-R/3)^2=\frac45+\frac{40}{81}=\frac{524}{405}$ at $R=6/5$. This was a
careful reading with exact computations and data checks, not a formal verification.

\section{Reproducibility and interpretation of the checks}\label{cps:log:sec:repro}
The accompanying reader package contains this editable \LaTeX\ source,
the PDF, exact symbolic and finite-enumeration checks, dependency
instructions, a source ledger and a strict checksum manifest. It
contains no third-party source PDFs. The package README gives the
build and replay commands.

The computational checks reconstruct both rational kernels and their
normalizations; verify stationary centering, correctors, effective and
phase-conditional covariances; check exponential-moment witnesses,
support edges, endpoint factors and shifts; verify finite sleeve topology and rooted recovery cases; and reproduce the
published finite prefixes from the original walk descriptions.
They run with ordinary Python and with optimization enabled, without
relying on removable \texttt{assert} statements. Deliberately damaged
fixtures are required to fail, and a fresh extracted archive is
replayed. The PDF build fixes its date and metadata inputs.

For corner paths, finite enumeration prunes a state at time $k$ when
$y>N-k$: only SE lowers $y$, by one, so the state cannot contribute
to any horizontal-axis endpoint by time $N$. This proves the requested
$p_n$ prefix through $N$, but it does \emph{not} certify an auxiliary
$(1,1)$ endpoint count at time $N$ with that same pruning. The checker
must either enlarge the height budget for such auxiliary endpoints
or omit the truncated boundary value.

For Schnyder paths, each SE-count layer is closed under face steps.
Every face strictly decreases $x$, so processing $x$ in descending
order is acyclic and exact. Binomial multiplicities and the SE time
parameter are retained. The four exceptional initial $s_n$ values
are entered explicitly; the positive relation is used only from
$n=4$. The formal auxiliary initialization $s'_0=1$ is not promoted
to a combinatorial size-zero count.

Finite checks certify the listed algebraic identities and enumerated
ranges, not the asymptotic theorem. The all-index conclusions depend
on the proofs in this report and the cited source bijections. Neither
matching a long prefix nor passing a mutation suite proves a cone
asymptotic. Conversely the proof's analytic steps do not depend on
numerical fitting or extrapolation from those prefixes.

\begin{writenotelog}
In this repository the package's PDF, README, source and checksum inventory
(\texttt{manifest.json}) are not shipped; its programs are shipped as
\texttt{code/47-logcone-*.py}, its fixtures and recorded outputs as \texttt{data/47-logcone-*}, and
its two check notes at the report root. The checker imports \texttt{verify\_sleeve} and reads
\texttt{fixtures.json} beside itself, and the integrity, seal, build and replay programs need the
delivered layout, the manifest or the original archive, so every rerun is made on a copy; the
report README gives the routes. Rerun on 5 October 2026 on such a copy, \texttt{checks/validate.py}
reproduced \texttt{data/47-logcone-checks-results-ordinary.log} and \texttt{.json} modulo line
endings. Besides the moment $189/13$ of \eqref{cps:log:eq:Smoment}, the checks certify the corner
quotient-displacement moment $(2/3)R+(4/27)R/(1-R/3)^2=524/405$ at $R=e^t=6/5$, not stated in the
text. The fixed-endpoint counts $e_L$ are enumerated with the enlarged height budget described
in the paragraph on corner paths above.
\end{writenotelog}

\appendix

\section{Exact computational checks}\label{cps:app:checks}
The scripts require Python 3 and SymPy. All paths in the reproducibility package are relative to the package root. Running \texttt{python3 scripts/run\_all.py} regenerates the verification outputs, and running \texttt{bash build.sh} recompiles the article.

\begin{writenote}
These are delivery
paths. In this repository the scripts are shipped as \texttt{code/45-cone-*.py}, their outputs
and the OEIS excerpts as \texttt{data/45-cone-*}, and \texttt{build.sh} as
\texttt{code/45-cone-build.sh}, which builds a delivery-layout \texttt{article.tex} (source 45
alone). The scripts resolve \texttt{scripts/}, \texttt{results/} and \texttt{sources/} relative to
their parent directory, so they must be run on a scratch copy in the delivered layout; the
report README gives the recipe.
\end{writenote}

\subsection{Symbolic identities}
The exact-arithmetic checks differentiate the two rational kernel matrices and their return generating functions. They verify stochasticity, the conditional drifts, both forward and reverse correctors, each conditional covariance matrix, the zero cycle means, both parity cycle covariances, the mean and variance of the cycle duration, the covariance angles, and the $S$ aggregate resolvent identity. No floating-point approximation is used for these assertions. Decimal exponents in the results table are evaluated only after the exact formulas have been established.

\subsection{Finite cycle geometry}
The supplied independent $P$ audit exhaustively tests the one-unit minimum bound and its reversed form for 25,900 nontrivial cycles with bounded face-step parameters and at most three middle parity-preserving steps. This finite test is useful for detecting sign errors. The proof for unbounded cycles is Lemma~\ref{cps:lem:buffer}, not an extrapolation from the test.

\subsection{Integer enumeration}
For $P$, an elementary dynamic program counts admissible physical walks. When counting walks with total length $N$, a state at time $t$ with ordinate greater than $N-t$ can be discarded: only a southeast step can decrease the ordinate, and it does so by one. Nonnegative horizontal coordinates bound possible face-step horizontal decrements. These restrictions make every counting sum finite without approximation.

For $S$, the package includes the weighted southeast/face recurrence and a separately implemented direct enumeration check. The weights $\binom{\ell+r}{r}$ are exact integers. The relation \eqref{cps:eq:Sconversion} converts the restricted counts $s'_n$ to $s_n$. Initial coefficients provide explicit indexing checks:
\begin{center}
\begin{tabular}{@{}r*{2}{r}@{}}
\toprule
$n$&$p_n$&$s_n$\\
\midrule
0&0&0\\
1&0&0\\
2&0&3\\
3&1&2\\
4&0&3\\
5&3&6\\
6&4&14\\
7&15&36\\
8&39&102\\
9&122&306\\
10&375&972\\
11&1212&3216\\
12&3980&11050\\
13&13413&39138\\
14&45966&142446\\
15&160295&530656\\
\bottomrule
\end{tabular}
\end{center}
The longer recorded Schnyder recurrence output reaches $n=30$ and agrees with the supplied OEIS export. Comparisons are exact coefficient comparisons; no numerical fit is used to infer either exponent.


\clearpage
\begin{thebibliography}{99}
\bibitem{FNS}
\'E. Fusy, E. Narmanli, and G. Schaeffer,
\emph{Enumeration of corner polyhedra and 3-connected Schnyder labelings},
Electronic Journal of Combinatorics \textbf{30}(2) (2023), P2.17,
\href{https://doi.org/10.37236/11174}{doi:10.37236/11174}; arXiv:2202.09172v3.
\href{https://arxiv.org/html/2202.09172v3}{arxiv.org/html/2202.09172v3}.
In particular, Propositions 13--14, Remarks 9 and 20, Section 4.2, Theorem 24, and Conjecture 25.

\bibitem{DWrevisited}
D. Denisov and V. Wachtel,
\emph{Random walks in cones revisited},
Annales de l'Institut Henri Poincar\'e, Probabilit\'es et Statistiques
\textbf{60}(1) (2024), 126--166.
\href{https://doi.org/10.1214/22-AIHP1331}{doi:10.1214/22-AIHP1331}; arXiv:2112.10244.
\href{https://arxiv.org/html/2112.10244}{arxiv.org/html/2112.10244}.
Theorems 2 and 3.

\bibitem{Whitt}
W. Whitt, \emph{Proofs of the martingale FCLT}, Probability Surveys
\textbf{4} (2007), 268--302, Theorem 2.1(ii), pp.\ 270--271.
\href{https://doi.org/10.1214/07-PS122}{doi:10.1214/07-PS122}.
\href{https://www.columbia.edu/~ww2040/WWproofs2007pub.pdf}{Primary text}.

\bibitem{DZ}
D. Denisov and K. Zhang,
\emph{Markov chains in the domain of attraction of Brownian motion in cones},
Journal of Theoretical Probability \textbf{38} (2025), article 14.
\href{https://doi.org/10.1007/s10959-024-01369-7}{doi:10.1007/s10959-024-01369-7}.
\href{https://arxiv.org/html/2309.16311}{Primary text}.

\bibitem{GLLP}
I. Grama, R. Lauvergnat and \'E. Le Page,
\emph{Conditioned local limit theorems for random walks defined on finite Markov chains},
Probability Theory and Related Fields \textbf{176} (2020).
\href{https://doi.org/10.1007/s00440-019-00948-8}{doi:10.1007/s00440-019-00948-8}.
\href{https://arxiv.org/html/1707.06129}{Primary text}.

\bibitem{PPS}
Thi da Cam Pham, Marc Peign\'e and Doan Thai Son,
\emph{A local limit theorem for nonlattice multidimensional random walks in cones},
arXiv:2603.26228 (27 March 2026).
\href{https://arxiv.org/abs/2603.26228}{Primary abstract and paper record}.

\bibitem{DWoriginal}
D. Denisov and V. Wachtel,
\emph{Random walks in cones},
Annals of Probability \textbf{43} (2015), 992--1044.
\href{https://doi.org/10.1214/13-AOP867}{doi:10.1214/13-AOP867}.

\bibitem{FR}
S. Fischler and T. Rivoal,
\emph{On the values of $G$-functions},
Commentarii Mathematici Helvetici \textbf{89} (2014), 313--341.
\href{https://doi.org/10.4171/CMH/321}{doi:10.4171/CMH/321}.
Theorem 6, Section 4.1, p. 328; Corollary 1, Section 4.2, p. 329.
In the preprint arXiv:1103.6022v2 the same results are Theorem 3 (Section 4.1) and Corollary 1
(Section 4.2) (both versions checked 5 October 2026).

\bibitem{BRS}
A. Bostan, K. Raschel, and B. Salvy,
\emph{Non-D-finite excursions in the quarter plane},
Journal of Combinatorial Theory, Series A \textbf{121} (2014), 45--63.
\href{https://specfun.inria.fr/bostan/publications/BoRaSa12.pdf}{Author-hosted manuscript}.
Section 2.2 and the proof of Theorem 2.

\bibitem{BS}
R. B\~anuelos and R. G. Smits,
\emph{Brownian motion in cones},
Probability Theory and Related Fields \textbf{108}(3) (1997), 299--319.
\href{https://doi.org/10.1007/s004400050111}{doi:10.1007/s004400050111}.
Lemma 1, equation (2.2); \href{https://www.math.purdue.edu/~banuelos/Papers/cones.pdf}{author-hosted PDF}.

\bibitem{OEISP}
The OEIS Foundation Inc., \emph{The On-Line Encyclopedia of Integer Sequences},
\href{https://oeis.org/A377922}{A377922}.
Official data export, \href{https://github.com/oeis/oeisdata/blob/main/seq/A377/A377922.seq}{oeis/oeisdata}.

\bibitem{OEISS}
The OEIS Foundation Inc., \emph{The On-Line Encyclopedia of Integer Sequences},
\href{https://oeis.org/A377920}{A377920} and
\href{https://oeis.org/A377921}{A377921}.
Official data exports, \href{https://github.com/oeis/oeisdata/tree/main/seq/A377}{oeis/oeisdata}.
\end{thebibliography}
\end{document}
